\documentclass[11pt]{article}
\usepackage[a4paper,margin=1in]{geometry}
\usepackage{amsmath,amssymb,mathtools,amscd}
\usepackage{mathrsfs}
\usepackage{microtype}
\usepackage[colorlinks=true,linkcolor=blue,citecolor=blue]{hyperref}

\newcommand{\R}{\mathbf R}
\newcommand{\FullMot}{\mathsf{Mot}_{\mathrm{full}}}
\newcommand{\RecOneMot}{\mathsf{Mot}_{\mathrm{rec},\le1}}
\newcommand{\DepthFib}{\mathsf{DepthFib}}
\newcommand{\SCart}{\mathsf{SCart}}
\newcommand{\Kak}{\mathsf{KakCart}}
\newcommand{\Mot}{\mathrm{Mot}}
\newcommand{\supp}{\mathrm{supp}}
\newcommand{\hofib}{\operatorname{hofib}}
\newcommand{\hocolim}{\operatorname*{hocolim}}
\newcommand{\holim}{\operatorname*{holim}}
\newcommand{\cofib}{\operatorname{cofib}}
\newcommand{\fib}{\operatorname{fib}}

\title{\textbf{Filtered Descent for the Physical Kakeya Incidence Source}}
\author{Chenxi Cai\\\texttt{Chenxi\_Cai@live.com}}
\date{11 September 2026}

\begin{document}
\setlength{\emergencystretch}{3em}
\allowdisplaybreaks
\maketitle

\begin{abstract}
We construct a finite, source-hereditary descent for shaded Kakeya incidence
measures.  A symmetric joint packet law keeps the forgotten translation fibre,
all support marginals, and every later restriction in one probability space.
Exact-support determinants, literal face deletion, root exchange, conditional
insertion kernels, and a reduced history tree then reduce the global filtered
estimate to one root-cross physical energy.

The analytic closure is geometric.  The planar C\'ordoba estimate
\cite{Cordoba77} gives the
two-dimensional base.  In dimension three the sticky Kakeya theorem and its
sticky/non-sticky reduction close the root gate.  At the next handoff,
directionwise sampling and an unbiased fibre lift reduce the remaining
power-small source to the marked four-dimensional sticky estimate of
\cite{StickyContact4D}.  Hence the scalar descent closes through dimension
four; in dimensions \(d\ge5\) the analogous root-cross gate remains an input.

The Selmer--Cartan tower \cite{UnifiedTower} is used only as structural
provenance for this concrete filtration: its all-support comparison identifies
the determinant successor and records coefficient, confluence, and
root--Moore labels without turning them into physical multiplicity.  It
supplies no Kakeya estimate.  The hereditary typed statement has the further,
independent requirement of a positive Hall/Radon--Nikodym capacity on
correlated pair sources; this is not implied by the tower or by the scalar
sticky input.
\end{abstract}

\paragraph{Standing inputs and scope.}
The physical source is a finite direction-separated family of shaded
\(\delta\)-tubes.  All branches and deletions are measured in the single
incidence law \eqref{eq:ordinary-incidence-atom-law} and its symmetric packet
extension \eqref{eq:symmetric-root-tuple-law}.  The initial hereditary shading
density obeys \(\lambda_{\rm in}^{-1}=\delta^{-o(1)}\) for the subpower
formulation; otherwise its reciprocal is retained explicitly.  The analytic
inputs are the hereditary planar C\'ordoba estimate \cite{Cordoba77}, the three-dimensional
sticky and sticky/non-sticky theorems
\cite{StickyKakeya,GuthWangZahl3D}, and the uniform marked
four-dimensional estimate \cite{StickyContact4D}.  These close
\eqref{eq:root-cross-energy-gate} in dimensions two, three, and four,
respectively.

For compatibility with \cite{UnifiedTower}, fix finite support, coefficient,
confluence, and Cartan ceilings and a finite set of odd coefficient primes.
The tower supplies the normalized all-support controller and its
Moore--Reedy reconstruction.  In particular, a \(q^2\) line records supplied
confluence provenance; it is not the physical pair multiplicity and is never
used as an analytic estimate.  The scalar proof below needs only the concrete
support filtration, the exact-support/face decomposition, the hereditary
source ledger, and proper-history recombination.  The additional motivic
labels are carried passively and then removed by the normalized physical
augmentation.

\section{The finite physical descent ledger}

The descent is carried by a concrete finite weighted incidence system.  A state
records a direction support \(I\), a retained insertion history \(\chi\), and
the ordinary incidence atom.  Coefficient, valuation, and Cartan labels may be
attached when comparison with \cite{UnifiedTower} is desired, but they do not
alter the underlying physical measure.  Support deletion forgets slots,
history reduction passes to a proper predecessor, and root projection retains
the forgotten-direction fibre density.

Write \(\mathsf{Inc}^{\rm FSC,rec}_{Q,T_0}\) for this finite support/history
system at a rooted cell.  The compatible full-motive decoration is only the
base change
\begin{equation}\label{eq:physical-fsc-full-base-change}
 \boxed{
 \mathsf{Inc}^{\rm FSC,full}_{Q,T_0}
 :=
 \mathsf{Inc}^{\rm FSC,rec}_{Q,T_0}
 \times_{\mathsf{SCart}^{\rm rec}}
 \mathsf{SCart}^{\rm full}
 \simeq
 \mathsf{Inc}^{\rm FSC,rec}_{Q,T_0}
 \times_{\RecOneMot}\FullMot.}
\end{equation}
Thus valuation depth belongs to the passive \(\FullMot\)-factor, whereas the
physical lift count \(\mathfrak q\) remains in the weighted middle fibre.  No
summation over a decoration fibre is allowed to create tube multiplicity.

Only four concrete properties of this ledger enter the scalar proof.

\begin{enumerate}
\item The exact-support term is the rooted transverse determinant, and deleting
a support slot gives its literal proper face.
\item Insertion kernels compose by conditional probability; hence insertion
order, confluence, restriction, and rerooting are compatible on one joint
source.
\item Every retained restriction of packet mass \(\alpha\) has one-point
marginals bounded by \(\alpha^{-1}\mu_{\rm inc}\).  We refer to this marginal
statement below as R5.
\item All histories in one load layer are summed before taking the physical
\(L^2/L^1\) estimate.  We refer to this global confluence estimate as R6.
\end{enumerate}

The first three properties are proved directly in the next four sections.
For the fourth, if \(\gamma\) ranges over retained history blocks and
\[
 \rho_\gamma
 :=
 \frac{d(\pi_\gamma)_*\mu_\gamma}{d\mu_{\rm inc}},
\]
then the unmarked observable is
\begin{equation}\label{eq:rn-corrected-evaluation}
 \mathcal E_{\rm inc}(f)
 :=
 \sum_\gamma\rho_\gamma\,\mathcal E_\gamma(f_\gamma).
\end{equation}
The sum in \eqref{eq:rn-corrected-evaluation} is not replaceable by one
normalized chart.  The analytic work below recombines all strict predecessor
histories and reduces the remaining part of R6 to
\eqref{eq:root-cross-energy-gate}.

The Selmer--Cartan input is limited to structural provenance.  Its normalized
controller has one universal exact-support generator \(b_I^{\rm univ}\) and
one latching generator \(\beta_I^{\rm univ}=2b_I^{\rm univ}\).  Their motivic
images satisfy
\begin{equation}\label{eq:latching-root-comparison}
 \alpha_I^{\rm lat,\Mot}=2\alpha_I^{\rm root,\Mot}.
\end{equation}
The physical images are respectively
\([\mathfrak Z^{(d)}_{Q,T_0}]\) and
\(2[\mathfrak Z^{(d)}_{Q,T_0}]\), as verified below.  This common-source
comparison records normalization and successor provenance; it neither
identifies a torsion line with a real determinant line nor infers physical
nonvanishing from a motivic class.

\section{The fixed-root exact-support bridge}

We now construct the first, source-compatible part of that realization.  Fix a
cube \(Q\), a root tube \(T_0\) meeting \(Q\), and the ordinary conditional law
\[
 \mu_{Q,T_0}(U)
 =\Pr(U\mid U\text{ is an admissible side tube at }(Q,T_0)).
\]
For the root-symmetric FSC source below, admissible means the full
\(Q\)-star; the atom \(U=T_0\) and all repeated side tuples have zero
determinant and therefore do not alter the exact-support detector.
Put \(W=e_{T_0}^{\perp}\), \(r=d-1\),
\[
 I=\{0,1,\ldots,r\},\qquad I^\circ=I\setminus\{0\},
\]
and \(s(U)=P_W e_U\).  Thus \(I\) is the unified tower's full direction support,
including the root, whereas \(I^\circ\) indexes the transverse slots.  An
ordered side tuple
\[
 \mathbf U=(U_1,\ldots,U_r)
 \quad\text{has law}\quad
 \mu_{Q,T_0}^{\otimes r}(\mathbf U),
\]
and its exact-support detector is
\begin{equation}\label{eq:physical-exact-support-detector}
 D_I(\mathbf U)
 =\det_W\bigl(s(U_1),\ldots,s(U_r)\bigr),
 \qquad |I|=d,\quad |I^\circ|=r.
\end{equation}

The physical packet cochain is not a chosen graph or a selected skeleton.  It
is the orthogonal direct sum over all ordered tuples,
\begin{equation}\label{eq:rooted-packet-cochain}
 \mathfrak Z^{(d)}_{Q,T_0}
 =\bigoplus_{\mathbf U}
 \mu_{Q,T_0}^{\otimes r}(\mathbf U)^{1/2}\,
 D_I(\mathbf U)\,z_{\mathbf U},
\end{equation}
where the \(z_{\mathbf U}\) are unit generators of the tuple blocks.  If
\[
 M_{Q,T_0}=\mathbb E_{U\sim\mu_{Q,T_0}}
       [s(U)s(U)^{\mathsf T}],
\]
then the ordered Cauchy--Binet identity gives the exact norm formula
\begin{equation}\label{eq:packet-cauchy-binet}
 \boxed{
 \|\mathfrak Z^{(d)}_{Q,T_0}\|^2
 =\sum_{\mathbf U}\mu_{Q,T_0}^{\otimes r}(\mathbf U)
       |D_I(\mathbf U)|^2
 =r!\det M_{Q,T_0}.}
\end{equation}
This is an identity on the ordinary incidence source.  No exceptional tuple
has been selected and no change of marginal has occurred.

\paragraph{Common-controller comparison of the exact-support lines.}
For the \(i\)-th side entry let \(L_i\) denote its transverse rank-one slot.
Root conditioning identifies the physical determinant line with the rooted
physical coefficient line
\begin{equation}\label{eq:determinant-motivic-line}
 \Lambda^r W
 \cong L_1\otimes\cdots\otimes L_r
 =:\mathcal J_{I^\circ}^{\rm phys}.
\end{equation}
It is not identified with the unified tower's torsion detector
\(\mathcal D_I^{\Mot}=\mathcal M_N[I]\otimes_{\mathbf Z}\mathbf Z[I]\).
Instead both receive maps from the integral controller line
\(\mathcal C_I^{\rm univ}=\mathbf Z b_I^{\rm univ}\).  For
\(|I|=d\ge3\), in the standing prime-power all-support range of the complete supportwise
Moore--Reedy realization, the comparison is
\begin{equation}\label{eq:normalized-controller-comparison}
 \begin{aligned}
 \Phi_{\rm phys}(b_I^{\rm univ})
   &=[\mathfrak Z^{(d)}_{Q,T_0}],
 &\Phi_{\Mot}(b_I^{\rm univ})
   &=\alpha_I^{\rm root,\Mot},\\
 \Phi_{\rm phys}(\beta_I^{\rm univ})
   &=2[\mathfrak Z^{(d)}_{Q,T_0}],
 &\Phi_{\Mot}(\beta_I^{\rm univ})
   &=\alpha_I^{\rm lat,\Mot}
     =2\alpha_I^{\rm root,\Mot},
 \end{aligned}
 \qquad
 \beta_I^{\rm univ}=2b_I^{\rm univ}.
\end{equation}
This records the normalization without requiring a nonexistent additive map
from \(\mathbf Z/N\) to \(\mathbf R\).

The two first nontrivial ranks display the structural comparison concretely.
In dimension \(3\), the top coefficient is
\[
 \alpha_1\beta_2-\beta_1\alpha_2
 =\det(s(U_1),s(U_2)).
\]
In dimension \(4\), the top \(UT_5\) latching coefficient is
\[
 \mu_a\kappa_{bc}-\mu_b\kappa_{ac}+\mu_c\kappa_{ab}
 =\det(s(U_1),s(U_2),s(U_3)).
\]
The degree-two boundary has zero projection to the corresponding central top
line, so these coefficients define the required closed exact-support classes.
These are the first two instances of the all-support bar--cobar comparison
constructed below.  Formula \eqref{eq:packet-cauchy-binet} is the corresponding
norm identity; it is not itself used to infer the higher-rank comparison.

\paragraph{Support deletion is literal.}
The tuple blocks in \eqref{eq:rooted-packet-cochain} are direct summands.
Deleting a side tube \(U\) deletes exactly the blocks whose tuple contains
\(U\).  The differential, every proper Hall-face map, and the detector all
preserve the remaining direct sum.  Consequently the unnormalized detector
energy can only decrease, and after deletion it is again exactly
\(r!\det M'_{Q,T_0}\) for the restricted raw measure.  At the coordinate level,
deleting \(i\in I^\circ\) replaces the top wedge by its
\(I^\circ\setminus\{i\}\) face; in the exact-support quotient every such proper
face is killed.  This proves all side-face compatibilities at a fixed root
without a geometric covering argument.  The remaining face
\(I\setminus\{0\}\) changes the chosen root and is not supplied by this direct
sum argument.

\paragraph{The source ledger.}
Each one-point marginal of
\(\mu_{Q,T_0}^{\otimes r}\) is exactly \(\mu_{Q,T_0}\).  If a later descent step
conditions the tuple law on an event \(A\), then for every side-tube set \(B\)
\begin{equation}\label{eq:conditional-marginal-ledger}
 \Pr(U_i\in B\mid A)
 \le \frac{\mu_{Q,T_0}(B)}{\Pr(A)}.
\end{equation}
Thus the R5 loss is at most \(\delta^{-o(1)}\) precisely when the retained event
has ordinary tuple mass at least \(\delta^{o(1)}\).  Before such a conditioning,
R5 holds with constant \(1\).  Equation
\eqref{eq:conditional-marginal-ledger}, rather than cohomological nonvanishing,
is the complete fixed-root source ledger.

\paragraph{Root projection and the antecedent.}
Let \(\pi_{T_0}=P_W\).  For a shaded side tube \(U\), do not immediately replace
its projection by a binary planar shading.  Retain the fibre density
\begin{equation}\label{eq:fibre-density}
 g_U(z)=\int_{\pi_{T_0}^{-1}(z)}\mathbf 1_{Y(U)}(x)\,
           d\mathcal H^1(x),
 \qquad
 \int_W g_U(z)\,dz=|Y(U)|.
\end{equation}
The second identity is Fubini, so projected mass is exact.  Moreover, in an
oriented orthonormal frame beginning with \(e_{T_0}\),
\begin{equation}\label{eq:root-wedge-projection}
 \det(e_{T_0},e_{U_1},\ldots,e_{U_r})
 =\det_W(P_W e_{U_1},\ldots,P_W e_{U_r}).
\end{equation}
Hence the physical antecedent of the root class is the measure-valued packet
\[
 \bigl(\{(\pi_{T_0}U,g_U)\}_U,\mu_{Q,T_0}\bigr),
\]
not merely the set of projected tube axes.  Equations
\eqref{eq:fibre-density} and \eqref{eq:root-wedge-projection} prove compatibility
with root projection while retaining both the forgotten-direction fibre and
the ordinary one-point law.  Converting \(g_U\) to a strong binary shading is a
later physical estimate and must record the discarded mass.

The construction in this section proves the root-conditioned determinant
identity, literal side-face compatibility, and the fixed-root marginal bound.
The following sections extend the weighted source realization to insertion,
confluence, the root face, rerooting, and global shaded restrictions.  The
extension to the all-support structural controller is given by the normalized
bar--cobar construction below.

\section{Insertion spans and confluence at one root}

Fix \((Q,T_0)\), the full support \(I\), and one joint ordinary law
\(\mathbf P_I\) on fully marked ordered side tuples.  The product law used in
\eqref{eq:rooted-packet-cochain} is one example, but independence is not needed
here.  For \(A\subseteq I\) containing the root slot, let \(X_A\) be the finite
set of \(A\)-marked physical packets and let \(\mathbf P_A\) be the marginal of
\(\mathbf P_I\).

For \(A\subseteq B\), define the insertion-history middle set
\begin{equation}\label{eq:physical-insertion-middle}
 H_{A,B}
 :=\{(x_A,x_B)\in X_A\times X_B:
          x_B|_A=x_A\},
\end{equation}
with its two projection legs.  On the support of \(\mathbf P_A\), give this
span the conditional kernel
\begin{equation}\label{eq:physical-insertion-kernel}
 K_{A,B}(x_B\mid x_A)
 :=\frac{\mathbf P_B(x_B)}{\mathbf P_A(x_A)}
    \mathbf 1_{\{x_B|_A=x_A\}}.
\end{equation}
Zero-mass states are discarded from the ordinary source.  Then
\[
 \sum_{x_B|_A=x_A}K_{A,B}(x_B\mid x_A)=1,
 \qquad
 \mathbf P_B=K_{A,B}\mathbf P_A.
\]
Thus every insertion span preserves the ordinary packet law exactly; R5 has
constant \(1\) before any later conditioning.

For \(A\subseteq B\subseteq C\), restriction gives a canonical identification
\[
 H_{A,B}\times_{X_B}H_{B,C}\cong H_{A,C},
\]
and the kernels obey
\begin{equation}\label{eq:physical-kernel-chain-rule}
 K_{A,C}(x_C\mid x_A)
 =
 K_{B,C}(x_C\mid x_C|_B)\,
 K_{A,B}(x_C|_B\mid x_A).
\end{equation}
This is the conditional-probability chain rule.  In particular, for two new
slots \(i,j\notin A\), both insertion orders give the same weighted span:
\begin{equation}\label{eq:physical-confluence-square}
 K_{A\cup\{i\},A\cup\{i,j\}}K_{A,A\cup\{i\}}
 =
 K_{A\cup\{j\},A\cup\{i,j\}}K_{A,A\cup\{j\}}
 =
 K_{A,A\cup\{i,j\}}.
\end{equation}
All longer confluence diagrams follow by associativity of
\eqref{eq:physical-kernel-chain-rule}.

The corresponding pull--push operator is
\begin{equation}\label{eq:weighted-span-linearization}
 \mathsf L_{A,B}e_{x_A}
 =
 \sum_{x_B|_A=x_A}
 K_{A,B}(x_B\mid x_A)e_{x_B}.
\end{equation}
Equations \eqref{eq:physical-kernel-chain-rule} and
\eqref{eq:physical-confluence-square} say exactly that weighted pullback
composition becomes matrix composition.  Hence
\[
 \mathsf L_{B,C}\mathsf L_{A,B}=\mathsf L_{A,C}
\]
strictly on the ordinary incidence module.

The raw lift count
\[
 \mathfrak q_{A,B}(x_A)
 :=\#\{x_B\in X_B:x_B|_A=x_A\}
\]
is the cardinality of the same middle fibre.  It is retained by the span even
though the conditional kernel has total mass \(1\).  After restricting to an
event \(E\), the new conditional marginal is bounded by
\(\mathbf P(E)^{-1}\) times the old one, exactly as in
\eqref{eq:conditional-marginal-ledger}.

This proves the insertion and confluence part of the finite physical
realization.  The
endpoint-invisible \(A_N\)-frontier of the arithmetic FSC carrier is not
identified with a \(\mathfrak q_{A,B}\)-element physical set.  It remains the Moore/root
coefficient decoration, and its comparison with the physical exact-support
line is the common-controller cospan
\eqref{eq:normalized-controller-comparison}.  The following sections construct
the root-face/rerooting span and compatibility of the rooted functors on cube
overlaps.

\section{The root-face span}

Take a measurable partition by physical cells \(Q\) of diameter \(O(\delta)\),
and let
\[
 Y_Q(T):=Y(T)\cap Q,
 \qquad
 w_{Q,T}:=|Y_Q(T)|,
 \qquad
 W_Q:=\sum_Tw_{Q,T},
 \qquad
 W:=\sum_QW_Q.
\]
Discard zero-mass cubes and tube cells.  On ordered \(d\)-tuples through one
cell, allowing repeated entries, define
\begin{equation}\label{eq:symmetric-root-tuple-law}
 \mathbf P_w(Q,T_0,\ldots,T_r)
 :=
 \frac{W_Q}{W}
 \prod_{j=0}^r\frac{w_{Q,T_j}}{W_Q}
 =
 \frac{\prod_{j=0}^r w_{Q,T_j}}{W\,W_Q^r},
 \qquad r=d-1.
\end{equation}
Equivalently, sample the coarse shaded incidence \((Q,T_0)\) with mass
\(w_{Q,T_0}/W\), and then sample \(r\) side tubes independently from the
weighted \(Q\)-star with probabilities \(w_{Q,U}/W_Q\).
Repeated entries lie on proper diagonal faces and their determinant is zero.

The expression in \eqref{eq:symmetric-root-tuple-law} is invariant under every
permutation of \((T_0,\ldots,T_r)\).  Moreover, every slot has exactly the
coarse shaded-incidence marginal:
\begin{equation}\label{eq:every-root-marginal}
 \sum_{T_0,\ldots,\widehat T_j,\ldots,T_r}
 \mathbf P_w(Q,T_0,\ldots,T_r)
 =\frac{w_{Q,T_j}}W.
\end{equation}
Conditional on \((Q,T_0)\), this is precisely the product law used in
\eqref{eq:rooted-packet-cochain}, with the weighted full \(Q\)-star as the side
law.

For \(j\in I^\circ\), let \(\rho_{0j}\) move \(T_j\) into the root slot and
retain the relative order of the other entries.  The doubly rooted middle
object
\[
 \mathscr R_{0j}
 :=\{(Q,T_0,\ldots,T_r)\}
\]
with its old-root and new-root legs is therefore a measure-preserving span.
The inverse span is \(\mathscr R_{j0}\), and permutation identities give the
root-exchange groupoid relations strictly.

The determinant marking is compatible with this span.  In the new order,
with \(T_j\) first, one has
\begin{equation}\label{eq:reroot-determinant-sign}
 \det(e_{T_j},e_{T_0},\ldots,\widehat e_{T_j},\ldots,e_{T_r})
 =
 (-1)^j\det(e_{T_0},\ldots,e_{T_r}).
\end{equation}
Applying \eqref{eq:root-wedge-projection} with root \(T_j\) shows that the new
transverse determinant is the old exact-support determinant multiplied by
the same orientation sign.  Thus \(\rho_{0j}\) carries
\(\mathfrak Z^{(d)}_{Q,T_0}\) to the \(T_j\)-rooted packet cochain, with no norm
or source loss.

Deleting the old root is now factored as
\begin{equation}\label{eq:root-face-factorization}
 \text{\(T_0\)-rooted \(I\)-packet}
 \xrightarrow{\ \rho_{0j}\ }
 \text{\(T_j\)-rooted \(I\)-packet}
 \xrightarrow{\ \mathrm{del}_{0}\ }
 \text{\(T_j\)-rooted \(I\setminus\{0\}\)-packet}.
\end{equation}
To avoid choosing \(j\), take the disjoint union of
\eqref{eq:root-face-factorization} over \(j\in I^\circ\) with weight \(1/r\).
Equation \eqref{eq:every-root-marginal} makes both endpoint laws ordinary.
The full-support determinant maps to zero in the exact-support quotient of
the proper face, while the remaining tuple is precisely the lower-support
physical carrier.  This is the missing root-face span.

If a branch condition retains only an event \(E\) of the symmetric tuple law,
the same construction has Radon--Nikodym loss at most
\(\mathbf P_w(E)^{-1}\).  Hence a retained mass
\(\delta^{o(1)}\) gives the declared \(\delta^{-o(1)}\) R5 loss.  No
root-degree or point-degree factor survives.

The local FSC functor is therefore complete on one cube: side insertion,
confluence, side deletion, root exchange, and the root face have compatible
weighted spans.  The remaining source-level interface is descent on overlaps
of cubes and shaded restrictions.

\section{Global incidence gluing and shaded restrictions}

Retain the partition and weights from the root-face construction.
The ordinary incidence space is
\[
 \Omega_{\rm inc}
 :=\{e=(Q,T,x):x\in Y_Q(T)\},
\]
with probability law
\begin{equation}\label{eq:ordinary-incidence-atom-law}
 d\mu_{\rm inc}(Q,T,x)
 :=\frac1W{\bf1}_{Y_Q(T)}(x)\,dx,
 \qquad
 m_a(x):=\sum_Ta(Q(x),T,x){\bf1}_{Y(T)}(x)
\end{equation}
for every incidence weight \(0\le a\le1\); here \(Q(x)\) is the unique cell
containing \(x\).  For the current support \(I=\{0,\ldots,r\}\), the global
ordered packet law is \eqref{eq:symmetric-root-tuple-law}; its coarse
one-slot marginal is \eqref{eq:every-root-marginal}.
To retain point-dependent hereditary restrictions, lift
\eqref{eq:symmetric-root-tuple-law} to
\begin{equation}\label{eq:global-lifted-incidence-packet-law}
 \begin{aligned}
 &d\widetilde{\mathbf P}_w
 \bigl(Q,(T_0,x_0),\ldots,(T_r,x_r)\bigr)\\
 &\qquad:=
 \mathbf P_w(Q,T_0,\ldots,T_r)
 \prod_{j=0}^r
 \frac{{\bf1}_{Y_Q(T_j)}(x_j)\,dx_j}{w_{Q,T_j}}.
 \end{aligned}
\end{equation}
Zero-mass tube cells have already been discarded.  Forgetting the points
recovers \(\mathbf P_w\), while every marked-incidence slot has marginal
\begin{equation}\label{eq:global-incidence-slot-marginal}
 d(\Pi_j)_*\widetilde{\mathbf P}_w(Q,T,x)
 =\frac1W{\bf1}_{Y_Q(T)}(x)\,dx
 =d\mu_{\rm inc}(Q,T,x).
\end{equation}
The lifted law is symmetric in the pairs \((T_j,x_j)\).  In particular
\(m_1\) is the ordinary shaded multiplicity and \(W=\int m_1\).

Since \(Q\) remains part of every state, all support deletions, insertions,
confluence maps, and root exchanges preserve the cube label.  The global
source-level functor is therefore the weighted direct sum
\begin{equation}\label{eq:global-local-fsc-sum}
 \Phi^{\rm rec}
 =\bigoplus_Q\Phi_Q^{\rm rec}.
\end{equation}

Now let \(0\le a(e)\le1\) be a hereditary shaded restriction.  Apply the
symmetric packet mask
\[
 A(Q,(T_0,x_0),\ldots,(T_r,x_r))
 :=\prod_{j=0}^r a(Q,T_j,x_j)
\]
to \(\widetilde{\mathbf P}_w\), and let
\(\alpha=\mathbb E_{\widetilde{\mathbf P}_w}A\) be its retained packet mass.  After
normalization,
\[
 d\widetilde{\mathbf P}_{w,A}
 =\alpha^{-1}A\,d\widetilde{\mathbf P}_w.
\]
For every slot \(j\) and every incidence set \(B\),
\begin{equation}\label{eq:global-restriction-rn}
 \widetilde{\mathbf P}_{w,A}\bigl((Q,T_j,x_j)\in B\bigr)
 \le
 \alpha^{-1}\mu_{\rm inc}(B).
\end{equation}
Thus a restriction with \(\alpha\ge\delta^{o(1)}\) has precisely the required
\(\delta^{-o(1)}\) R5 domination.  Because the mask is symmetric, rerooting
remains measure preserving.  Taking all support marginals from the single
restricted joint law makes the conditional kernels satisfy
\eqref{eq:physical-kernel-chain-rule}, so insertion and confluence also survive
the restriction.

Finally, root projection glues by ordinary pushforward.  In a chart rooted at
\(T_0\), let \(g_{Q,U}\) be the \(e_{T_0}\)-fibre density of
\(Y(U)\cap Q\).  Then
\begin{equation}\label{eq:global-projected-mass}
 \sum_Q\int_{e_{T_0}^\perp}g_{Q,U}(z)\,dz
 =
 \sum_Q|Y(U)\cap Q|
 =
 |Y(U)|.
\end{equation}
Overlapping projected cells contribute additively; none is selected or
discarded.  Equations \eqref{eq:every-root-marginal},
\eqref{eq:global-restriction-rn}, and
\eqref{eq:global-projected-mass} close the global source and marginal part of
the local FSC realization.

\section{Structural provenance and passive decorations}

The preceding construction already gives the physical support filtration.  We
now record, in one place, its compatibility with the all-support controller of
\cite{UnifiedTower}.  Let \(X_A\) be the finite weighted packet space in slots
\(A\subseteq I\), and set
\[
 \mathscr H_A^{\rm phys}:=\ell^2(X_A)\otimes\Lambda^{|A|}\R^d,
 \qquad
 \mathscr H_I^{\rm phys}:=\bigoplus_{A\subseteq I}\mathscr H_A^{\rm phys}.
\]
For \(i\notin A\), use the conditional kernel of the preceding section to
define
\begin{equation}\label{eq:physical-koszul-operator}
 \begin{aligned}
 \mathsf P_i(e_{x_A})
 &:=
 \sum_{x_{A\cup\{i\}}|_A=x_A}
 K_{A,A\cup\{i\}}(x_{A\cup\{i\}}\mid x_A)^{1/2}
 e_{x_{A\cup\{i\}}},\\
 \mathsf D_i(e_{x_A}\otimes\omega)
 &:=
 \sum_{x_{A\cup\{i\}}|_A=x_A}
 K_{A,A\cup\{i\}}(x_{A\cup\{i\}}\mid x_A)^{1/2}
 e_{x_{A\cup\{i\}}}\otimes
 \bigl(e_i(x_{A\cup\{i\}})\wedge\omega\bigr).
 \end{aligned}
\end{equation}
Both operators vanish after slot \(i\) has been inserted.  The kernel chain
rule \eqref{eq:physical-kernel-chain-rule} gives
\(\mathsf P_i\mathsf P_j=\mathsf P_j\mathsf P_i\); exterior creation therefore
gives
\[
 \mathsf D_i^2=0,
 \qquad
 \mathsf D_i\mathsf D_j+\mathsf D_j\mathsf D_i=0.
\]
Conditioning on a root \((Q,T_0)\) and contracting the top wedge with
\(e_{T_0}\) yields
\[
 \det\bigl(P_{e_{T_0}^{\perp}}e_{U_1},\ldots,
 P_{e_{T_0}^{\perp}}e_{U_{d-1}}\bigr),
\]
hence exactly the packet \(\mathfrak Z^{(d)}_{Q,T_0}\).

The normalized cobar--Rees controller in \cite{UnifiedTower} is generated by
singleton insertions, with nonsingleton differentials equal to their signed
commutators.  Sending its singleton generator in slot \(i\) to
\(\mathsf P_i\), and its exterior totalization to \(\mathsf D_i\), is therefore
a dg representation:
\begin{equation}\label{eq:global-fsc-physical-koszul-comparison}
 \boxed{
 \Xi_I:
 \operatorname{Tot}_{\Lambda}
 \bigl(\mathscr U^{\rm rec}_{I,N}\bigr)\otimes_{\mathbf Z}\R
 \longrightarrow
 \operatorname{End}(\mathscr H_I^{\rm phys}).}
\end{equation}
Its maximal bar word is the antisymmetrized product of the singleton
insertions, so its normalized physical image is the top determinant just
computed.  This proves the common-controller comparison
\eqref{eq:normalized-controller-comparison}; deletion and permutation
compatibility follow because both sides literally delete or relabel the same
operators.  No other part of the cobar or Moore presentation is used in the
scalar estimate.

For completeness, the passive typed decorations can be audited without
turning them into analytic hypotheses.  At a bounded state
\(\mathcal K^{(n)}_{\eta;I,\chi,\nu}\), the confluence matching fibre is
\begin{equation}\label{eq:confluence-matching-defect}
 M^{\rm conf}_{I,\chi}
 :=
 \holim_{\chi'\prec\chi}
 \mathcal K_{\eta;I,\chi',\nu}^{(n)},
 \qquad
 F^{\rm conf}_{I,\chi}
 :=
 \fib\left(
 \mathcal K_{\eta;I,\chi,\nu}^{(n)}
 \longrightarrow M^{\rm conf}_{I,\chi}
 \right),
\end{equation}
where \(\chi'\prec\chi\) ranges over strict predecessor histories.  The full
structural audit is
\begin{equation}\label{eq:complete-obstruction-list}
 \boxed{
 C^{\supp}_{I,\nu},\quad
 F^{\rm val}_{\eta;I},\quad
 F^{\rm depth}_{I,\nu},\quad
 F^{\rm conf}_{I,\chi},\quad
 \Theta^{\supp,\rm val},\quad
 \Theta^{\supp,\rm depth},\quad
 \operatorname{Gr}^{\rm Cart}_{n}.}
\end{equation}
Only the exact-support determinant, the history/load fibre, and the physical
root antecedent survive as scalar observables.

Indeed, over an incidence atom \(e\), let \(\mathcal A(e)\) be the finite fibre
of compatible decorations and let
\(\kappa(a\mid e)\ge0\) satisfy
\(\sum_{a\in\mathcal A(e)}\kappa(a\mid e)=1\).  Define
\begin{equation}\label{eq:physical-vertical-augmentation}
 \epsilon_{\rm phys}f(e)
 :=
 \sum_{a\in\mathcal A(e)}\kappa(a\mid e)f(e,a).
\end{equation}
Every transition in valuation, coefficient, or Cartan decoration fixes \(e\)
and preserves this normalized kernel.  It therefore acts as the identity
after \(\epsilon_{\rm phys}\); support deletion commutes with it because the
two maps act on separate coordinates.  Thus vertical and mixed defects create
no physical multiplicity.  Confluence histories are different: they change
the physical middle fibre and remain in the sum
\eqref{eq:rn-corrected-evaluation}.  In particular, the supplied \(q^2\)
provenance labels those histories but does not estimate their positive
root-cross overlap.

\section{Quantitative thickening of support latching}

Exact support must be separated from a near-face packet at the physical scale.
This is done inside the support filtration.  Put \(r=d-1\) and, for an ordered
rooted packet, define its Gram--Schmidt pivots by
\begin{equation}\label{eq:physical-pivots}
 p_1:=|s_1|,
 \qquad
 p_j:=
 \operatorname{dist}
 \bigl(s_j,\operatorname{span}(s_1,\ldots,s_{j-1})\bigr)
 \quad(2\le j\le r).
\end{equation}
Then
\begin{equation}\label{eq:determinant-pivot-product}
 |D_I(\mathbf U)|=\prod_{j=1}^r p_j.
\end{equation}
Fix a geometric threshold \(\kappa\in[\delta,c_d]\) for a small dimensional
constant \(c_d\).  The packet belongs to the
\(\kappa\)-broad exact-support stratum when every \(p_j\ge\kappa\).  Otherwise
assign it to the first failed pivot \(j\), with lexicographic tie breaking.
These \(r+1\) events are disjoint and exhaust the rooted packet law.

On the broad stratum,
\begin{equation}\label{eq:broad-determinant-floor}
 |D_I(\mathbf U)|\ge\kappa^r.
\end{equation}
On the \(j\)-th near-face stratum let \(P_{j-1}\) be orthogonal projection to
\(\operatorname{span}(s_1,\ldots,s_{j-1})\) and put
\(\widetilde s_j=P_{j-1}s_j\).  Then
\begin{equation}\label{eq:near-face-error}
 |s_j-\widetilde s_j|<\kappa,
 \qquad
 |D_I(s_1,\ldots,s_r)
      -D_I(s_1,\ldots,\widetilde s_j,\ldots,s_r)|
 \le\kappa.
\end{equation}
The second determinant on the left is zero.  Thus the packet is a
\(\kappa\)-thickening of the proper face \(I\setminus\{j\}\), together with the
linear-combination/history datum expressing \(\widetilde s_j\) in the earlier
span.

This statement also has the required affine form.  Write
\[
 v_j:=\langle e_{U_j},e_{T_0}\rangle e_{T_0}+\widetilde s_j.
\]
Then \(v_j\) belongs to the span of the root direction and the preceding side
directions, while
\(\lvert e_{U_j}-v_j\rvert=\lvert s_j-\widetilde s_j\rvert<\kappa\).
For \(\kappa<c_d\), \(v_j\ne0\) and
\(\lvert e_{U_j}-v_j/|v_j|\rvert=O(\kappa)\).  Since every participating tube
meets the current root cell, its unit-length segment lies in an
\(O(\kappa+\delta)\)-neighbourhood of the affine carrier generated by the root
direction and the surviving pivot directions.  Hence the failed pivot is
realized by the proper-support stage with carrier width
\begin{equation}\label{eq:near-face-geometric-scale}
 a_{\rm face}:=\delta+\kappa.
\end{equation}
The physical tube scale remains \(\delta\).  The width \(a_{\rm face}\) is a
parameter of the physical realization, not valuation depth, coefficient depth,
confluence multiplicity or obstruction height.

The carrier just constructed is initially a packet observable, not a cube
observable.  Give the broad stratum one formal mark \(c_{\rm br}\).  If
\(\omega=(Q,T_0,U_1,\ldots,U_r)\) lies instead in the \(j\)-th failed-pivot
stratum, quantize the normal and affine offset of the span in the preceding
paragraph and write
\begin{equation}\label{eq:packet-carrier-label}
 \mathfrak c_{\rm pkt}(\omega)
 :=(j,\text{normal cap},\text{offset bin}).
\end{equation}
In general \(\mathfrak c_{\rm pkt}\) depends on the completing tubes
\(U_1,\ldots,U_r\); it is not determined by \(Q\), nor by one terminal
incidence \((Q,T,x)\).  A cube label \(\mathfrak c_Q\) is available only after
one proves a coherent narrow-field statement: on the retained incidence set
there is a unit normal \(\nu_Q\), depending only on \(Q\), such that
\begin{equation}\label{eq:cube-carrier-coherence}
 |\nu_Q\cdot e_T|\lesssim\kappa
 \qquad\text{for every retained }(Q,T).
\end{equation}
Quantizing \(\nu_Q\) and \(\nu_Q\cdot x_Q\) at widths
\(\kappa\) and \(a_{\rm face}\) then gives \(\mathfrak c_Q\).  Every retained
tube through \(Q\) lies in the corresponding fixed enlarged slab, and all
terminal incidences at that cube have the same carrier label.  Thus
\eqref{eq:cube-carrier-coherence} is the precise passage from packetwise
latching to a moving cube-carrier field; it is a physical source theorem, not a
formal consequence of support deletion.

That passage can in fact be made on the intrinsic cube-narrow branch.  With
\(\mu_Q(T):=w_{Q,T}/W_Q\), define the unrooted direction Gram operator
\[
 G_Q:=\mathbb E_{T\sim\mu_Q}[e_Te_T^{\mathsf T}].
\]
The conditional form of the symmetric packet law
\eqref{eq:symmetric-root-tuple-law} and Cauchy--Binet give
\begin{equation}\label{eq:unrooted-cube-cauchy-binet}
 \mathbb E_{T_0,\ldots,T_{d-1}\sim\mu_Q}
 \left|\det(e_{T_0},\ldots,e_{T_{d-1}})\right|^2
 =d!\det G_Q.
\end{equation}
Choose \(\rho=\rho(\delta)\to0\) with
\(\rho^{-1}=\delta^{-o(1)}\), and put
\(\beta_{\rm Gr}:=\rho^d\kappa^{2d}\).  If
\(\det G_Q>\beta_{\rm Gr}\), then, since the squared determinant is at most
one, the packets satisfying
\begin{equation}\label{eq:cube-broad-packet-event}
 \left|\det(e_{T_0},\ldots,e_{T_{d-1}})\right|^2
 \ge \tfrac{d!}{2}\beta_{\rm Gr}
\end{equation}
have conditional mass at least \(c_d\beta_{\rm Gr}\).  This is a
source-subpower broad packet event for subpower \(\rho^{-1}\) and
\(\kappa^{-1}\), and its determinant is exactly the physical image of the top
exact-support word.

If \(\det G_Q\le\beta_{\rm Gr}\), let \(\nu_Q\) be a unit eigenvector for the
smallest eigenvalue.  Since \(\operatorname{tr}G_Q=1\),
\[
 \mathbb E_{T\sim\mu_Q}|\nu_Q\cdot e_T|^2
 =\lambda_{\min}(G_Q)
 \le(\det G_Q)^{1/d}
 \le\rho\kappa^2.
\]
Markov's inequality deletes at most a \(\rho\)-fraction of the conditional
incidence mass and leaves
\begin{equation}\label{eq:intrinsic-cube-normal-hardening}
 |\nu_Q\cdot e_T|\le\kappa.
\end{equation}
This is \eqref{eq:cube-carrier-coherence} with a normal depending only on
\(Q\).  Summing over cubes preserves the same deletion fraction.  Thus the
cube-first Gram split gives a source-compatible dichotomy: either retain a
source-subpower broad packet population, or descend on a coherent moving
cube-carrier field.  The packetwise first-failed-pivot rule above remains the
exact support-latching description, but it is not used to manufacture a
cube label.  After routing all packets in the intrinsic narrow branch to this
valid cube carrier, the label is a function of the cube.  The global
coarse-carrier count below then controls both packet-completion marks and
longitudinal fragmentation.

The first-failed-pivot partition is hereditary: after deleting tubes, shaded
incidences or a support slot, recompute \eqref{eq:physical-pivots} on the
restricted joint law.
Every surviving atom still belongs to exactly one pivot stratum.  Its
conditional pushforward is therefore a normalized kernel, so R5 loses only
the reciprocal mass of the retained stratum.  Since \(r\) is fixed, either the
broad stratum has the declared retained mass, or one proper-face stratum does,
with only a dimensional loss.

Consequently the analytic support step has the precise dichotomy
\begin{equation}\label{eq:quantitative-latching-dichotomy}
 \boxed{
 \begin{array}{ll}
 \text{broad exact support:}
   & |D_I|\ge\kappa^r
     \ \text{and use the determinant/root projection};\\[0.2em]
 \text{thickened proper face:}
   & p_j<\kappa
     \ \text{and descend to }I\setminus\{j\}
       \text{ in width }\delta+\kappa.
 \end{array}}
\end{equation}
Choose \(\kappa=\delta^{\varepsilon(\delta)}\), where
\(\varepsilon(\delta)\to0\) and \(\kappa\to0\).  Then the broad determinant
floor \(\kappa^r\) is source-subpower.  On the near-face branch the ordinary
directional pair argument has only
\[
 L_{a_{\rm face}}
 :=1+\log_+(a_{\rm face}/\delta)
 \le1+\log(1/\delta)
\]
in place of \(L\); it pays no power of the carrier width.  Since support drops
and \(d\) is fixed, the total number of such latching steps is finite and all
threshold, chart and logarithmic losses remain \(\delta^{-o(1)}\).  This is the
quantitative realization of support latching, not another motivic level.

\section{From fibre densities to strong projected shadings}

\subsection{Binary hardening and the planar base}

The measure-valued antecedent can be converted to a binary shading without
choosing one lift.  The conversion must be performed after aggregating all
lifts in one projected line cell.  Fix a root \(T_0\), write
\(\pi=\pi_{T_0}:\R^d\to W=e_{T_0}^{\perp}\), and work on a dyadic projected-speed
shell
\begin{equation}\label{eq:projected-speed-shell}
 \sigma\le |\pi e_U|<2\sigma,
 \qquad \delta\le\sigma\le1.
\end{equation}
Let \(C\) be an \(O(\delta)\) projected line cell, let
\(\mathbb U_C\) be its physical lifts, and regularize
\(m\le|\mathbb U_C|<2m\).  If
\(|Y(U)|\ge\lambda\delta^{d-1}\), put
\begin{equation}\label{eq:cell-fibre-density}
 n_C(x):=\sum_{U\in\mathbb U_C}{\bf1}_{Y(U)}(x),
 \qquad
 g_C(z):=\int_{\pi^{-1}(z)}n_C(x)\,d\mathcal H^1(x).
\end{equation}
Thus \(g_C=\sum_{U\in\mathbb U_C}g_U\); no new measure has been introduced.

One root-parallel fibre meets a lift in length \(O(\delta/\sigma)\), the
projected segment has \((d-1)\)-dimensional volume
\(O(\sigma\delta^{d-2})\), and Fubini gives
\begin{equation}\label{eq:cell-fibre-bounds}
 0\le g_C\lesssim m\frac\delta\sigma,
 \qquad
 \int_Wg_C\gtrsim\lambda m\delta^{d-1}.
\end{equation}
Choose a sufficiently small dimensional constant \(c>0\) and define
\begin{equation}\label{eq:binary-strong-shading}
 A_C:=\left\{z:
 g_C(z)\ge c\lambda m\frac\delta\sigma\right\}.
\end{equation}
The complement of \(A_C\) carries at most a small fixed fraction of the second
quantity in \eqref{eq:cell-fibre-bounds}.  The first quantity then yields
\begin{equation}\label{eq:strong-shading-bounds}
 \boxed{
 |A_C|\gtrsim\lambda\sigma\delta^{d-2},
 \quad
 g_C\gtrsim\lambda m\frac\delta\sigma\ \hbox{ on }A_C,
 \quad
 \int_{A_C}g_C\gtrsim\lambda m\delta^{d-1}.}
\end{equation}
This is the required binary strong shading.  The factors \(\sigma\) and
\(\sigma^{-1}\) compensate exactly.

The associated fibre cutoff is equally source-compatible.  Inside one
projected line cell only the forgotten scalar direction parameter remains.
For two full directions separated by angle \(\theta\), the intersection of
their unit \(\delta\)-tubes is \(O_d(\delta^d/\theta)\).  Direction separation
and dyadic summation therefore give
\begin{equation}\label{eq:root-fibre-cordoba-energy}
 \int_{\R^d}n_C(x)^2\,dx
 \lesssim_d m\delta^{d-1}L,
 \qquad L:=1+\log(1/\delta),
\end{equation}
uniformly in \(m\) and \(\sigma\).  To restart the three-dimensional argument
from the plane, the projected translation load must be kept rather than
discarded.  Put \(d=3\), write \(W=e_{T_0}^{\perp}\cong\mathbb R^2\), and on
the occupancy/speed shell set
\begin{equation}\label{eq:projected-capacity-weight}
 a_C:=\frac{m_C\delta}{\sigma}.
\end{equation}
If \(\mathcal C_\theta\) consists of the projected cells with direction
\(\theta\), direction separation upstairs gives
\begin{equation}\label{eq:projected-direction-capacity}
 0<a_C\lesssim1,
 \qquad
 \sup_\theta\sum_{C\in\mathcal C_\theta}a_C\lesssim1.
\end{equation}
Indeed \(\sum_{C\in\mathcal C_\theta}m_C\lesssim\sigma/\delta\): after the
projected direction and speed shell are fixed, only one \(\delta\)-separated
scalar direction parameter remains.  The factor \(a_C\) is also forced by
Fubini, since
\(a_C\sigma\delta=m_C\delta^2\), the original three-dimensional source mass.

The plane controls the weighted projected marginal.  Define
\begin{equation}\label{eq:weighted-planar-projected-load}
 W_{\rm pr}(z):=\sum_Ca_C{\bf1}_{A_C}(z).
\end{equation}
Parallel translated cells have bounded overlap, and two cells whose projected
directions make angle \(\vartheta\) intersect in area
\(O(\delta^2/(\vartheta+\delta/\sigma))\).  Summing first over translations by
\eqref{eq:projected-direction-capacity} and then over the
\(\delta/\sigma\)-separated projected directions gives the geometric planar
estimate
\begin{equation}\label{eq:weighted-planar-cordoba}
 \boxed{
  \int_W W_{\rm pr}^2
  \lesssim
  \frac{L_\sigma}{\lambda}\int_WW_{\rm pr},
  \qquad
  L_\sigma:=1+\log(\sigma/\delta).}
\end{equation}
No Fourier estimate is used here.  Hence, with
\(K_{\rm pr}=L_\sigma/(\tau\lambda)\) and
\(H_{\rm pr}:=\{W_{\rm pr}\le K_{\rm pr}\}\),
\begin{equation}\label{eq:weighted-planar-pruning}
 \int_{H_{\rm pr}}W_{\rm pr}
 \ge(1-O(\tau))\int_WW_{\rm pr},
 \qquad
 \|W_{\rm pr}{\bf1}_{H_{\rm pr}}\|_\infty
 \le K_{\rm pr}=\delta^{-o(1)}.
\end{equation}

An unweighted projected multiplicity bound cannot replace
\eqref{eq:weighted-planar-pruning}.  Tile a fixed planar square, for each of
\(\asymp\delta^{-1}\) directions, by \(\asymp\delta^{-1}\) parallel
\(\delta\)-tubes and give every cell weight \(a_C=\delta\).  Then
\eqref{eq:projected-direction-capacity} holds and \(W_{\rm pr}\simeq1\), but
\(\sum_C{\bf1}_{A_C}\simeq\delta^{-1}\) everywhere.  Any subsets \(B_C\)
with unweighted overlap at most \(K\) satisfy
\begin{equation}\label{eq:false-unweighted-projected-target}
 \sum_Ca_C|B_C|\lesssim\delta K,
 \qquad
 \sum_Ca_C|A_C|\simeq1.
\end{equation}
Thus subpower retention would force \(K\gtrsim\delta^{-1+o(1)}\).  The earlier
unweighted translation-stable target was therefore too strong; it erased the
fibre coordinate and was false already over the planar base.

\subsection{The fibre-aware three-dimensional descent}

The correct three-dimensional descendant keeps that coordinate.  Write
\(x=(z,t)\in W\times\mathbb R\) and put \(\ell:=\delta/\sigma\).  Over its
projected cell, each lift \(U\in\mathbb U_C\) contributes an interval of length
\(O(\ell)\) centred on an affine height function \(t=h_U(z)\).  Thus, up to
fixed enlargements,
\begin{equation}\label{eq:affine-fibre-disintegration}
 n_C(z,t)
 =\sum_{U\in\mathbb U_C}{\bf1}_{\{|t-h_U(z)|\lesssim\ell\}},
 \qquad
 \rho_C(z,t):=\frac{n_C(z,t)}{a_C}.
\end{equation}
On \(A_C\), equations \eqref{eq:cell-fibre-bounds}--%
\eqref{eq:strong-shading-bounds} say
\begin{equation}\label{eq:normalized-affine-fibre-marginal}
 c\lambda
 \lesssim\int_{\mathbb R}\rho_C(z,t)\,dt
 \lesssim1.
\end{equation}
In particular the actual lifted load is the weighted mixture
\begin{equation}\label{eq:affine-fibre-lifted-load}
 N_3(z,t):=
 \sum_C n_C(z,t){\bf1}_{A_C}(z)
 =\sum_Ca_C\rho_C(z,t){\bf1}_{A_C}(z),
\end{equation}
and therefore
\begin{equation}\label{eq:affine-fibre-projected-marginal}
 c\lambda W_{\rm pr}(z)
 \lesssim\int_{\mathbb R}N_3(z,t)\,dt
 \lesssim W_{\rm pr}(z).
\end{equation}
High unweighted
projected overlap can therefore be harmless when the short fibre intervals
are separated in \(t\); inequality \eqref{eq:false-unweighted-projected-target}
was detecting a projection artefact.

The exact stronger three-dimensional target is now fibre-aware.  After the
planar mask \(H_{\rm pr}\), one needs a single hereditary
\(H_3\subseteq\pi^{-1}(H_{\rm pr})\) such that
\begin{equation}\label{eq:affine-fibre-three-dimensional-refinement}
 \boxed{
  \int_{H_3}N_3
  \ge s_3\int_{\pi^{-1}(H_{\rm pr})}N_3,
  \qquad
  \|N_3{\bf1}_{H_3}\|_\infty\le K_3,
  \qquad
  s_3^{-1},K_3=\delta^{-o(1)}.}
\end{equation}
The planar calculation proves the marginal part of this statement completely.
The remaining assertion is an affine-fibre confluence packing estimate:
short intervals centred on the functions \(h_U\) may not align with
power-sized multiplicity on a source-subpower portion of their total mass.
Pairwise alignment is controlled by \eqref{eq:root-fibre-cordoba-energy}; the
unresolved sum is over repeatedly reused, source-correlated cycle cores.  In
the notation below, it is precisely the global packing requirement
\eqref{eq:dense-core-deletion-packing}, now with the affine heights retained.

The usual broad/narrow split can be expressed entirely in this fibre-aware
source.  For three lifts put
\begin{equation}\label{eq:affine-fibre-transversality}
 J(U_1,U_2,U_3)
 :=|\det(e_{U_1},e_{U_2},e_{U_3})|.
\end{equation}
The projected directions, speeds and affine height slopes together determine
\(J\).  On \(J\ge\kappa\), the common intersection has volume
\(O(\delta^3/\kappa)\).  After partitioning directions into
\(\kappa\)-caps, the number of transverse cap triples is
\(\kappa^{-O(1)}=\delta^{-o(1)}\), and endpoint multilinear Kakeya
\cite{BCT06,Guth10} controls the sum of the pointwise-broad descendants.  If no such
triple occurs at a high-load point, elementary linear algebra puts its active
directions in an \(O(\kappa)\)-neighbourhood of one plane.

The latter alternative is not yet a planar problem.  A unit tube with normal
direction component \(O(\kappa)\) still moves through a \(\kappa\)-thick slab;
collapsing that coordinate forgets the affine heights in
\eqref{eq:affine-fibre-disintegration}.  Partitioning into \(\kappa\)-slabs
and dilating the normal coordinate by \(\kappa^{-1}\) produces the same
three-dimensional affine-fibre problem at relative width
\begin{equation}\label{eq:affine-fibre-narrow-rescaling}
 \delta'=\frac{\delta}{\kappa},
\end{equation}
with an inherited carrier and confluence history.  Thus the narrow branch is
a scale iteration of \eqref{eq:affine-fibre-three-dimensional-refinement},
not an application of the bare planar union theorem.  To sum that iteration
one needs a source-compatible packing bound for the reused narrow carriers;
this is the same cycle-core issue, now organized by carrier depth.

This identifies what a proof ``from dimension two'' must add.  Planar tube
geometry gives \eqref{eq:weighted-planar-cordoba}--%
\eqref{eq:weighted-planar-pruning}; the full motivic filtration records the
affine heights, their confluence histories and cycle provenance; the imported
sticky theorem and sticky/non-sticky reduction supply the positive terminal
packing in \eqref{eq:terminal-sticky-carrier-gate}.  Together these give
\eqref{eq:affine-fibre-three-dimensional-refinement} with subpower retained
mass.  A near-full hereditary weighted version would be maximal-function
strength and is not supplied by these inputs.

\subsection{The marked three-to-four-dimensional handoff}

The next handoff shows exactly why the three-dimensional theorem cannot simply
be quoted again in dimension four.  Put now \(d=4\), so that
\(W=e_{T_0}^{\perp}\cong\mathbb R^3\), and retain the same dyadic
occupancy/speed shell.  The projected cells \(A_C\subset W\) are translated
three-dimensional tubes, while
\begin{equation}\label{eq:marked-three-dimensional-source}
 W_3(z):=\sum_Ca_C{\bf1}_{A_C}(z),
 \qquad
 N_4(z,t):=\sum_Ca_C\rho_C(z,t){\bf1}_{A_C}(z)
\end{equation}
obey
\begin{equation}\label{eq:four-dimensional-fibre-marginal}
 c\lambda W_3(z)
 \lesssim \int_{\mathbb R}N_4(z,t)\,dt
 \lesssim W_3(z).
\end{equation}
Thus the projected datum is a weighted multi-translation family carrying the
affine interval kernels \(\rho_C\).  The weights and parallel translations can
nevertheless be coalesced by the ordinary three-dimensional sticky output.

For each projected direction \(\theta\), after a harmless fixed
normalization put
\(A_\theta=\sum_{C\in\mathcal C_\theta}a_C\le1\) and independently choose at
most one cell in that direction, choosing \(C\) with probability \(a_C\) and
choosing none with probability \(1-A_\theta\).  Every realization \(\omega\)
is direction separated.  After rescaling length \(\sigma\) to one, apply the
three-dimensional sticky theorem and its sticky/non-sticky reduction to the
shadings \(A_C\) in that realization.  They give retained shadings
\(B_C^\omega\subseteq A_C\).  Put
\[
 M_\omega^{\rm ret}
 :=\sum_{C\,\mathrm{chosen}}{\bf1}_{B_C^\omega},
 \qquad
 M_\omega:= \sum_{C\,\mathrm{chosen}}{\bf1}_{A_C}.
\]

With constants uniform in \(\omega\),
\begin{equation}\label{eq:sampled-sticky-output}
 \int_WM_\omega^{\rm ret}
 \ge s_{\rm p}\int_WM_\omega,
 \qquad
 \|M_\omega^{\rm ret}\|_\infty\le K_{\rm p},
 \qquad
 s_{\rm p}^{-1},K_{\rm p}=\delta^{-o(1)}.
\end{equation}

Define the coalesced fractional shading
\[
 w_C(z):=\Pr_\omega\{C\ \text{ is chosen},\ z\in B_C^\omega\},
 \qquad
 \widetilde W_3(z):=\sum_Cw_C(z).
\]
Taking expectations in \eqref{eq:sampled-sticky-output} gives
\begin{equation}\label{eq:marked-sticky-projected-source}
 0\le w_C\le a_C{\bf1}_{A_C},
 \qquad
 \int_W\widetilde W_3\ge s_{\rm p}\int_WW_3,
 \qquad
 \|\widetilde W_3\|_\infty\le K_{\rm p}.
\end{equation}
This is one coalesced fractional physical subsource, not a sum of local masks.
Such Radon--Nikodym weights are native to the global source ledger; finite
rational approximation changes only subpower constants.  Thus ordinary
three-dimensional sticky geometry completely closes the weighted
multi-translation projected part.

What it does not see is the correlation among the fibre kernels.  Couple the
projected subsource to them by
\begin{equation}\label{eq:projected-retained-four-dimensional-source}
 N_4^{\rm pr}(z,t):=\sum_Cw_C(z)\rho_C(z,t),
 \qquad
 c\lambda\widetilde W_3(z)
 \lesssim\int_{\mathbb R}N_4^{\rm pr}(z,t)\,dt
 \lesssim\widetilde W_3(z).
\end{equation}

One sufficient remaining target is a single hereditary physical subsource
\(0\le\widetilde N_4\le N_4^{\rm pr}\), obtained by deleting complete affine
interval atoms rather than by separately chosen point masks for the cells, and
requires
\begin{equation}\label{eq:marked-fibre-frostman-capacity}
 \begin{split}
  &\int_{W\times\mathbb R}\widetilde N_4
   \ge s_{\rm f}\int_{W\times\mathbb R}N_4^{\rm pr},
   \qquad s_{\rm f}^{-1}=\delta^{-o(1)},\\
  &\sup_{\substack{z\in W,\ I\subset\mathbb R\\
                          |I|\ge\ell}}
   \frac{\displaystyle\int_I\widetilde N_4(z,t)\,dt}
        {|I|+\ell}
   \le K_{\rm f},
   \qquad K_{\rm f}=\delta^{-o(1)},
   \qquad \ell=\delta/\sigma .
 \end{split}
\end{equation}

This absolute capacity is unchanged by splitting a weight among parallel
translates and asks only for the multiplicity needed in the four-dimensional
conclusion; dividing additionally by \(\widetilde W_3(z)\) would impose an
unneeded stronger relative equidistribution.  Since every summand in
\eqref{eq:affine-fibre-disintegration} is an interval of length comparable to
\(\ell\), fixed enlargement of \(I\) gives, away from harmless interval
endpoints,
\begin{equation}\label{eq:fibre-capacity-implies-four-dimensional-cutoff}
 \widetilde N_4(z,t)
 \lesssim \ell^{-1}\!\int_{|u-t|\lesssim\ell}
                    \widetilde N_4(z,u)\,du
 \lesssim K_{\rm f}.
\end{equation}

Equations \eqref{eq:marked-sticky-projected-source} and
\eqref{eq:marked-fibre-frostman-capacity}, followed by the usual dyadic choice
of a projected mass shell if necessary, therefore give the four-dimensional
analogue of \eqref{eq:affine-fibre-three-dimensional-refinement}, retaining at
least a \(c\lambda s_{\rm p}s_{\rm f}\) fraction of the shell mass.  This is a
marked-fibre interface, not a second application of the unmarked
three-dimensional theorem.  Indeed, concentrating all normalized fibre mass
in one common interval of length \(\ell\) leaves \(\widetilde W_3\) unchanged
but makes \(N_4^{\rm pr}\simeq \widetilde W_3/\ell\) there.  Projection
alone has lost exactly the datum that
\eqref{eq:marked-fibre-frostman-capacity} measures.


For the actual scalar conclusion one can use a weaker quadratic form, which
interfaces directly with the confluence filtration.  Define the marked fibre
collision energy
\begin{equation}\label{eq:marked-fibre-collision-energy}
 \mathcal E_{\rm fib}
 :=\int_{W\times\mathbb R}
       \bigl(N_4^{\rm pr}(z,t)\bigr)^2\,dz\,dt .
\end{equation}
The sufficient estimate is
\begin{equation}\label{eq:marked-fibre-energy-gate}
 \boxed{
 \mathcal E_{\rm fib}
 \le A_{\rm fib}\int N_4^{\rm pr},
 \qquad A_{\rm fib}=\delta^{-o(1)}.}
\end{equation}
Indeed, for the global deletion budget \(\tau=\tau(\delta)\) put
\begin{equation}\label{eq:marked-fibre-size-biased-mask}
 H_{\rm fib}
 :=\{(z,t):N_4^{\rm pr}(z,t)
       \le A_{\rm fib}\tau^{-1}\},
 \qquad
 \widehat N_4:=N_4^{\rm pr}{\bf1}_{H_{\rm fib}} .
\end{equation}
Then size-biased Markov truncation gives
\begin{equation}\label{eq:marked-fibre-energy-pruning}
 \int_{H_{\rm fib}^c}N_4^{\rm pr}
 \le\frac{\tau}{A_{\rm fib}}\mathcal E_{\rm fib}
 \le\tau\int N_4^{\rm pr},
 \qquad
 \|\widehat N_4\|_\infty
 \le A_{\rm fib}\tau^{-1}=\delta^{-o(1)}.
\end{equation}
This is one global physical mask.  It gives the same absolute interval
inequality as \eqref{eq:marked-fibre-frostman-capacity}, with
\(K_{\rm f}=A_{\rm fib}\tau^{-1}\), and gives the desired cutoff
directly; unlike the atomwise Frostman route, it need not retain complete fibre
atoms.

\subsection{History recombination and the marked terminal source}

Disintegrate the convexified source by its already existing reduced confluence
histories,
\[
 N_4^{\rm pr}=\sum_\gamma v_\gamma,
 \qquad v_\gamma\ge0,
\]
using the same Radon--Nikodym-corrected joint law.  The prefix-tree calculation
\eqref{eq:lca-positive-charge}--\eqref{eq:r6-no-repeat-charge} applies
without change.  Conditional only on the established proper-predecessor
cutoff \(B_{\rm fib}^{\rm pred}\), it gives
\begin{equation}\label{eq:marked-fibre-proper-lca-energy}
 \sum_{a(\gamma,\gamma')\ne r_\ast}
   \int v_\gamma v_{\gamma'}
 \lesssim
 \delta^{-o(1)}B_{\rm fib}^{\rm pred}
 \int N_4^{\rm pr}.
\end{equation}

Consequently the only new quadratic term is
\begin{equation}\label{eq:marked-fibre-root-collision-gate}
 \boxed{
 \sum_{a(\gamma,\gamma')=r_\ast}
   \int v_\gamma v_{\gamma'}
 \lesssim
 \delta^{-o(1)}B_{\rm fib}^{\rm pred}
 \int N_4^{\rm pr}.}
\end{equation}

At this handoff the convexified three-dimensional cutoff
\eqref{eq:marked-sticky-projected-source}, together with the already
established strict-predecessor induction, gives
\(B_{\rm fib}^{\rm pred}=\delta^{-o(1)}\) on the retained shell.
Hence \eqref{eq:marked-fibre-proper-lca-energy} and
\eqref{eq:marked-fibre-root-collision-gate} imply
\eqref{eq:marked-fibre-energy-gate}.

Thus full motivic descent really does remove every proper-history fibre
collision, but it does not estimate the root-born correlated collision.
Equation \eqref{eq:marked-fibre-root-collision-gate} is the
three-to-four-dimensional specialization of
\eqref{eq:root-cross-energy-gate}; the structural \(q^2\) confluence
provenance labels its pair atoms but supplies no positive bound for their
physical overlap.

There is still a nontrivial unconditional packing gain inside this root term.
Regularize the marked lift atoms to comparable physical source mass, call the
resulting weighted cardinality \(N_{\rm fib}\), and select disjoint dyadic
carrier boxes \(Q\times J\subset W\times\mathbb R\) at scale
\(r\), where \(Q\) has side \(r\) and \(J\) has length \(r\).
The affine heights are uniformly Lipschitz on the speed shell, so a marked
unit tube crosses such a box in longitudinal length \(O(r)\).  If
\(\mathbb V_{Q,J}\) is the set of atoms assigned to the box, put
\begin{equation}\label{eq:marked-fibre-carrier-incidence}
 I_r^{\rm fib}
 :=\sum_{Q\times J}|\mathbb V_{Q,J}|.
\end{equation}
Exactly the incidence-source calculation in
\eqref{eq:tube-core-carleson-charge}--%
\eqref{eq:scale-average-core-charge}, now in four dimensions, gives
\begin{equation}\label{eq:marked-fibre-average-carrier-charge}
 \boxed{
 \overline{\mathfrak b}_{r}^{\rm fib}
 \lesssim
 \min\!\left\{1,
   \frac{rI_r^{\rm fib}}{\lambda N_{\rm fib}}\right\}.}
\end{equation}
Indeed the portion of one shaded \(\delta\)-tube inside an
\(r\)-box has volume \(O(r\delta^3)\), whereas its current source mass
is at least \(\lambda\delta^3\).  Fractional weights from
\eqref{eq:marked-sticky-projected-source} cancel on the two sides, so no
inverse power of \(a_C\) is introduced.

Consequently all fibre-carrier shells satisfying
\begin{equation}\label{eq:marked-fibre-nonsaturated-shells}
 \sum_r\frac{rI_r^{\rm fib}}
                   {\lambda N_{\rm fib}}=o(1)
\end{equation}
can be deleted inside the existing global source budget.  Failure of this
packing can occur only on scales with
\begin{equation}\label{eq:marked-fibre-sticky-saturation}
 I_r^{\rm fib}
 \gtrsim\frac{\lambda N_{\rm fib}}r .
\end{equation}
If saturation persists down a nested chain of boxes, the source has the same
number of available longitudinal slots as an all-scale sticky family, but now
the boxes retain the affine fibre mark.  This is the terminal marked-sticky
part of \eqref{eq:marked-fibre-root-collision-gate}.  The
three-dimensional theorem controls its projection and therefore already
accounts for \eqref{eq:marked-sticky-projected-source}; it does not control
which projected atoms occupy the same intervals \(J\).  A regular
tube--box bipartite reuse graph in which each tube meets
\(D\simeq\lambda/r\) selected boxes saturates
\eqref{eq:marked-fibre-sticky-saturation}, independently of its cycle-rank
statistics.  Thus neither unlabelled cycle rank nor the motivic history count
proves the missing positive estimate.

The saturated marked source has a final geometric split.  For four active
lifts define
\begin{equation}\label{eq:four-dimensional-marked-transversality}
 J_4(U_1,U_2,U_3,U_4)
 :=|\det(e_{U_1},e_{U_2},e_{U_3},e_{U_4})|.
\end{equation}
On \(J_4\ge\kappa\), four tubes have common intersection
\(O(\delta^4/\kappa)\); after partitioning directions into
\(\kappa\)-caps, endpoint multilinear Kakeya \cite{BCT06,Guth10} controls the whole broad
contribution with loss
\(\kappa^{-O(1)}=\delta^{-o(1)}\).  If no transverse quadruple occurs,
elementary linear algebra places all active directions in an
\(O(\kappa)\)-neighbourhood of one three-dimensional hyperplane.
Partition into \(\kappa\)-slabs and dilate the normal coordinate by
\(\kappa^{-1}\).  The relative width becomes
\begin{equation}\label{eq:four-dimensional-marked-narrow-rescaling}
 \delta'=\frac\delta\kappa,
\end{equation}
and the affine interval marks, confluence histories and saturated carrier
chain are all retained.  Thus this narrow branch is the same
four-dimensional marked-fibre problem at a coarser scale; it is not an
ordinary three-dimensional source.

After the broad terms and the non-saturated shells have been removed, the
exact terminal estimate is
\begin{equation}\label{eq:terminal-marked-sticky-gate}
 \boxed{
 |U(\mathbb V,Y;\mathscr Q_{\rm fib})|
 \ge\delta^{o(1)}\sum_{V\in\mathbb V}|Y(V)|,
 \qquad
 \mu(\mathbb V,Y;\mathscr Q_{\rm fib})
 \le\delta^{-o(1)}.}
\end{equation}
Here \(\mathscr Q_{\rm fib}\) is the inherited nested carrier chain
satisfying \eqref{eq:marked-fibre-sticky-saturation}; it is a mark on the
same physical union, not a second measure.  The low-multiplicity truncation
used after \eqref{eq:terminal-sticky-carrier-gate} closes the
four-dimensional scalar refinement directly, bypassing the need to establish
\eqref{eq:marked-fibre-root-collision-gate} on the unpruned source.
The uniform marked source-hereditary theorem of \cite{StickyContact4D}
applies with weights \(w_V\), arbitrary fractional restrictions
\(0\le\widetilde w_V\le w_V\), the affine fibre as a conditional mark, and
the \(\delta\)-dependent tree \(\mathscr Q_{\rm fib}\).  Its weighted
\(L^2\) estimate gives the volume bound, and its half-source-mass
low-multiplicity mask gives the second conclusion in
\eqref{eq:terminal-marked-sticky-gate} after the permitted relabelling of the
retained shading.  Hence the four-dimensional scalar
refinement is closed.

The strength of this input is essential.  The earlier four-dimensional
dimension-\(13/4\) estimate \cite{ChoudhuriSticky4D} permits a power loss
of order \(\delta^{3/4+o(1)}\) in union volume, rather than the
\(\delta^{o(1)}\) retention required in
\eqref{eq:terminal-marked-sticky-gate}.  It therefore cannot be inserted
into the subpower descent ledger.
More generally, if the terminal input has only
\begin{equation}\label{eq:terminal-marked-sticky-power-input}
 |U(\mathbb V,Y;\mathscr Q_{\rm fib})|
 \ge\delta^{\eta_4+o(1)}
       \sum_{V\in\mathbb V}|Y(V)|,
\end{equation}
then the same low-average-multiplicity truncation gives only
\begin{equation}\label{eq:terminal-marked-sticky-power-output}
 s_4\ge\delta^{\eta_4+o(1)},
 \qquad
 K_4\le\delta^{-\eta_4-o(1)}.
\end{equation}
At the level of finite-scale exponent bookkeeping, the known \(13/4\)
dimension corresponds to \(\eta_4=3/4\); the cited Hausdorff-dimension
statement is not being asserted here in the stronger uniform hereditary form
\eqref{eq:terminal-marked-sticky-power-input}.  In either form the
filtration preserves a positive exponent honestly; it does not convert a
power loss into a subpower one.

One can now test the proposed lower-dimensional recursion quantitatively.
Let \(\mathfrak R_4(\delta)\) denote the worst normalized root-fibre
collision ratio
\[
 \frac{\mathcal E_{\rm fib}}
      {B_{\rm fib}^{\rm pred}\int N_4^{\rm pr}}
\]
over the retained terminal marked sources at width \(\delta\).  The
four-linear broad estimate and one narrow rescaling give, with the presently
available carrier coalescing, at best the following schematic but
loss-explicit inequality:
\begin{equation}\label{eq:failed-four-dimensional-recursion}
 \mathfrak R_4(\delta)
 \lesssim \delta^{-o(1)}
 +\kappa^{-C}\mathfrak R_4(\delta/\kappa).
\end{equation}
The polynomial factor is not a count of motivic histories.  It is the physical
congestion of the possible narrow hyperplane carriers and their affine fibre
translations.  Replacing it by the number of structural operation words would
again confuse a geometric scale label with a confluence level.

Along a purely narrow chain choose scales
\(\delta_{j+1}=\delta_j/\kappa_j\) until
\(\delta_L\simeq1\).  Necessarily
\begin{equation}\label{eq:narrow-scale-product}
 \prod_{j=0}^{L-1}\kappa_j\simeq\delta,
 \qquad
 \prod_{j=0}^{L-1}\kappa_j^{-C}\simeq\delta^{-C}.
\end{equation}
Thus iteration of \eqref{eq:failed-four-dimensional-recursion} has a power
loss even when every three-dimensional projected source is closed perfectly.
Taking \(\kappa_j=\delta_j^{\varepsilon}\) merely distributes that
same power among more scales; it does not make the product subpower.

A successful \(3\to4\) recursion would therefore have to replace the
narrow term by
\begin{equation}\label{eq:contractive-marked-fibre-recursion}
 \mathfrak R_4(\delta)
 \lesssim\delta^{-o(1)}
 +G(\kappa)\mathfrak R_4(\delta/\kappa),
 \qquad
 \prod_jG(\kappa_j)=\delta^{-o(1)}.
\end{equation}
The required improvement from \(\kappa^{-C}\) to such a
scale-summable coefficient is exactly a coherent packing estimate for the
affine fibre carriers.  In the present notation it is
\eqref{eq:terminal-marked-sticky-gate}, or on the quadratic route
\eqref{eq:marked-fibre-root-collision-gate}.  For the common-interval
example after \eqref{eq:marked-fibre-frostman-capacity}, the fibre length
changes from \(\ell\) to \(\ell/\kappa\) under rescaling, so its
density obeys
\[
 \ell^{-1}=\kappa^{-1}(\ell/\kappa)^{-1}.
\]
Thus this example alone forces a factor at least \(\kappa^{-1}\) in any
recursion based only on the projected load.  It rules out a scale-summable
coefficient without an additional fibre-spreading estimate.

This also explains the fixed exponent in the known planebrush induction.
Its trilinear estimate already contains a \(\delta^{3/4}\) union-volume
factor, while its plany branch is reassembled across scales using stickiness
\cite{ChoudhuriSticky4D}.  The induction propagates the exponent
\(\eta_4=3/4\); it does not define a map
\(\eta_4\mapsto c\eta_4\) with \(c<1\) that could be iterated to
zero.  Hence the existing two- and three-dimensional theorems plus the
published four-dimensional planebrush recursion do not prove full
\eqref{eq:terminal-marked-sticky-gate}.

\subsection{Randomized fibre lifting and occupancy}

There is a sharper direct test of whether the three-dimensional convexification
can be lifted without the missing terminal input.  Restrict to a dyadic shell
\begin{equation}\label{eq:occupancy-weight-shell}
 a\le a_C<2a.
\end{equation}
In the sampling construction preceding
\eqref{eq:sampled-sticky-output}, define the unbiased lifted random load
\begin{equation}\label{eq:unbiased-random-fibre-lift}
 Z_\omega(z,t)
 :=\sum_{C\,\mathrm{chosen}}
     a_C^{-1}n_C(z,t){\bf1}_{B_C^\omega}(z).
\end{equation}
The factor \(a_C^{-1}\) is forced: taking expectation gives exactly the
convexified physical subsource,
\begin{equation}\label{eq:unbiased-random-fibre-expectation}
 \mathbb E_\omega Z_\omega=N_4^{\rm pr}.
\end{equation}
Without this factor the expectation retains only an \(a_C\)-fraction of
the original cell source and is unusable on a power-small occupancy shell.

For each realization the retained projected multiplicity is at most
\(K_{\rm p}\).  Hence Cauchy--Schwarz over its active projected cells,
followed by \eqref{eq:root-fibre-cordoba-energy}, gives
\[
 \int Z_\omega^2
 \lesssim K_{\rm p}a^{-2}
 \sum_{C\,\mathrm{chosen}}\int n_C^2
 \lesssim K_{\rm p}La^{-2}
 \sum_{C\,\mathrm{chosen}}m_C\delta^3 .
\]
In dimension four,
\(m_C\delta^3=a_C\sigma\delta^2\);
\eqref{eq:strong-shading-bounds} gives
\(|A_C|\gtrsim\lambda\sigma\delta^2\) and
\(g_C\gtrsim\lambda a_C\) on \(A_C\); and
\eqref{eq:sampled-sticky-output} gives
\(\sum_C|B_C^\omega|
 \ge s_{\rm p}\sum_C|A_C|\).  Combining these relations yields
\begin{equation}\label{eq:random-lift-occupancy-loss}
 \boxed{
 \int Z_\omega^2
 \lesssim
 \frac{K_{\rm p}L}
      {a\lambda^2s_{\rm p}}
 \int Z_\omega.}
\end{equation}
Jensen's inequality and
\eqref{eq:unbiased-random-fibre-expectation} therefore imply
\begin{equation}\label{eq:convexified-high-occupancy-energy}
 \mathcal E_{\rm fib}
 \lesssim
 \frac{K_{\rm p}L}
      {a\lambda^2s_{\rm p}}
 \int N_4^{\rm pr}.
\end{equation}
Consequently every shell satisfying
\begin{equation}\label{eq:subpower-occupancy-closure}
 a^{-1}=\delta^{-o(1)}
\end{equation}
obeys \eqref{eq:marked-fibre-energy-gate} using only the completed
three-dimensional sticky theorem, the geometric fibre energy and the global
source convexification.

This argument also identifies why the remaining low-occupancy shell is not a
bookkeeping artefact.  Removing \(a_C^{-1}\) from
\eqref{eq:unbiased-random-fibre-lift} loses an \(a_C\)-fraction of the
physical source; keeping it produces exactly the factor \(a^{-1}\) in
\eqref{eq:random-lift-occupancy-loss}.  Selecting each cell more often
creates \(a^{-1}\) parallel-translation color classes and returns the
same congestion after coalescing.  Thus the direct \(2/3\to4\) recursion
is now proved on the high-occupancy source, and its unresolved part is reduced
further to
\begin{equation}\label{eq:low-occupancy-marked-source}
 a=\delta^{\beta+o(1)}
 \quad\text{for some fixed }\beta>0,
\end{equation}
with correlated affine fibre intervals.  This is precisely the source on
which a new fibre-spreading gain must recover the factor \(a^{-1}\).

\subsection{The low-occupancy broad/narrow branch}

The low-occupancy source admits a source-compatible pointwise broad/narrow
split which sharpens this description.  Let \(\Gamma_x\) be the normalized
physical direction law of the active marked lift atoms at \(x=(z,t)\), with
all Radon--Nikodym weights retained, and put
\begin{equation}\label{eq:low-occupancy-pointwise-gram}
 G_x^{\rm fib}:=\int e_Ue_U^{\mathsf T}\,d\Gamma_x(U).
\end{equation}
Fix subpower parameters \(\theta,\kappa\) and split the ordinary
incidence mass according to
\begin{equation}\label{eq:low-occupancy-pointwise-broad-narrow}
 E_{\rm br}:=\{x:\det G_x^{\rm fib}\ge\theta\},
 \qquad
 E_{\rm nar}:=\{x:\det G_x^{\rm fib}<\theta\}.
\end{equation}
If \(E_{\rm br}\) carries a source-subpower fraction, pointwise
Cauchy--Binet gives a source-subpower fraction of transverse four-tuples at
the same physical point.  After the usual cap regularization, endpoint
multilinear Kakeya closes that descendant.  This avoids the false implication
from the averaged cube Gram operator noted in
\eqref{eq:cube-to-point-broad-transfer}: broadness is tested directly on
the common-point source.

On \(E_{\rm nar}\), choose measurably a unit eigenvector \(n(x)\) for
the least eigenvalue.  Since \(\operatorname{tr}G_x^{\rm fib}=1\),
\[
 \lambda_{\min}(G_x^{\rm fib})
 \le(\det G_x^{\rm fib})^{1/4}
 <\theta^{1/4}.
\]
Consequently Markov's inequality in the conditional direction law gives
\begin{equation}\label{eq:pointwise-normal-direction-pruning}
 \Gamma_x\{U:|\langle e_U,n(x)\rangle|>\kappa\}
 \le\frac{\theta^{1/4}}{\kappa^2}.
\end{equation}
Choose \(\theta^{1/4}\le\tau\kappa^2\) with
\(\tau^{-1},\kappa^{-1}=\delta^{-o(1)}\) and delete this
source-small part.  Quantize \(n(x)\) by a \(\kappa\)-net in
\(S^3\), of cardinality \(O(\kappa^{-3})=\delta^{-o(1)}\), and
retain one normal bin of maximal ordinary incidence mass.  This is one global
source selection.  On the resulting subpower descendant there is a fixed
unit normal \(n_\nu\) such that
\begin{equation}\label{eq:coherent-fixed-hyperplane-descendant}
 |\langle e_U,n_\nu\rangle|
 \lesssim\kappa
 \quad\text{for every retained active incidence atom.}
\end{equation}
Thus the pointwise-narrow part really does enter one coherent
three-dimensional hyperplane carrier before rescaling; no point-dependent
family of local masks is being summed.

However, applying
\eqref{eq:four-dimensional-marked-narrow-rescaling} to
\eqref{eq:coherent-fixed-hyperplane-descendant} returns the same
marked-fibre problem at width \(\delta/\kappa\).  If its new occupancy
is subpower it closes by \eqref{eq:subpower-occupancy-closure}; otherwise it
remains in \eqref{eq:low-occupancy-marked-source}.  Repeating the latter
alternative gives precisely
\eqref{eq:failed-four-dimensional-recursion}.  Hence all low-occupancy
mass that is pointwise broad at some scale closes; the only surviving
configuration is an all-scale chain of coherent pointwise-narrow hyperplane
carriers whose affine fibre density pays the unavoidable factor in
\eqref{eq:narrow-scale-product}.  This is a smaller and fully physical
form of the terminal marked-sticky obstruction.

Pulling the successive narrow normals back to the original coordinates gives
a useful flag dichotomy.  After rotating
\eqref{eq:coherent-fixed-hyperplane-descendant}, write
\(n_\nu=e_4\) and let
\[
 A_\kappa:=\operatorname{diag}(1,1,1,\kappa^{-1})
\]
be the narrow dilation.  If \(n'=v+\alpha e_4\) is the next narrow
normal in the dilated coordinates, its pullback constraint on an original
direction is
\begin{equation}\label{eq:pulled-back-narrow-normal}
 |\langle A_\kappa^{\mathsf T}n',e_U\rangle|
 \lesssim\kappa'.
\end{equation}
If \(|v|\gtrsim\kappa'\), this covector is quantitatively independent
of \(e_4\); the retained directions lie near the intersection of two
hyperplanes.  This is a proper two-dimensional direction plane and passes to
the already closed lower-support branch.  Such transverse flag drops can
occur only finitely many times.

If \(|v|\lesssim\kappa'\), then
\eqref{eq:pulled-back-narrow-normal} does not introduce a new flag
direction; instead it improves the original normal aperture from
\(\kappa\) to
\begin{equation}\label{eq:coherent-normal-aperture-product}
 |\langle e_U,e_4\rangle|
 \lesssim\kappa\kappa'.
\end{equation}
Hence a purely coherent narrow chain eventually places every direction within
\(O(\delta)\) of one fixed hyperplane.  That terminal flag is genuinely
three-dimensional.  Partition physical space into normal
\(\delta\)-slabs; each retained tube meets \(O(1)\) adjacent slabs,
and inside one slab its projection to \(e_4^\perp\) has only bounded
direction multiplicity.  Fubini strong-shading selection followed by the
three-dimensional sticky theorem therefore gives
\begin{equation}\label{eq:terminal-fixed-hyperplane-closure}
 |U(\mathbb V,Y)|
 \ge\delta^{o(1)}\sum_{V\in\mathbb V}|Y(V)|
\end{equation}
on this terminal fixed-hyperplane descendant.

The obstruction lies between
\eqref{eq:coherent-fixed-hyperplane-descendant} and
\eqref{eq:terminal-fixed-hyperplane-closure}.  Under
\(A_\kappa\), a round four-dimensional \(\delta\)-tube becomes an
anisotropic plank with transverse widths comparable to
\begin{equation}\label{eq:narrow-anisotropic-plank-widths}
 (\delta,\delta,\delta/\kappa).
\end{equation}
Enlarging the two thinner transverse axes to \(\delta/\kappa\) in order
to invoke an isotropic theorem pays a polynomial power of
\(\kappa^{-1}\) at every coherent scale.  This is the physical factor in
\eqref{eq:failed-four-dimensional-recursion}; it is not caused by choosing
new normals, since the coherent branch uses the same pulled-back flag.

Thus a lower-dimensional proof would close if one could establish the
following eccentricity-uniform bridge for the inherited coherent flag:
\begin{equation}\label{eq:flag-anisotropic-sticky-gate}
 \boxed{
 |U(\mathbb P,Y;\mathscr F)|
 \ge\delta^{o(1)}\sum_{P\in\mathbb P}|Y(P)|,
 \qquad
 \text{uniformly in all intermediate plank aspect ratios}.}
\end{equation}
Here \(\mathscr F\) is the single pulled-back normal flag, not a new
descent level.  The ordinary three-dimensional sticky theorem proves the
terminal aspect ratio
\eqref{eq:terminal-fixed-hyperplane-closure}, but it does not provide
uniform control while the normal width lies strictly between
\(\delta\) and one.  The flag-anisotropic gate is therefore a still
sharper formulation of the missing fibre-spreading estimate: proving it would
make the coherent aperture products telescope to the three-dimensional base
without the power loss in \eqref{eq:narrow-scale-product}.

\subsection{The flag tensorization obstruction}

The anisotropic gate has a lower-dimensional tensorization interpretation.
In coordinates adapted to the fixed flag, a coherent plank
atom is determined by a tangential three-dimensional tube \(T_c\) and a
two-dimensional strip \(S_j\) in the longitudinal--normal plane, with the
same longitudinal parameter.  Up to fixed enlargements,
\begin{equation}\label{eq:flag-fibre-product-incidence}
 P_{c,j}\simeq T_c\times_{[0,1]}S_j .
\end{equation}
The three-dimensional sticky theorem controls the \(T_c\) marginal and
the geometric planar C\'ordoba estimate controls the \(S_j\) marginal.
Neither marginal estimate controls a physical selector whose two outputs
depend jointly on \((c,j)\).

There is nevertheless a precise sufficient tensorization.  Suppose the
retained flag source has a positive fractional decomposition into modes
\(M_\eta\), and on every mode its physical kernel is endpointwise dominated
by
\[
 K_\eta^{T}(h\mid c)K_\eta^{S}(h'\mid j).
\]
If the two unary kernels have inherited RN constants
\(A_{\eta,T}\) and \(A_{\eta,S}\), require
\begin{equation}\label{eq:flag-fractional-mode-capacity}
 \boxed{
 \sum_\eta t_\eta A_{\eta,T}A_{\eta,S}
 \le\delta^{-o(1)},
 \qquad
 \sum_{\eta:(c,j)\in M_\eta}t_\eta\ge1
 \quad\text{on the positive flag source}.}
\end{equation}
Thus ``mode'' means an actual endpointwise factorizing positive source, not
an arbitrary rectangle containing nonexistent pairs.  The external-product
calculation \eqref{eq:external-span-pair-capacity} applies on each mode;
positive coalescence proves the RN/energy form of
\eqref{eq:flag-anisotropic-sticky-gate}.  This is the flag realization of
the fractional criterion \eqref{eq:fractional-mode-capacity}.

The general low-occupancy source need not admit a subpower decomposition of
this kind.  A finite model is the full grid
\((c_i,j_j)\), \(i,j\in\mathbf Z/R\), with deterministic joint selector
\begin{equation}\label{eq:flag-diagonal-correlation}
 F_{\rm flag}(c_i,j_j)=(i+j,i+j).
\end{equation}
Every endpointwise factorizing mode is a matching, whereas the \(R\) level
sets \(i+j=s\) are such matchings.  Hence its positive deterministic mode
rank is \(R\), by \eqref{eq:deterministic-mode-conflict-rank}.  Fractional
positive splitting does not improve this number.  With unit mode capacities,
the \(R\) level sets give the primal upper bound; in the dual
\eqref{eq:mode-conflict-fractional-dual}, assign weight \(1/R\) to every one
of the \(R^2\) grid vertices.  Every independent mode is a matching of size at
most \(R\), so this is dual feasible and has value \(R\).  Therefore
\begin{equation}\label{eq:flag-diagonal-fractional-obstruction}
 \boxed{\mathfrak r_{+}^{\rm mode}
 =\kappa_{+}^{\rm mode}=R.}
\end{equation}
Both unlabelled marginals nevertheless have uniform loads.  On an occupancy
shell one may have \(R\simeq a^{-1}\), reproducing the loss in
\eqref{eq:random-lift-occupancy-loss}.  The motivic operation word and the
structural \(q^2\) provenance can remain fixed while this physical
correlation rank grows.

Consequently the desired \(2D/3D\) tensorization is proved on the
external-product and fractional low-conflict flag sources.  Its unresolved
part is the same kind of correlated positive-selector/Hall obstruction as
the one isolated in \eqref{eq:coherent-physical-capacity-system}, now
realized geometrically by intermediate anisotropic planks.  This is a
sufficient geometric realization of that interface, not an identification
of the scalar and typed conditions.  Marginal lower-dimensional Kakeya
estimates alone cannot bound it.
The finite model does not assert failure of the geometric plank estimate;
it proves only that no argument using the two marginal estimates plus
positive (even fractional) coalescence can establish the general gate.  A
successful four-dimensional step must exploit additional joint plank
geometry or assume the corresponding capacity.

Two genuine pieces of this interface do reduce to the three-dimensional
output.  In a factorized diffuse-fibre block, partition the cells into modes
\(\mathcal C_\eta\) so that, on the retained source,
\begin{equation}\label{eq:factorized-diffuse-fibre-block}
 0\le\widetilde\rho_C\le\rho_C,
 \qquad
 \widetilde\rho_C(z,t)=r_\eta(z,t)
 \quad(C\in\mathcal C_\eta),
 \qquad
 \sum_\eta r_\eta(z,t)\le K_{\rm f}=\delta^{-o(1)}.
\end{equation}

Writing \(\widetilde W_{3,\eta}
=\sum_{C\in\mathcal C_\eta}w_C\), one has
\begin{equation}\label{eq:factorized-fibre-lift}
 \widetilde N_4=\sum_\eta r_\eta \widetilde W_{3,\eta}
 \le K_{\rm f}\widetilde W_3.
\end{equation}

Hence the projected subsource
\eqref{eq:marked-sticky-projected-source} lifts immediately, with
\(\|\widetilde N_4\|_\infty\le K_{\rm f}K_{\rm p}\).  In the
fibre-disjoint block, put
\(b_C=w_C\widetilde\rho_C\), so that
\(\widetilde N_4=\sum_Cb_C\).  If the one-cell pruning gives
\begin{equation}\label{eq:fibre-disjoint-lift}
 \|b_C\|_\infty\le K_0,
 \qquad
 \sup_{z,t}\#\{C:b_C(z,t)>0\}\le Q,
 \qquad K_0,Q=\delta^{-o(1)},
\end{equation}
then \(\widetilde N_4\le K_0Q\).  These are respectively the scalar shadows of
the external-product and conflict-free parts of the typed source below.  They
also satisfy \eqref{eq:marked-fibre-energy-gate}, since a pointwise bound
\(\widetilde N_4\le B\) gives
\(\int\widetilde N_4^2\le B\int\widetilde N_4\).

They do not settle the correlated marked block: the latter is precisely the
root collision \eqref{eq:marked-fibre-root-collision-gate}, or more strongly
the coherent atomwise capacity
\eqref{eq:marked-fibre-frostman-capacity}.  Claiming that the ordinary
three-dimensional sticky theorem already supplies either statement would hide
the four-dimensional Kakeya difficulty in
the word ``projection.''

\subsection{Comparison with a stronger projection hypothesis}

For comparison only, suppose one imposes the unavailable stronger condition
that there are \(B_C\subset A_C\) with
\begin{equation}\label{eq:lower-backend-density}
 |B_C|\ge\lambda_{d-1}\sigma\delta^{d-2},
 \qquad
 \left\|\sum_C{\bf1}_{B_C}\right\|_\infty\le B_{d-1}.
\end{equation}
Then the following factorized fibre calculation is valid, but
\eqref{eq:false-unweighted-projected-target} shows that it is not a general
route to \eqref{eq:affine-fibre-three-dimensional-refinement}.

Equations \eqref{eq:strong-shading-bounds} and
\eqref{eq:lower-backend-density} imply
\begin{equation}\label{eq:lifted-cell-incidence-mass}
 \int_{\pi^{-1}(B_C)}n_C
 =\int_{B_C}g_C
 \gtrsim\lambda\lambda_{d-1}m\delta^{d-1}.
\end{equation}
Under the physical incidence law on this lifted set, the size-biased mean is
bounded by
\begin{equation}\label{eq:size-biased-root-fibre}
 \frac{\int_{\pi^{-1}(B_C)}n_C^2}
      {\int_{\pi^{-1}(B_C)}n_C}
 \lesssim_d\frac{L}{\lambda\lambda_{d-1}}.
\end{equation}
Hence, for every fixed \(\zeta>0\), Markov deletion of at most an
\(O_d(\delta^\zeta)\) fraction of the same physical shaded incidence mass gives
\begin{equation}\label{eq:root-fibre-hard-cutoff}
 \boxed{
 n_C(x)\lesssim_{d,\zeta}
 \delta^{-\zeta}\frac{L}{\lambda\lambda_{d-1}}.}
\end{equation}
There is no power of the cell occupancy \(m\) or of the projected speed
\(\sigma\).

For the subpower formulation, choose a slowly varying
\(\zeta=\zeta(\delta)\downarrow0\) with
\(\zeta(\delta)\log(1/\delta)\to\infty\), slowly enough that the implicit
constant in \eqref{eq:root-fibre-hard-cutoff} is also \(\delta^{-o(1)}\), and
with \(\delta^\zeta\) below the per-level share of the deletion budget fixed
below.
Then the discarded fraction tends to zero and the cutoff loss is
\(\delta^{-o(1)}\).  This is the only diagonal choice of exponents used below.

If the lower-dimensional retained multiplicity is
\[
 m_{d-1}(z):=\sum_C{\bf1}_{B_C}(z)\le B_{d-1},
\]
then partitioning the lifts by \(C\) gives the pointwise identity
\begin{equation}\label{eq:base-fibre-multiplicity}
 m_d(x)
 \le m_{d-1}(\pi x)
 \max_{C:\,\pi x\in B_C}n_C(x),
\end{equation}
and consequently the rootwise cutoff
\begin{equation}\label{eq:graded-dimensional-handoff}
 \boxed{
 B_d\lesssim_{d,\zeta}
 B_{d-1}\,
 \delta^{-\zeta}\frac{L}{\lambda\lambda_{d-1}}.}
\end{equation}
All deletions here are measured in the ordinary physical incidence law and
sum over the disjoint cells \(C\).

Finally, if the rooted determinant moment satisfies
\(\mathbb E|s_1\wedge\cdots\wedge s_{d-1}|^2\ge a\), then
\begin{equation}\label{eq:determinant-speed-regularization}
 \Pr\left\{|s_1\wedge\cdots\wedge s_{d-1}|^2\ge a/2\right\}
 \ge\frac{a}{2-a}\ge\frac a2,
\end{equation}
and every projected speed on that event is at least \((a/2)^{1/2}\).
Thus when \(a=\delta^{o(1)}\), determinant thresholding and the finitely many
speed shells cost only source-subpower mass.  Equation
\eqref{eq:binary-strong-shading} and the planar estimate
\eqref{eq:weighted-planar-pruning} close the projected marginal in dimension
three.  The lift to a three-dimensional multiplicity cutoff is exactly
\eqref{eq:affine-fibre-three-dimensional-refinement}.  The factorized shortcut
\eqref{eq:graded-dimensional-handoff} is valid only under the unavailable
unweighted hypothesis \eqref{eq:lower-backend-density}, and is not used as an
unconditional closure.

\section{Filtered gluing on the global incidence source}

\subsection{The global first-failure decomposition}

Fix a deletion budget \(\tau(\delta)\to0\) so slowly that
\(\tau(\delta)^{-1}=\delta^{-o(1)}\) and the sum of the budgets over all
support, load and history levels still tends to zero.  For any retained
physical filtered stage \(X\), let \(m_X(x)\) be its Radon--Nikodym-corrected
point multiplicity and define its hereditary cutoff by
\begin{equation}\label{eq:hereditary-cutoff-definition}
 B_X
 :=
 \inf\left\{B:
 \begin{array}{l}
  \text{there is a measurable hereditary set }E_X\subset\R^d\text{ with}\\[0.1em]
  \displaystyle
  \int_{E_X^c}m_X\le\tau(\delta)\int m_X,
  \quad m_X{\bf1}_{E_X}\le B\ \text{a.e.}
 \end{array}
 \right\}.
\end{equation}
Thus every deletion is measured in the size-biased incidence mass of the same
global source.  Write \(\lambda_I\) for the hereditary shading-density
parameter at support \(I\); at one dimensional handoff,
\((\lambda_I,\lambda_{I\setminus\{0\}})\) is exactly the pair
\((\lambda,\lambda_{d-1})\) in
\eqref{eq:size-biased-root-fibre}.

The preceding calculation is not exported as one certificate for each root.
It realizes the fresh exact-support cofiber inside the same global filtered
source.  At support \(I\), the physical realization of the Reedy triangle is
\begin{equation}\label{eq:physical-reedy-triangle}
 L_I^{\rm phys}
 \longrightarrow
 K_I^{\rm phys}
 \longrightarrow
 C_I^{\rm phys}.
\end{equation}
The first term is assembled from proper support faces, while the last term is
the exact-support determinant block decomposed by
\eqref{eq:binary-strong-shading}--%
\eqref{eq:affine-fibre-lifted-load}.  Its positive three-dimensional cutoff is
the affine-fibre gate
\eqref{eq:affine-fibre-three-dimensional-refinement}, not the projected
unweighted shortcut.
The chosen root describes only the local chart used to evaluate
\(C_I^{\rm phys}\); it does not replace
\eqref{eq:physical-reedy-triangle} by a family of independently counted root
outputs.

Concretely, use the single symmetric joint law
\eqref{eq:symmetric-root-tuple-law}.  Every root contraction is a conditional
kernel, but its target pushforward carries the density \(\rho_\gamma\) from
\eqref{eq:rn-corrected-evaluation}.  Normalization of the kernel does not permit
one to erase this density.  The one-cell matching example corresponds exactly
to many histories with small chartwise multiplicity but a large summed
\(\rho_\gamma\)-corrected target multiplicity.

The stable cofiber notation alone does not provide a positive measure
decomposition, so one must choose it on the physical source.  Fix the support
order and the determinant thresholds used by the current branch.  Apply
Gram--Schmidt in that order to every packet: assign the packet to the first
maximal subset \(J\subseteq I\) whose declared pivots survive, using
lexicographic tie breaking, and assign it to \(J=I\) exactly when every pivot
survives.  Failed pivots are carried as the confluence/history fibre over
\(J\).  This gives a disjoint measurable partition of the single joint packet
law into its proper-face strata and its fresh exact-support stratum.  Changing
the support order only permutes finitely many such charts, and the symmetric
law averages them without changing a one-point marginal.

Let \(h_J(e)\) be the pushforward retention weight of an ordinary incidence
atom \(e\) from the \(J\)-stratum.  Because the strata partition one joint law,
\begin{equation}\label{eq:graded-source-partition}
 0\le h_J(e)\le1,
 \qquad
 \sum_{J\subseteq I}h_J(e)=1
\end{equation}
before later branch deletion.  Thus the support strata themselves do not
duplicate an incidence.  The history pushforward inside one stratum still has
to be descended.

Dyadically filter a history block by its target density:
\[
 \Gamma_k
 :=
 \{(\gamma,e):2^k\le\rho_\gamma(e)<2^{k+1}\},
\]
and let \(\operatorname{Gr}^{\rm conf}_k C_I^{\rm phys}\) be the corresponding
associated graded of \eqref{eq:confluence-matching-defect}.  Finite geometric
capacity gives only \(O_d(\log(1/\delta))\) nonempty load shells.  Put
\begin{equation}\label{eq:global-load-observable}
 \mathsf N_{I,k}(x)
 :=
 \sum_{\gamma}
 \mathcal E_\gamma
 \left(
 \rho_\gamma{\bf1}_{\Gamma_k} f_{\gamma,I}
 \right)(x),
\end{equation}
where \(f_{\gamma,I}\) is the shaded multiplicity observable on the
exact-support \(I\)-block.  The local calculations above estimate a fixed
history chart.  The required global confluence statement is the associated
graded estimate
\begin{equation}\label{eq:r6-confluence-energy-target}
 \boxed{
 \frac{\int \mathsf N_{I,k}(x)^2\,dx}
      {\int \mathsf N_{I,k}(x)\,dx}
 \lesssim
 \delta^{-o(1)}
 \frac{L}{\lambda_I\lambda_{I\setminus\{0\}}}
 B_{K_{I\setminus\{0\}}},}
\end{equation}
uniformly in \(k\).  This is R6.  Its left side contains every cross-history
term in the same load layer; proving only the fixed-root diagonal terms is not
enough.

Put
\[
 A_I
 :=
 \delta^{-o(1)}
 \frac{L}{\lambda_I\lambda_{I\setminus\{0\}}}
 B_{K_{I\setminus\{0\}}}.
\]
On the event \(\mathsf N_{I,k}>\tau(\delta)^{-1}A_I\),
\eqref{eq:r6-confluence-energy-target} gives
\[
 \int_{\{\mathsf N_{I,k}>\tau^{-1}A_I\}}\mathsf N_{I,k}
 \le
 \frac{\tau}{A_I}\int\mathsf N_{I,k}^2
 \le
 \tau\int\mathsf N_{I,k}.
\]
Thus, once \eqref{eq:root-cross-energy-gate} supplies R6, it produces exactly
the hereditary cutoff of
\eqref{eq:hereditary-cutoff-definition}; the factor \(\tau^{-1}\) is absorbed
in \(\delta^{-o(1)}\).

Once \eqref{eq:r6-confluence-energy-target} holds, the confluence filtration
and Cauchy--Schwarz over its \(O(\log(1/\delta))\) graded layers give the
support-step recurrence
\begin{equation}\label{eq:global-support-recurrence}
 \boxed{
 B_{K_I}
 \lesssim
 B_{L_I}
 +
 \delta^{-o(1)}
 \frac{L}{\lambda_I\lambda_{I\setminus\{0\}}}
 B_{K_{I\setminus\{0\}}}.}
\end{equation}
This is a descent statement: the lower-support term, the confluence/load
graded term and the successor are evaluated in one filtered object.  It is not
obtained by independently exporting a cutoff from each root.

\subsection{Support recurrence and the source-mass ledger}

The latching term has the same recombination issue.  The
first-failed-pivot rule of \eqref{eq:quantitative-latching-dichotomy} assigns
every packet to one proper face.  If \(m_{I\to J}\) is the multiplicity of the
stratum assigned to \(J\subsetneq I\), then
\[
 m_{L_I}(x)
 \le
\sum_{J\subsetneq I}m_{I\to J}(x).
\]

There are only \(O_d(1)\) support faces, but each face may contain many packet
completions.  Applying the lower-support cutoff independently to those packets
would be the forbidden local-to-global step.  They are therefore kept
positively aggregated inside each structural history node, not made into
additional vertices of the prefix tree below.  The packet-carrier mark is used
later only to disintegrate one common source before recombination.  On the
coherent narrow branch it coarsens to the unique cube label, and the entire
coarse label alphabet has subpower size by
\eqref{eq:global-coarse-carrier-count}.

The proper-predecessor charge controls all nonroot interactions, while
\eqref{eq:root-cross-energy-gate} controls their root interaction.  Conditional
on that gate, the width, shell and final hardening losses are source-subpower,
and the already established proper-face cutoffs give
\begin{equation}\label{eq:latching-cutoff}
 \boxed{
 B_{L_I}
 \lesssim_d
 \delta^{-o(1)}
 \max_{J\subsetneq I}B_{K_J}.}
\end{equation}

Substitution into \eqref{eq:global-support-recurrence} yields, under the same
root-cross hypothesis, the closed support recurrence
\begin{equation}\label{eq:closed-support-recurrence}
 \boxed{
 B_{K_I}
 \lesssim_d
 \delta^{-o(1)}
 \left(
  \max_{J\subsetneq I}B_{K_J}
  +
  \frac{L}{\lambda_I\lambda_{I\setminus\{0\}}}
  B_{K_{I\setminus\{0\}}}
 \right).}
\end{equation}
Thus no unknown \(I\)-support quantity remains on the right once the physical
root-cross energy is controlled.  If the retained
shading densities at the finitely many stages satisfy
\(\lambda_J\ge\delta^{o(1)}\), then \(L=\delta^{-o(1)}\) and fixed-dimensional
iteration gives
\begin{equation}\label{eq:support-induction-closure}
 B_{K_I}
 \lesssim_d
 \delta^{-o(1)}B_{K_2}.
\end{equation}
Without this density condition one keeps the displayed \(\lambda\)-factors;
the descent does not erase sparse shading.

The required branch and density bounds follow from one source-mass ledger.

Across a fixed-dimensional descent, the first-failed-pivot partitions have
\(O_d(1)\) cells, while the speed, occupancy, target-density and confluence-load
partitions each have \(O_d(L)\) cells.  Include the usual dyadic per-tube
conditional-density regularization, but keep it distinct from selection of a
tuple stratum.  The cube-Gram split adds only its broad/narrow alternative.
When a quantitative broad packet is retained, its conditional mass is at
least \(c_d\beta_{\rm Gr}\) by \eqref{eq:cube-broad-packet-event}, and its
reciprocal is source-subpower.  Quantizing a proper affine carrier at angular
scale \(\kappa\) and offset scale \(a_{\rm face}\ge\kappa\) contributes at
most \(C_d\kappa^{-C_d}\) labels over the bounded physical region.
Support drops at most \(d-2\) times, and
the history tree is positively recombined rather than branch-selected.  Hence
only \(O_d(1)\) dyadic partition types are selected along a branch, and the
product of their cardinalities satisfies
\begin{equation}\label{eq:source-partition-count}
 \mathcal P_\delta
 \le C_d L^{C_d}
       \beta_{\rm Gr}^{-C_d}\kappa^{-C_d}
 =\delta^{-o(1)}.
\end{equation}

Here is the normalization needed to compare the weighted Markov budgets with
ordinary source mass.  Every conditional history kernel has mass one, and
integration of its Radon--Nikodym-corrected pushforward recovers the mass of
its marked source.  A fixed source atom occurs in only \(O_d(1)\) reduced
insertion orders; load shells partition rather than duplicate a history, and
all selected chart refinements contribute at most a dimensional power of
\(\mathcal P_\delta\).  Hence, if \(\mathcal I_{\rm cur}\) is the ordinary
marked-incidence mass before the Markov cutoffs, then
\begin{equation}\label{eq:weighted-ordinary-mass-ledger}
 \sum_{\text{retained stages }(I,k)}
 \int\mathsf N_{I,k}(x)\,dx
 \lesssim_d
 \mathcal P_\delta^{C_d}\mathcal I_{\rm cur}.
\end{equation}

By positivity, the ordinary marked-incidence mass removed at a stage is no
larger than the corresponding removed filtered mass.  Choose the diagonal
budget above also so that
\(\tau(\delta)\mathcal P_\delta^{C_d}\to0\).  Summed over all stages,
\eqref{eq:weighted-ordinary-mass-ledger} then makes the ordinary source loss
\(o(\mathcal I_{\rm cur})\).

For each tube choose a heaviest conditional-density cell and let
\(a(Q,T,x)\in\{0,1\}\) be the resulting hereditary incidence mask.  Put
\[
 \bar a_{Q,T}
 :=\frac1{w_{Q,T}}\int_{Y_Q(T)}a(Q,T,x)\,dx
\]
on the nonzero tube cells.  Its ordinary
one-slot retained mass is
\[
 \beta
 :=\frac1W\sum_{Q,T}w_{Q,T}\bar a_{Q,T}
 \ge\mathcal P_\delta^{-1}.
\]
Writing
\(r_Q=W_Q^{-1}\sum_Tw_{Q,T}\bar a_{Q,T}\), the symmetric \(|I|\)-slot packet mask
has mass
\[
 \alpha_{\rm mask}
 =\sum_Q\frac{W_Q}{W}r_Q^{|I|}
 \ge
 \left(\sum_Q\frac{W_Q}{W}r_Q\right)^{|I|}
 =\beta^{|I|}
\]
by convexity.  Tuple-stratum selection and the at most \(d-2\) support drops
add only a dimensional power of \(\mathcal P_\delta\).  Consequently
\begin{equation}\label{eq:source-mass-ledger}
 \alpha\gtrsim_d\mathcal P_\delta^{-C_d}=\delta^{o(1)},
 \qquad
 \lambda_J^{-1}\lesssim_d
 \lambda_{\rm in}^{-1}\mathcal P_\delta^{C_d}=\delta^{-o(1)}
\end{equation}
whenever \(\lambda_{\rm in}^{-1}=\delta^{-o(1)}\).  The global Markov deletions
consume only an \(o(1)\) fraction of ordinary source mass by
\eqref{eq:weighted-ordinary-mass-ledger}, so they do not alter these bounds.
Thus \(\alpha\) and the later \(\lambda_J\) are outputs of the same source
filtration, not independent assumptions.

\subsection{Terminal unweighting}

It remains to pass from the weighted descended source to an ordinary
descendant once, at the end.

On the lifted \(|I|\)-slot incidence law, let
\(A_{\rm tot}\in[0,1]\) be the product of all actual deletion and retention
masks along the selected descent stratum; these masks are lifted to the
marked incidences and symmetrized in the surviving slots.  Confluence/history
multiplicities are not included in \(A_{\rm tot}\); they remain in the filtered
observable \(m_{K_I}\).  Choose in
\eqref{eq:hereditary-cutoff-definition} an admissible final cutoff event
\(G_I\subset\R^d\) with
\[
 m_{K_I}{\bf1}_{G_I}\le2B_{K_I},
\]
and include \(\prod_{j\in I}{\bf1}_{G_I}(x_j)\) in \(A_{\rm tot}\).
The corresponding deletion has already been charged to the global budget.

Let \(\Pi_j\) be the \(j\)-th marked-incidence slot projection.  Define the
averaged slot density
\begin{equation}\label{eq:terminal-rn-weight}
 h(e)
 :=
 \frac1{|I|}\sum_{j\in I}
 \frac{d(\Pi_j)_*(A_{\rm tot}\widetilde{\mathbf P}_w)}{d\mu_{\rm inc}}(e).
\end{equation}
Every unmasked slot marginal of \(\widetilde{\mathbf P}_w\) is \(\mu_{\rm inc}\), and
\(A_{\rm tot}\le1\).  Hence \(0\le h\le1\), while
\[
 \int h\,d\mu_{\rm inc}
 =\int A_{\rm tot}\,d\widetilde{\mathbf P}_w
 =\alpha.
\]

Consequently
\[
 \int_{\{h<\alpha/2\}}h\,d\mu_{\rm inc}\le\alpha/2,
\]
so, for \(E=\{h\ge\alpha/2\}\),
\[
 \int_Eh\,d\mu_{\rm inc}\ge\frac\alpha2,
 \qquad
 \mu_{\rm inc}(E)\ge\int_Eh\,d\mu_{\rm inc},
 \qquad
 E\subseteq\widehat G_I:=\{(Q,T,x):x\in G_I\}.
\]

With \(m_h\) and \(m_E\) defined by
\eqref{eq:ordinary-incidence-atom-law}, one has
\begin{equation}\label{eq:terminal-unweighting}
 \mu_{\rm inc}(E)\ge\frac\alpha2,
 \qquad
 m_E(x)\le\frac2\alpha m_h(x)
 \le\frac2\alpha m_{K_I}(x){\bf1}_{G_I}(x)
 \le\frac4\alpha B_{K_I}.
\end{equation}

The last inequality follows from positivity and normalization of every history
kernel: the filtered observable contains the retained incidence contribution
and possibly additional history load.  Therefore, after R6 has supplied the
filtered weighted estimate, a descent with
\(\alpha=\delta^{o(1)}\) produces a genuine ordinary-incidence descendant with
only \(\delta^{-o(1)}\) loss.  This single terminal threshold is the correct
globalization; no intermediate root-certificate assembly is used.

\section{Positive confluence recombination and the root gate}

\subsection{The history tree and least-common-ancestor charge}

The confluence estimate is global, but its globalization is carried by the
history filtration.  The unified tower's reduced history normal form uses the terminal
structural object once and retains every genuine intermediate frontier label
once.  Thus a marked reduced insertion chain is represented by a word
\[
 \gamma=(h_1,\ldots,h_s;I),
\]
in its actual insertion order, where \(I\) is written only once and no genuine
intermediate frontier is duplicated.  Different admissible orders remain
different marked histories.  Every prefix is again a legal reduced insertion
history, and deleting its last label is its strict predecessor.  This prefix
property, rather than an artificial global sorting of the labels, is the only
part of the reduced normal form used below.

Arrange the retained words by common prefix and append a terminal marker to
every complete word so that no leaf is a proper prefix of another.  The
retained histories on one load shell then form a finite rooted prefix tree
\(\mathscr T_{I,k}\): its leaves are the full histories, and every internal
vertex is a strict confluence predecessor.  A proper support face is removed
before entering this tree and belongs to \(L_I^{\rm phys}\).
The tree contains only these structural history labels.  In particular the
packet-carrier \(\mathfrak c_{\rm pkt}\) of
\eqref{eq:packet-carrier-label} is not made into another tree level: all of its
values remain positively aggregated inside the corresponding structural leaf.
They are disintegrated only when the root-cross measure is analysed.  This
placement is essential; applying a predecessor cutoff separately to tiny
carrier fragments would be the forbidden local-to-global argument.

On one confluence/load shell write
\[
 u_\gamma(x)
 :=
 \mathcal E_\gamma
 \left(
 \rho_\gamma{\bf1}_{\Gamma_k}f_{\gamma,I}
 \right)(x)\ge0.
\]
Then
\begin{equation}\label{eq:double-history-expansion}
 \int\mathsf N_{I,k}^2
 =
 \sum_{\gamma,\gamma'}
 \int u_\gamma(x)u_{\gamma'}(x)\,dx.
\end{equation}
For a vertex \(a\), let \(\operatorname{Desc}(a)\) be its descendant leaves and
put
\begin{equation}\label{eq:history-node-mass}
 W_a(x):=\sum_{\gamma\in\operatorname{Desc}(a)}u_\gamma(x).
\end{equation}
The children of \(a\) partition its descendant leaves, hence
\begin{equation}\label{eq:history-martingale-partition}
 W_a=\sum_{b\in\operatorname{ch}(a)}W_b.
\end{equation}
This is an identity of nonnegative functions on the original incidence source,
not a stable-category decomposition.  The unified tower supplies only the legal prefix
predecessors.  Positivity and additivity here come from the single physical
joint law: the cylinder sets of the child prefixes are disjoint and exhaust
the parent cylinder, and their conditional kernels obey
\eqref{eq:physical-kernel-chain-rule}.  More precisely, if \(a\) has
confluence type \(\chi(a)\prec\chi\), pushforward of that disjoint cylinder
partition gives, for the descended shading \(f_a\),
\begin{equation}\label{eq:node-predecessor-realization}
 W_a
 =\mathcal E_{\chi(a)}(f_a).
\end{equation}
Thus an internal node is an observable in a strictly lower filtration stage,
not a relabeling of the current full-history sum.

Let \(r_\ast\) be the root of \(\mathscr T_{I,k}\).  The valid induction
invariant applies to every \emph{proper} vertex \(a\ne r_\ast\):
\[
 W_a(x)\le B_I^{\rm pred},
 \qquad
 B_I^{\rm pred}
 :=
 \delta^{-o(1)}
 \frac{L}{\lambda_I\lambda_{I\setminus\{0\}}}
 B_{K_{I\setminus\{0\}}},
\]
after the usual source-small Markov deletion.  For a leaf this is
\eqref{eq:root-fibre-hard-cutoff}, with R5 absorbing its source-subpower
Radon--Nikodym density.  For a proper internal vertex it is the induction
hypothesis applied to \eqref{eq:node-predecessor-realization}; a fresh Cartan
vertex first passes to its root--Moore antecedent and then has smaller support.
This uniform form is applied only to the whole retained structural node, whose
total pushforward has source-subpower RN density.  It is never invoked
separately on its possibly power-many packet completions; those completions have
already been summed in \(W_a\), and their remaining pair interaction is treated
at the root below.

The well-founded analytic order for these proper predecessors is
\begin{equation}\label{eq:analytic-descent-order}
 (|I|,\ell(\chi),n),
\end{equation}
ordered lexicographically with support outermost, confluence length next and
Cartan height last.  A strict confluence predecessor lowers
\(\ell(\chi)\), a proper face lowers support, and a Cartan predecessor lowers
\(n\).  Supports of size at most one are trivial latching faces.  The minimal
nontrivial leaf state is \((2,0,0)\), where the planar C\'ordoba estimate
applies.  This establishes the cutoff on proper vertices; it does not establish
one on \(r_\ast\), because
\begin{equation}\label{eq:history-root-is-global-sum}
 W_{r_\ast}=\sum_\gamma u_\gamma=\mathsf N_{I,k}.
\end{equation}
Using the desired cutoff on \(r_\ast\) would therefore be circular.

For two leaves whose least common ancestor
\(a(\gamma,\gamma')\ne r_\ast\), define, with value zero when the denominator
vanishes,
\begin{equation}\label{eq:lca-positive-charge}
 c_{\gamma,\gamma'}(x)
 :=
 \frac{u_\gamma(x)u_{\gamma'}(x)}
      {W_{a(\gamma,\gamma')}(x)}.
\end{equation}
The proper-node cutoff gives
\begin{equation}\label{eq:r6-pair-charge}
 u_\gamma u_{\gamma'}
 \le B_I^{\rm pred}c_{\gamma,\gamma'}.
\end{equation}
For a proper internal vertex,
\[
 \sum_{a(\gamma,\gamma')=a}u_\gamma u_{\gamma'}
 =W_a^2-\sum_{b\in\operatorname{ch}(a)}W_b^2
 \le W_a^2,
\]
and a leaf contributes its diagonal \(W_a^2\).  Since the proper nodes at
each depth have disjoint descendant sets,
\begin{equation}\label{eq:r6-no-repeat-charge}
 \sum_{a(\gamma,\gamma')\ne r_\ast}
 c_{\gamma,\gamma'}(x)
 \le
 H_{I,k}\,\mathsf N_{I,k}(x),
\end{equation}
where \(H_{I,k}=\delta^{-o(1)}\) is the refined tree height.  Thus all
proper-predecessor pairs satisfy
\[
 \sum_{a(\gamma,\gamma')\ne r_\ast}
 \int u_\gamma u_{\gamma'}
 \lesssim
 \delta^{-o(1)}B_I^{\rm pred}\int\mathsf N_{I,k}.
\]

What remains is exactly the cross term born at the root:
\begin{equation}\label{eq:root-cross-energy}
 \mathsf X^{\rm root}_{I,k}
 :=
 \sum_{a(\gamma,\gamma')=r_\ast}
 \int u_\gamma u_{\gamma'}
 =
 \int\left(
  W_{r_\ast}^2-
  \sum_{b\in\operatorname{ch}(r_\ast)}W_b^2
 \right).
\end{equation}
Consequently the unconditional output of the tree calculation is
\begin{equation}\label{eq:r6-root-reduction}
 \int\mathsf N_{I,k}^2
 \lesssim
 \delta^{-o(1)}B_I^{\rm pred}\int\mathsf N_{I,k}
 +\mathsf X^{\rm root}_{I,k}.
\end{equation}
R6 follows precisely from the additional physical estimate
\begin{equation}\label{eq:root-cross-energy-gate}
 \boxed{
 \mathsf X^{\rm root}_{I,k}
 \lesssim
 \delta^{-o(1)}B_I^{\rm pred}
 \int\mathsf N_{I,k}.}
\end{equation}

\subsection{Fixed-carrier decomposition of the root term}

The physical content of this gate can be separated without changing the
history filtration.  Refine each terminal history leaf by the ordinary tube
to which its target incidence belongs.  Positivity and the disjoint
partition of the ordinary incidence source give
\begin{equation}\label{eq:terminal-tube-refinement}
 u_\gamma=\sum_Tu_{\gamma,T},
 \qquad u_{\gamma,T}\ge0,
\end{equation}
where zero summands are omitted.  This adds only a terminal label; it is not
a new motivic coordinate and it does not duplicate physical mass.  Hence the
root term has the disjoint decomposition
\begin{align}
 \mathsf X^{\rm root}_{I,k}
 &=\mathsf X^{\rm geom}_{I,k}
   +\mathsf X^{\rm dup}_{I,k},\label{eq:root-physical-split}\\
 \mathsf X^{\rm geom}_{I,k}
 &:={}
 \sum_{\substack{a(\gamma,\gamma')=r_\ast\\T\ne T'}}
 \int u_{\gamma,T}u_{\gamma',T'},\nonumber\\
 \mathsf X^{\rm dup}_{I,k}
 &:={}
 \sum_{\substack{a(\gamma,\gamma')=r_\ast\\T=T'}}
 \int u_{\gamma,T}u_{\gamma',T}.\nonumber
\end{align}

For a root child \(b\) define its terminal-tube load and the total terminal
load by
\begin{equation}\label{eq:root-child-terminal-load}
 D_{b,T}(x):=
 \sum_{\gamma\in\operatorname{Desc}(b)}u_{\gamma,T}(x),
 \qquad
 D_T(x):=\sum_{b\in\operatorname{ch}(r_\ast)}D_{b,T}(x).
\end{equation}

Then \(\sum_TD_T=\mathsf N_{I,k}\), and expanding only pairs of distinct
root children gives the exact positive identities
\begin{align}
 \mathsf X^{\rm dup}_{I,k}
 &=\int\sum_T
 \left(D_T^2-\sum_bD_{b,T}^2\right),
 \label{eq:duplicate-root-load-identity}\\
 \mathsf X^{\rm geom}_{I,k}
 &=\int\sum_{T\ne T'}
 \left(D_TD_{T'}-\sum_bD_{b,T}D_{b,T'}\right).
 \label{eq:geometric-root-load-identity}
\end{align}

Every bracket is nonnegative: it is the sum over \(b\ne b'\).  In particular,
the same-terminal-tube subgate follows from the source-compatible hardening
\begin{equation}\label{eq:terminal-history-hard-gate}
 \boxed{
 D_T(x)\lesssim
 \delta^{-o(1)}B_I^{\rm pred}
 \quad\text{for every retained }(x,T),}
\end{equation}
because
\[
 \mathsf X^{\rm dup}_{I,k}
 \le\int\sum_TD_T^2
 \le \|D_T\|_{L^\infty_{x,T}}\int\mathsf N_{I,k}.
\]

Thus the duplicate branch asks for a terminal-incidence history-load theorem,
not a tube-angle estimate and not a coefficient-depth lift.

Indeed, set
\[
 F_T(x):=\#\{b\in\operatorname{ch}(r_\ast):D_{b,T}(x)>0\}.
\]
Every root child retained in the hereditary branch-density range is a proper vertex, so
\(D_{b,T}\le W_b\le B_I^{\rm pred}\) by the induction invariant.  Therefore
\begin{equation}\label{eq:history-fragment-count-suffices}
 D_T(x)\le F_T(x)B_I^{\rm pred}.
\end{equation}
Consequently a source-compatible restriction with
\(F_T(x)\le\delta^{-o(1)}\) proves
\eqref{eq:terminal-history-hard-gate}.  This is a sufficient raw-fragment
criterion.  The common-source argument below replaces the raw history count by
an aggregate over the globally finite coarse-carrier alphabet; the needed cube
coherence on the narrow branch comes from the physical Gram split, not from
the formal tower.

\subsection{Common-source aggregation}

The full descent does remove multiplicity coming only from different history
charts on one physical carrier.  Let \(\vartheta(b)\) be the typed predecessor
of a root child, let \(\jmath(b)\in\{0,\ldots,r\}\) record its terminal slot
after the measure-preserving root exchange.  The structural child \(b\) does
not have one carrier label.  Instead disintegrate its cylinder by the measurable
packet mark:
\[
 C_{b,c}:=C_b\cap\{\mathfrak c_{\rm pkt}=c\},
 \qquad
 W_b=\sum_cW_{b,c},
 \qquad
 D_{b,T}=\sum_cD_{b,c,T}.
\]

For fixed \((\theta,j,c)\), put
\[
 \overline W_{\theta,j,c}
 :=\sum_{\substack{\vartheta(b)=\theta\\ \jmath(b)=j}}W_{b,c},
 \qquad
 \overline D_{\theta,j,c,T}
 :=\sum_{\substack{\vartheta(b)=\theta\\ \jmath(b)=j}}D_{b,c,T}.
\]
The sets \(C_{b,c}\) are disjoint restrictions of
\(\widetilde{\mathbf P}_{w,A}\), and for fixed \((\theta,j,c)\) every terminal
map is the same coordinate projection \(\Pi_j\).  Hence, for every ordinary
incidence set \(B\),
\begin{align}
 &\sum_{\substack{\vartheta(b)=\theta\\\jmath(b)=j}}
   (\Pi_j)_*\bigl({\bf1}_{C_{b,c}}\widetilde{\mathbf P}_{w,A}\bigr)(B)
   \nonumber\\
 &\qquad=
 \sum_{\substack{\vartheta(b)=\theta\\\jmath(b)=j}}
 \widetilde{\mathbf P}_{w,A}
 \bigl(C_{b,c}\cap\Pi_j^{-1}(B)\bigr)
 \le
 \widetilde{\mathbf P}_{w,A}\bigl(\Pi_j^{-1}(B)\bigr)
 \le\alpha^{-1}\mu_{\rm inc}(B).
 \label{eq:typed-aggregate-rn-proof}
\end{align}

Thus \eqref{eq:global-restriction-rn} proves the aggregate comparison
\begin{equation}\label{eq:typed-aggregate-rn-gate}
 \boxed{
 \frac{d\left(
   \sum_{\substack{\vartheta(b)=\theta\\\jmath(b)=j}}
   (\Pi_j)_*\bigl({\bf1}_{C_{b,c}}
                  \widetilde{\mathbf P}_{w,A}\bigr)
 \right)}{d\mu_{\rm inc}}
 \le\alpha^{-1}=\delta^{-o(1)}.}
\end{equation}

Let \(\overline W_{\theta,j}:=\sum_c\overline W_{\theta,j,c}\).  This is one
disjoint aggregate of proper structural nodes, so the predecessor invariant
applies to \(\overline W_{\theta,j}\); positivity then gives
\(\overline W_{\theta,j,c}\le\overline W_{\theta,j}\lesssim
B_I^{\rm pred}\) without applying induction separately to the carrier
fragment.  This is the exact gain: charts sharing a typed predecessor and
packet carrier are aggregated, but the carrier is not thereby made a function
of the terminal incidence.

No entropy hypothesis is needed at the chosen carrier resolution.  Let
\(\mathcal C_\kappa\) contain \(c_{\rm br}\) and all quantized proper affine
carriers occurring in \eqref{eq:packet-carrier-label}, over the finitely many
pivot positions and dimensions.  A \(\kappa\)-net in each relevant projective Grassmannian has
\(O_d(\kappa^{-C_d})\) elements.  Since the physical region is bounded and
\(a_{\rm face}=\delta+\kappa\ge\kappa\), the corresponding affine offsets
have the same type of bound.  Thus
\begin{equation}\label{eq:global-coarse-carrier-count}
 \boxed{\lvert\mathcal C_\kappa\rvert
 \le C_d\kappa^{-C_d}=\delta^{-o(1)}.}
\end{equation}
In particular, for a fixed terminal endpoint there are at most this many
carrier marks produced by packet completions; arbitrarily many packets with
the same mark remain in one common-source aggregate.  Combining this count with
\eqref{eq:typed-aggregate-rn-gate} and the fixed number of typed
predecessor/slot pairs gives directly
\begin{equation}\label{eq:carrier-fragment-load-bound}
 D_T(x)
 \le\sum_{\theta,j,c}\overline W_{\theta,j,c}(x)
 \lesssim\lvert\mathcal C_\kappa\rvert B_I^{\rm pred}
 \lesssim\delta^{-o(1)}B_I^{\rm pred}.
\end{equation}
This closes the duplicate subgate on both broad and narrow histories, without
summing independent chartwise estimates.

On the intrinsic cube-narrow branch,
\eqref{eq:intrinsic-cube-normal-hardening} makes the carrier label
\(\mathfrak c_Q\) a function of \(Q\).  For each tube, partition its ordinary
incidence mass among the labels in \(\mathcal C_\kappa\).  Given any slowly
vanishing source-subpower \(\tau\), delete the fragments having conditional
mass less than \(\tau/\lvert\mathcal C_\kappa\rvert\).  Their total mass on
each tube, and hence globally, is at most \(\tau\), while every retained
fragment has relative mass
\begin{equation}\label{eq:coarse-carrier-density-floor}
 \frac{\tau}{\lvert\mathcal C_\kappa\rvert}=\delta^{o(1)}.
\end{equation}
The hereditary lower-dimensional estimate therefore applies to each fixed
carrier with only subpower density loss.  The carrier pieces partition the
unary incidence mass; at each cube the label is unique, so their quadratic
outputs are disjoint up to the bounded overlap of the cube decomposition.
Consequently they recombine with at most the subpower factor in
\eqref{eq:global-coarse-carrier-count}.  Thus moving coarse carriers create no
additional gate.  This conclusion uses the scale
\(\kappa=\delta^{o(1)}\); it would be false for labels resolved at a fixed
power of \(\delta\).

The first summand is interaction between distinct physical tubes.  The
second is repeated history load over the same unmarked terminal tube and has
no angular separation.  Thus a concrete sufficient, and branchwise natural,
form of the root gate is
\begin{equation}\label{eq:two-root-subgates}
 \boxed{
 \mathsf X^{\rm geom}_{I,k}
 \lesssim\delta^{-o(1)}B_I^{\rm pred}\int\mathsf N_{I,k},
 \qquad
 \mathsf X^{\rm dup}_{I,k}
 \lesssim\delta^{-o(1)}B_I^{\rm pred}\int\mathsf N_{I,k}.}
\end{equation}
On one fixed projected line cell, the unweighted geometric kernel underlying
the first estimate is the C\'ordoba calculation
\eqref{eq:root-fibre-cordoba-energy}; after the load-shell and RN
regularizations, this is the fixed-carrier input.  Equations
\eqref{eq:global-coarse-carrier-count}--
\eqref{eq:coarse-carrier-density-floor} permit its carrierwise application on
the coherent cube-narrow branch.  The second estimate is not a
tube-intersection estimate: \eqref{eq:typed-aggregate-rn-proof} and
\eqref{eq:carrier-fragment-load-bound} close it by common-source aggregation.

The common-source aggregation is load-bearing.  If one keeps only individual
R5 bounds and forgets the disjoint cylinder source, a one-point model with
\(R\) leaf histories, all
having the same terminal tube, and \(u_\gamma=1\), every proper leaf has cutoff
\(1\), but
\[
 \mathsf N=R,\qquad
 \mathsf X^{\rm root}=R^2-R,\qquad
 \frac{\int\mathsf N^2}{\int\mathsf N}=R.
\]
Here \(D_T=R\), so the failure lies in the duplicate subgate.  This model
cannot satisfy \eqref{eq:typed-aggregate-rn-proof}: \(R\) unit-mass children
are not disjoint restrictions of one unit-mass common source.  It shows why
individual chartwise R5 is insufficient, not an obstruction to the
fixed-carrier aggregate argument.
Thus the duplicate subgate is closed globally, and the coherent cube-narrow
part of the geometric subgate is also closed.  What is not supplied by the
carrier count is the positive assembly of broad exact-support pairs of
distinct physical tubes into strict predecessors.  The unified tower names
those destinations, but the aggregate descent alone does not dominate their
physical pair measure.

A rotating-star construction is an obstruction only to a much finer carrier
alphabet.  Indeed, one root tube may meet locally narrow stars whose normals
visit power-many caps at angular resolution \(\eta=\delta^c\) with fixed
\(c>0\).  This shows that local Gram narrowness does not control fine-scale
normal variation.  It does not obstruct the present decomposition: at
resolution \(\kappa=\delta^{o(1)}\), the whole normal-cap space, not merely the
labels visited by one tube, has cardinality \(\kappa^{-O_d(1)}
=\delta^{-o(1)}\).  Refining the tube-normal covariance into spectral ranks or
principal-angle strata would therefore add non-load-bearing levels and is not
used.  The remaining issue is instead the broad exact-support pair assembly
described next.
\subsection{Positive routing to strict predecessors}

For the full-motive route, a sufficient remaining interface needs no new
descent coordinate.  Remove the coherent cube-narrow part of the ordered
terminal root-collision source, already controlled by
\eqref{eq:carrier-fragment-load-bound} and the fixed-carrier recombination
following \eqref{eq:coarse-carrier-density-floor}, and denote the broad
exact-support event by \(\mathscr B_{I,k}\).  The cube-Gram alternative and
the retained broad-packet stratum were fixed before the history tree was
formed, so this is a measurable mark of each retained history; no unrecorded
packet completion is being chosen here.

On a collision atom,
\({\bf1}_{\mathscr B_{I,k}}\) means that the common cube \(Q(x)\) is on the
Gram-broad alternative and both histories carry the retained broad-packet
mark.  What is still required is a positive
pair-realization into already existing strictly lower typed strata
\(\xi\prec(\eta;I,\chi,\nu,n)\).  Put
\begin{equation}\label{eq:root-collision-measure}
 d\Lambda^{\rm root}_{I,k}
 (\gamma,T;\gamma',T';x)
 :={\bf1}_{\mathscr B_{I,k}}
   {\bf1}_{\{T\ne T'\}}
   {\bf1}_{\{a(\gamma,\gamma')=r_\ast\}}
   u_{\gamma,T}(x)u_{\gamma',T'}(x)\,dx.
\end{equation}

Write \(\mathsf X^{\rm br}_{I,k}:=\|\Lambda^{\rm root}_{I,k}\|\).  A
positive routing consists of kernels
\[
 \mathcal R_\xi(dh,dh'\mid z)\ge0,
 \qquad z=(\gamma,T;\gamma',T';x),
 \qquad
 \sum_\xi\mathcal R_\xi(1\mid z)=1
\]
for \(\Lambda^{\rm root}_{I,k}\)-almost every \(z\).  Each kernel lands in a
strict predecessor \(\xi\), is the pair version of the R5-compatible unary
maps, and preserves their common lifted physical point.  It must vanish unless
\(\xi\) occurs in the finite typed support of the algebraic routing
\eqref{eq:algebraic-root-collision-routing}.  Put
\[
 \Lambda_\xi^{\rm rout}(A)
 :=\int\mathcal R_\xi(A\mid z)\,d\Lambda^{\rm root}_{I,k}(z).
\]

Write \(v_{\xi,h}(x)\ge0\) for the complete RN-corrected terminal-history
loads on the lifted predecessor and set
\begin{equation}\label{eq:predecessor-pair-measure}
 \mathsf N_\xi(x):=\sum_hv_{\xi,h}(x),
 \qquad
 d\Lambda_\xi^{\rm pred}(h,h';x)
 :=v_{\xi,h}(x)v_{\xi,h'}(x)\,dx.
\end{equation}

Thus \(\|\Lambda_\xi^{\rm pred}\|=\int\mathsf N_\xi^2\).  The exact remaining
target is
\begin{equation}\label{eq:root-collision-realization-target}
 \boxed{
 \Lambda_\xi^{\rm rout}
 \le \delta^{-o(1)}\Lambda_\xi^{\rm pred}
 \qquad\text{for every strict }\xi.}
\end{equation}

The kernels must be defined on the physical pair source and respect the
support, confluence and Cartan predecessors supplied by the full motivic
filtration; a stable-category decomposition alone is insufficient.

The unary destination ledger is already contained in the source-mass ledger,
provided equal typed predecessors are coalesced rather than copied once for
each source pair.  Indeed, the complete loads \(\mathsf N_\xi\) are existing
retained predecessor stages, so \eqref{eq:weighted-ordinary-mass-ledger} and
\eqref{eq:source-mass-ledger} give
\begin{equation}\label{eq:coalesced-predecessor-unary-ledger}
 \sum_\xi\int\mathsf N_\xi
 \lesssim_d
 \mathcal P_\delta^{C_d}\mathcal I_{\rm cur}
 \lesssim_d
 \mathcal P_\delta^{C_d}\alpha^{-1}
 \int\mathsf N_{I,k}
 =\delta^{-o(1)}\int\mathsf N_{I,k}.
\end{equation}
Creating pair-dependent copies of the same \(\xi\) would invalidate this
estimate and is not allowed.  Since
every \(\xi\) is strict, the already established predecessor R6 estimate gives
\[
 \|\Lambda_\xi^{\rm pred}\|
 \lesssim B_I^{\rm pred}\int\mathsf N_\xi.
\]
Summing \eqref{eq:root-collision-realization-target} gives
\[
 \mathsf X^{\rm br}_{I,k}
 \lesssim\delta^{-o(1)}B_I^{\rm pred}
          \int\mathsf N_{I,k},
\]
which, together with the closed cube-narrow contribution, is the root gate.


\subsection{The Hall capacity formulation}

For fixed predecessor loads, kernel existence has an exact capacity
formulation.  Let \(\mathscr Z^{\rm br}_{I,k}\) be the source of
\eqref{eq:root-collision-measure}, let
\(\mathscr Y^{\rm pred}:=\bigsqcup_\xi\mathscr Y_\xi^{[2]}\), and write
\(z\leadsto y\) when \(y\) is an allowed physical pair image of \(z\) in a
typed predecessor belonging to
\eqref{eq:algebraic-root-collision-routing}.  For
\(E\subseteq\mathscr Z^{\rm br}_{I,k}\), put
\(\mathcal A_\xi(E):=\{y\in\mathscr Y_\xi^{[2]}:
z\leadsto y\text{ for some }z\in E\}\).  Then
\eqref{eq:root-collision-realization-target}, with constant \(C_\delta\), is
equivalent to the measurable Hall inequalities
\begin{equation}\label{eq:pair-routing-hall-cut}
 \boxed{
 \Lambda^{\rm root}_{I,k}(E)
 \le C_\delta\sum_\xi
       \Lambda_\xi^{\rm pred}(\mathcal A_\xi(E))
 \quad\text{for every measurable }E,
 \qquad C_\delta=\delta^{-o(1)}.}
\end{equation}

Necessity follows by sending the routed mass of \(E\) into its allowed
neighbourhood.  For sufficiency, disintegrate over the preserved physical
point \(x\).  At fixed \(\delta\), the tube, retained-history and destination
alphabets in each fibre are finite.  Equation
\eqref{eq:pair-routing-hall-cut} for all measurable \(E\) is therefore
equivalent to its finite Hall cut inequalities for almost every \(x\): if a
fibre cut failed on a positive-measure set, one of the finitely many source
label subsets would fail measurably there and would give a violating global
\(E\).  On each fibre, max-flow/min-cut produces a flow with destination
capacities \(C_\delta\Lambda_{\xi,x}^{\rm pred}\).  Order the finite edges and
choose the lexicographically first feasible basic flow; its entries are
measurable functions of the fibre weights.  Integrating these flows and
disintegrating by the source weights gives the kernels \(\mathcal R_\xi\).

Thus unary R5 bounds only special one-slot cuts; the
unresolved analytic content is the full family of pair-source cuts in
\eqref{eq:pair-routing-hall-cut}.

The single choice \(E=\mathscr Z^{\rm br}_{I,k}\), together with the unary
ledger, is enough for the scalar root estimate.  The reason to require every
\(E\) is hereditary stability: subsequent support, shading and history
restrictions replace the source by such subsets, and the same routing must
remain valid without being reselected from target mass.


For a broad edge write
\[
 \mathcal S_{\rm alg}(e,e')
 :=\supp_{\rm typ}\mathfrak R_{\rm alg}(e,e'),
\]
the set of strict types having nonzero coefficient in the algebraic routing
\eqref{eq:algebraic-root-collision-routing}.  The direct full-filtration route
has a useful deterministic sufficient form which does not require a common
unary predecessor.  Let \(E_x^{\rm br}\) be the finite once-oriented fibre of
positive broad distinct-tube source pairs at \(x\), with the fixed order on
root children used to choose the orientation.  Let
\[
 \Theta_x:E_x^{\rm br}\longrightarrow
 \bigsqcup_\xi\mathscr Y_{\xi,x}^{[2]},
 \qquad
 \Theta_x(e,e')=(\xi,h,h'),
\]
be a measurable joint lift of the first nonzero typed obstruction, with
\(\xi\in\mathcal S_{\rm alg}(e,e')\).  It may depend jointly on \((e,e')\).

Writing \(u_e=u_{\gamma,T}\), it is enough to prove the weighted fibre bound
\begin{equation}\label{eq:joint-typed-fibre-capacity}
 \boxed{
 \sum_{\substack{(e,e')\in E_x^{\rm br}:\\
       \Theta_x(e,e')=(\xi,h,h')}}
 u_e(x)u_{e'}(x)
 \le C_\delta v_{\xi,h}(x)v_{\xi,h'}(x),
 \qquad C_\delta=\delta^{-o(1)},}
\end{equation}
for almost every \(x\) and every destination.  Pushforward by \(\Theta_x\)
then gives \eqref{eq:root-collision-realization-target} directly.  This is a
weighted physical codegree estimate for the joint FSC lift.  Neither the raw
cardinality of its fibre nor the structural \(q^2\) confluence provenance
proves \eqref{eq:joint-typed-fibre-capacity}.  When no deterministic lift has
this bound, the Hall formulation permits the source mass to split among
several admissible destinations.

The distinction is already visible with a genuine product source.  Give every
\((i,j)\in(\mathbf Z/R)^2\) mass \(R^{-2}\), and let a joint confluence map
send
\[
 (i,j)\longmapsto(i+j,i+j).
\]
Each target slot is still uniform, so both unary RN constants are one.  The
routed pair measure, however, has mass \(1/R\) at each diagonal point, whereas
the product of its unary target loads has mass \(1/R^2\) there.  The diagonal
cut therefore forces \(C_\delta\ge R\) in
\eqref{eq:pair-routing-hall-cut}.  Tensoring two unary R5 statements cannot
control a predecessor map which depends jointly on the two endpoints; the
missing input is precisely control of correlated confluence mass.

This example also marks the limit of the unary local-FSC data.  If, after a
strict type is selected, its physical pair map happens to be the external
product of two existing unary R5 maps, positivity and Fubini give the required
pair domination.  But the root-cross source is defined by
\(a(\gamma,\gamma')=r_\ast\): the two histories have no common strict prefix
other than the current root, whose cutoff is the assertion being proved.
Sending both endpoints to an arbitrary proper support face is no substitute,
because the broad exact-support determinant dies there.  Thus neither unary
R5 nor the structural pullback compatibility of the all-support comparison
supplies the required pair lift.  The general route must prove the joint
weighted-fibre estimate \eqref{eq:joint-typed-fibre-capacity}, or, more
generally, the splittable Hall estimate \eqref{eq:pair-routing-hall-cut}; no separate
unary-factorization hypothesis is imposed.

\subsection{The scalar cycle-space reduction}

There is a sharper scalar cycle-space reduction.  At a fixed \(x\), let the vertices be the
active children \(b\in\operatorname{ch}(r_\ast)\).  Write \(U_b(x)\) for the
broad endpoint load in that child and, for \(b\ne b'\), let
\(M_{bb'}(x)\) be the corresponding unordered broad distinct-tube collision
mass; explicitly,
\[
 \begin{aligned}
 U_b(x)&:=\sum_{\gamma(e)\in\operatorname{Desc}(b)}
          {\bf1}_{\rm br}(e;x)u_e(x),\\
 M_{bb'}(x)&:=
 \sum_{\substack{\gamma(e)\in\operatorname{Desc}(b),\
                  \gamma(e')\in\operatorname{Desc}(b')\\
                  T(e)\ne T(e')}}
 {\bf1}_{\rm br}(e;x){\bf1}_{\rm br}(e';x)u_e(x)u_{e'}(x).
 \end{aligned}
\]
where \({\bf1}_{\rm br}\) is the broad mark already fixed in
\eqref{eq:root-collision-measure}.  Then
\begin{equation}\label{eq:broad-child-edge-majorant}
 0\le M_{bb'}(x)\le U_b(x)U_{b'}(x),
 \qquad
 U_b(x)\le W_b(x)\le B_I^{\rm pred}.
\end{equation}

Let \(G_x^{\rm br}\) be the graph with an edge \(\{b,b'\}\) precisely when
\(M_{bb'}(x)>0\).  For a spanning forest \(F_x\), orient every tree away from
one root.  Each nonroot vertex is then charged once, so
\[
 \sum_{\{b,b'\}\in F_x}M_{bb'}(x)
 \le B_I^{\rm pred}\sum_bU_b(x).
\]

Thus the only uncharged mass lies on the independent chords.  Define the
weighted cycle excess
\begin{equation}\label{eq:broad-weighted-cycle-excess}
 \mathfrak C_{I,k}^{\rm br}
 :=2\int
 \min_{F_x\ {\rm spanning\ forest}}
 \sum_{\{b,b'\}\in E(G_x^{\rm br})\setminus E(F_x)}
 M_{bb'}(x)\,dx.
\end{equation}

The finite graph permits a measurable lexicographic choice of a minimizing
forest, and the preceding charge gives
\begin{equation}\label{eq:broad-cycle-rank-reduction}
 \boxed{
 \mathsf X_{I,k}^{\rm br}
 \le
 2B_I^{\rm pred}\int\mathsf N_{I,k}
 +\mathfrak C_{I,k}^{\rm br}.}
\end{equation}

Consequently, the scalar root gate closes if, at every retained support and
shading restriction, one has
\begin{equation}\label{eq:broad-cycle-budget}
 \boxed{
 \mathfrak C_{I,k}^{\rm br}
 \le\delta^{-o(1)}B_I^{\rm pred}
       \int\mathsf N_{I,k}}
\end{equation}

A coarser sufficient condition is
\begin{equation}\label{eq:broad-cycle-rank-condition}
 \operatorname*{ess\,sup}_x\beta_1(G_x^{\rm br})
 \le\delta^{-o(1)},
 \qquad
 \beta_1(G)=|E(G)|-|V(G)|+c(G),
\end{equation}
because a spanning forest has exactly \(\beta_1(G_x^{\rm br})\) chords and
each chord is bounded by
\(B_I^{\rm pred}\sum_bU_b(x)\).  The weighted form
\eqref{eq:broad-cycle-budget} is preferable: it allows many cycles provided
their physical load is small.  The raw cycle rank is monotone under inherited
graph deletion.  The normalized weighted budget must be verified uniformly
after every retained restriction, just like the root gate itself.

The forest decomposition is a localization of the difficulty, not a
cancellation theorem.  For example, suppose at one \(x\) there are \(n\)
active children, \(U_b=a=B_I^{\rm pred}\), and their terminal-tube supports
are pairwise disjoint.  Then \(G_x^{\rm br}=K_n\),
\(M_{bb'}=a^2\), and
\[
 2\min_{F_x}\sum_{E(K_n)\setminus E(F_x)}M_{bb'}
 =(n-1)(n-2)a^2,
\]
whereas the whole ordered root-cross mass is \(n(n-1)a^2\).  Thus a spanning
forest removes only \(O(n)a^2\) from an \(O(n^2)a^2\) interaction.  Any useful
cycle budget must come from independent geometric sparsity or small chord
weights; selecting an optimal forest does not create it.


Pairwise tube geometry alone cannot supply that sparsity in dimensions
\(d\ge3\).  Take a direction-separated bush of
\(N\simeq\delta^{-(d-1)}\) unit \(\delta\)-tubes through one common
\(\delta\)-ball and write \(m=\sum_T{\bf1}_T\).  The angular-shell count
\[
 \#\{T':\angle(T,T')\simeq\theta\}
 \lesssim(\theta/\delta)^{d-1},
 \qquad
 |T\cap T'|\lesssim\frac{\delta^d}{\theta},
\]
gives a contribution \(O(\delta\theta^{d-2})\) per dyadic shell and per
fixed \(T\).  Its shell sum is \(O(\delta)\) for \(d\ge3\), whereas it is
\(O(\delta\log(1/\delta))\) for \(d=2\).  In \(d\ge3\), the common ball
contributes \(N^2\delta^d\simeq N\delta\) after summation and gives the
matching lower order.  Thus
\begin{equation}\label{eq:bush-pairwise-power-loss}
 \boxed{
 \int m\simeq N\delta^{d-1}\simeq1,
 \qquad
 \int m^2\simeq N\delta\simeq\delta^{-(d-2)}
 \quad(d\ge3).}
\end{equation}

In \(d=2\) the same upper-bound calculation loses only a logarithm, which is
the C\'ordoba regime \cite{Cordoba77};
in \(d\ge3\) the loss is a genuine power and comes from broad angular shells.
Thus neither direction separation nor a repeated pairwise planar estimate
proves \eqref{eq:broad-cycle-budget} on the unpruned source.  Deleting only the
common \(\delta\)-ball does not fix this: the outer angular shells still have
large multiplicity.  A single bush can nevertheless be pruned cheaply.  If
\(R_\delta\to0\) and \(R_\delta^{-1}=\delta^{-o(1)}\), delete the concentric
\(R_\delta\)-ball.  This loses only an \(O(R_\delta)\)-fraction of every tube,
while direction packing gives
\[
 m(x)\lesssim (|x-x_0|+\delta)^{-(d-1)}
 \lesssim R_\delta^{-(d-1)}=\delta^{-o(1)}
 \qquad (|x-x_0|\ge R_\delta).
\]

What is missing in general is therefore a single source-compatible pruning or
redistribution of all dense cycle cores.  Precisely, on the current marked
incidence source one would need one measurable mask \(H\), with
\begin{equation}\label{eq:global-cycle-pruning-target}
 \boxed{\begin{gathered}
 \mu_{\rm inc}(H)
 \ge\delta^{o(1)}\mu_{\rm inc}(\text{current source}),\\
 \mathfrak C_{I,k}^{\rm br}[a]
 \le\delta^{-o(1)}B_I^{\rm pred}[a]
       \int\mathsf N_{I,k}[a]
 \quad\text{for every }0\le a\le{\bf1}_H.
 \end{gathered}}
\end{equation}

Here brackets mean the inherited restriction of every source, history and
predecessor load by the same mask.  This is one condition on all later
restrictions, not a new mask chosen for each bad core.  If
\eqref{eq:global-cycle-pruning-target} were proved at each of the finitely many
descent stages, the source ledger would retain subpower incidence mass and
\eqref{eq:broad-cycle-budget} would close the scalar proof.  For the hereditary
typed route, the corresponding exact target is the all-subset Hall condition
\eqref{eq:pair-routing-hall-cut} on the source induced by \(H\).  The bush
shows why deleting one visible core can be cheap; pairwise geometry gives no
simultaneous choice of \(H\) for all overlapping cores.  Selecting a different
deletion separately for each core would therefore not suffice.

\subsection{Pruning, bushes, and the maximal-function alternative}

The direct pruning attempts stop at exactly this point.  Bernoulli thinning of
root children does not help: in the complete-graph example above, retaining a
fraction \(p\) leaves about \(pn\) active children and cycle ratio of order
\(pn\), so a subpower ratio forces \(p\lesssim\delta^{-o(1)}/n\); for
power-sized \(n\) this discards a power fraction of the incidence source.
Deleting the region of the currently worst core is cheaper for a single bush,
but is not hereditary: a later restriction may expose a different dense
subcore.  Iterating that deletion would require a packing estimate of the form
\begin{equation}\label{eq:dense-core-deletion-packing}
 \mu_{\rm inc}\!\left(
   \bigcup_{\mathcal C}
   \{(T,x):x\in\Omega_{\mathcal C},\ T\in V(\mathcal C)\}
 \right)
 \le \tau(\delta)\mu_{\rm inc}(\text{current source}),
 \qquad \tau(\delta)\to0,
\end{equation}

\subsection{Average carrier packing}

The correct quantity for a multi-bush system is an average tube--core
Carleson charge, not a uniform bound for every tube.  Index the selected
cores by \(j\), let \(\Omega_j\) lie in a ball of radius \(R_j\ge\delta\),
and let \(V_j\) be the tubes charged to that core.  If every current shading
has density at least \(\lambda\), define
\begin{equation}\label{eq:tube-core-carleson-charge}
 \mathfrak b(T)
 :=\min\!\left\{1,
  C\lambda^{-1}\sum_{j:\,T\in V_j}R_j\right\}.
\end{equation}
A unit \(\delta\)-tube meets an \(R_j\)-ball in volume
\(O(R_j\delta^2)\).  Subadditivity of the union, followed by summation in
the ordinary shaded incidence law, therefore gives
\begin{equation}\label{eq:average-carrier-packing-criterion}
 \frac{
  \mu_{\rm inc}\!\left(
   \bigcup_j\{(T,x):T\in V_j,\ x\in\Omega_j\}\right)}
 {\mu_{\rm inc}(\text{current source})}
 \le
 \frac{\sum_T|Y(T)|\mathfrak b(T)}{\sum_T|Y(T)|}.
\end{equation}
Thus \eqref{eq:dense-core-deletion-packing} follows from smallness of the
average on the right; a few repeatedly charged spine tubes are allowed.

This closes the canonical multi-bush test on the current per-tube shading
shell.  Suppose \(M\) disjoint
\(R\)-balls support \(q\)-bushes, all bushes share one spine tube, and all
other spokes occur in only one bush.  There are \(1+M(q-1)\) tubes.  The
spine contributes at most one full tube of source mass, while the spokes lose
at most \(O(R/\lambda)\) each.  Hence
\begin{equation}\label{eq:spine-multibush-charge}
 \frac{\sum_T|Y(T)|\mathfrak b(T)}{\sum_T|Y(T)|}
 \lesssim
 \frac{R}{\lambda}+\frac1{Mq}.
\end{equation}
Choose the cutoff \(K\) slowly enough that the broad-bush deletion radius
\(R=K^{-1/2}\) satisfies \(R/\lambda\to0\); this is compatible with
\(K=\delta^{-o(1)}\).  Then arbitrarily many bushes on one spine are
harmless in the global source ledger.  The earlier objection based on the
worst tube was therefore too strong.

For a general system, first split the radii dyadically and write
\begin{equation}\label{eq:scale-core-incidence-count}
 I_R:=\sum_{j:\,R_j\simeq R}|V_j|
     =\sum_T d_R(T),
 \qquad
 d_R(T):=|\{j:R_j\simeq R,\ T\in V_j\}|.
\end{equation}
On the same per-tube shading shell, \(|Y(T)|\) is comparable across active
tubes.  Using \(\min(1,u)\le u\) in
\eqref{eq:tube-core-carleson-charge}, the contribution of the \(R\)-shell to
the average in \eqref{eq:average-carrier-packing-criterion} is at most
\begin{equation}\label{eq:scale-average-core-charge}
 \boxed{
  \overline{\mathfrak b}_R
  \lesssim
  \min\!\left\{1,\frac{RI_R}{\lambda N}\right\},
  \qquad N:=|\mathbb T|.}
\end{equation}
Consequently every collection of non-saturated shells satisfying
\begin{equation}\label{eq:nonsaturated-core-shell-sum}
 \sum_R\frac{RI_R}{\lambda N}=o(1)
\end{equation}
already obeys \eqref{eq:dense-core-deletion-packing}.  This conclusion is
purely geometric and uses no sticky theorem.  Failure can occur only on
scales for which
\begin{equation}\label{eq:sticky-core-saturation}
 I_R\gtrsim\frac{\lambda N}{R}.
\end{equation}
Since a unit tube meets at most \(O(R^{-1})\) disjoint \(R\)-balls, the last
display says that a positive source fraction is close to saturating its
available longitudinal core slots.  This is the precise carrier-tree meaning
of stickiness in the present source ledger.

The same argument does not close a general reused-spoke network.  Abstractly,
take a bipartite incidence graph between tubes and cores in which every core
has \(q\) tubes and every tube belongs to \(D\simeq\lambda/R\) cores.
Then \(\mathfrak b(T)\simeq1\) for essentially every tube, independently
of how large the graph cycle rank is.  Consequently neither cycle rank nor
the number of cores bounds the right side of
\eqref{eq:average-carrier-packing-criterion}.  Geometry must rule out,
factor, or estimate this regular reuse pattern.

In tube language that pattern is the sticky alternative: the same source
factors through nested thick carriers at many scales.  A multiscale
Frostman/Katz--Tao nonconcentration estimate \cite{KatzTao02} bounds the average charge in the
non-sticky branches; when the estimate persists at every scale, the required
terminal multiplicity bound is exactly the sticky Kakeya input
\cite{StickyKakeya,GuthWangZahl3D}, which in the present normalization is
\begin{equation}\label{eq:terminal-sticky-carrier-gate}
 \boxed{
 |U(\mathbb T,Y)|
 \ge\delta^{o(1)}\sum_{T\in\mathbb T}|Y(T)|,
 \qquad
 \mu(\mathbb T,Y)\le\delta^{-o(1)}.}
\end{equation}

We invoke at the same time the sticky/non-sticky factoring theorem from
\cite{GuthWangZahl3D}.  Its multiplicity factorization is performed before
passing to terminal carrier families and therefore coalesces overlapping
physical images; it is not a sum of independently chosen local masks.  The
ball-like non-saturated factors are charged by
\eqref{eq:scale-average-core-charge}; the reduction's slab and plank factors
are controlled by its stated multiplicity estimates; and every all-scale
sticky terminal factor is covered by
\eqref{eq:terminal-sticky-carrier-gate}.

Indeed, put \(m=\sum_T{\bf1}_{Y(T)}\),
\(\bar m=\int m/|U|\), and
\(H_{\rm st}:=\{x\in U:m(x)\le2\bar m\}\).  Markov's inequality in Lebesgue measure and
the fact that \(m\ge1\) on \(U\) give
\[
 |H_{\rm st}|\ge\tfrac12|U|,
 \qquad
 \int_{H_{\rm st}}m\ge|H_{\rm st}|
 \ge\tfrac12\delta^{o(1)}\int m,
 \qquad
 \|m{\bf1}_{H_{\rm st}}\|_\infty\le2\delta^{-o(1)}.
\]
Thus the low-multiplicity mask gives
\eqref{eq:affine-fibre-three-dimensional-refinement}.  This closes the scalar
three-dimensional affine-fibre gate.  The imported terminal estimate is not a
consequence of \eqref{eq:weighted-planar-cordoba} and is not used for the typed
Hall routing.

The motivic carrier history records the
factorization and makes the average in
\eqref{eq:average-carrier-packing-criterion} hereditary, but it supplies no
positive estimate for the regular sticky network.  Thus the multi-bush
calculation removes spine reuse as an obstruction and isolates the remaining
three-dimensional input to sticky carrier packing, rather than to cycle rank
alone.

Here the family of regions must be source-compatible and must hit every inherited
\(\delta^{-o(1)}\)-dense cycle core \(\mathcal C\).  The two-tube intersection
bound controls one selected \(\Omega_{\mathcal C}\), as in the bush, but does
not sum \eqref{eq:dense-core-deletion-packing} over cores met successively by
the same tube.  A disjoint-ball selection does not repair the sum: at radius
\(R\), one unit tube may meet \(O(R^{-1})\) disjoint selected balls and hence
lose \(O(1)\), rather than \(o(1)\), of its length.  Thus the greedy proof of
\eqref{eq:global-cycle-pruning-target} needs a new Carleson-type packing input;
neither the source ledger nor planar C\'ordoba supplies it.

\subsection{Maximal and three-dimensional set inputs}

There is, however, a clean conditional closure by the Kakeya maximal
conjecture; see \cite{Wolff95,KatzTao02,ZahlKakeyaSurvey} for the classical
maximal-function framework and its relation to the set problem.  Put
\(p_d=d/(d-1)\).  After the existing load and per-tube
shading regularizations, its weighted hereditary form says that, for every
current direction-separated terminal shell and every collection
\(0\le f_T\le{\bf1}_{Y(T)}\),
\begin{equation}\label{eq:hereditary-weighted-kakeya-maximal}
 \boxed{
  \int\left(\sum_Tf_T\right)^{p_d}
  \le K_\delta\int\sum_Tf_T,
  \qquad K_\delta=\delta^{-o(1)}.}
\end{equation}
This is the usual conjectural maximal estimate in the form needed here.  To
pass from tube indicators to \eqref{eq:hereditary-weighted-kakeya-maximal},
decompose the coefficients dyadically.  On a retained cell
\(f_T\simeq a{\bf1}_{Y(T)}\), the ordinary maximal bound is applied to the
active tubes, while
\(|Y(T)|\ge\lambda|T|\), \(\lambda^{-1}=\delta^{-o(1)}\), converts their total
tube volume to at most \(\lambda^{-1}\sum_T|Y(T)|\); the factors \(a^{p_d}\le
a\) and the subpower number of load cells give the displayed form.

Indeed, \eqref{eq:carrier-fragment-load-bound} gives, on the current shell,
\[
 0\le D_T(x)\le A_\delta B_I^{\rm pred}{\bf1}_{Y(T)}(x),
 \qquad A_\delta=\delta^{-o(1)}.
\]
Set
\[
 f_T:=\frac{D_T}{A_\delta B_I^{\rm pred}},
 \qquad M:=\sum_Tf_T,
\]
and fix the deletion budget \(\tau(\delta)\) from the source ledger.  With
\begin{equation}\label{eq:maximal-pruning-threshold}
 L_\delta:=\left(\frac{K_\delta}{\tau(\delta)}\right)^{1/(p_d-1)}
 =\left(\frac{K_\delta}{\tau(\delta)}\right)^{d-1}
 =\delta^{-o(1)},
 \qquad H:=\{M\le L_\delta\},
\end{equation}
the maximal estimate and size-biased Markov inequality give
\begin{equation}\label{eq:maximal-pruning-retained-mass}
 \int_{H^c}M
 \le L_\delta^{1-p_d}\int M^{p_d}
 \le\tau(\delta)\int M.
\end{equation}

Thus this single point mask retains \(1-\tau(\delta)\) of the current terminal
load; \eqref{eq:weighted-ordinary-mass-ledger} places its loss inside the same
ordinary physical source-deletion budget.  On \(H\),
\[
 \mathsf N_{I,k}=\sum_TD_T
 \le A_\delta L_\delta B_I^{\rm pred}
 =\delta^{-o(1)}B_I^{\rm pred}.
\]

Every later mask \(0\le a\le{\bf1}_H\) only decreases this positive load, so
\begin{equation}\label{eq:maximal-implies-root-gate}
 \boxed{
 \mathsf X_{I,k}^{\rm root}[a]
 \le\int\mathsf N_{I,k}[a]^2
 \le\delta^{-o(1)}B_I^{\rm pred}
       \int\mathsf N_{I,k}[a].}
\end{equation}

Consequently the Kakeya maximal conjecture implies the hereditary scalar root
gate, \eqref{eq:broad-cycle-budget}, and
\eqref{eq:global-cycle-pruning-target}; no core-by-core deletion is needed.
This implication is scalar.  It gives neither a common-point typed lift nor
the destination-by-destination cuts \eqref{eq:pair-routing-hall-cut}.

The distinction from the recent three-dimensional result is essential.
Wang--Zahl prove the Kakeya \emph{set} conjecture in \(\mathbb R^3\)
\cite{WangZahl3D}; this does not establish the Kakeya maximal function
conjecture, which remains a separate stronger statement
\cite{ZahlKakeyaSurvey}.  Hence their theorem cannot be substituted for
\eqref{eq:hereditary-weighted-kakeya-maximal} in the near-full-mass argument
above.

It does nevertheless suffice for the weaker retention required by the present
source ledger in dimension three.  In the shaded form of the Wang--Zahl
theorem, stated explicitly in \cite{GuthWangZahl3D}, every active
direction-separated family with shading density \(\delta^{o(1)}\) satisfies
\begin{equation}\label{eq:three-dimensional-kakeya-set-volume}
 \left|U(\mathbb T,Y)\right|
 \ge S_\delta\,|\mathbb T|\,|T|,
 \qquad S_\delta=\delta^{o(1)}.
\end{equation}
Here \(U(\mathbb T,Y):=\bigcup_{T\in\mathbb T}Y(T)\).
Work on one of the existing dyadic load cells, so that
\(f_T\simeq a{\bf1}_{Y(T)}\), and put
\[
 M:=\sum_Tf_T,
 \qquad I:=\int M,
 \qquad U:=U(\mathbb T,Y),
 \qquad \overline M:=I/|U|.
\]

Since \(I\le a|\mathbb T||T|\),
\eqref{eq:three-dimensional-kakeya-set-volume} gives
\(\overline M\le aS_\delta^{-1}=\delta^{-o(1)}\).  Lebesgue Markov inequality
and the lower bound \(M\gtrsim a\) on \(U\) show that the single mask
\begin{equation}\label{eq:wang-zahl-low-multiplicity-mask}
 H_{\rm WZ}:=\{x\in U:M(x)\le2\overline M\}
\end{equation}
satisfies
\begin{equation}\label{eq:wang-zahl-retained-load}
 |H_{\rm WZ}|\ge\tfrac12|U|,
 \qquad
 \int_{H_{\rm WZ}}M
 \gtrsim S_\delta I,
 \qquad
 M{\bf1}_{H_{\rm WZ}}\le2aS_\delta^{-1}=\delta^{-o(1)}.
\end{equation}

Thus the mask retains the source-subpower fraction required by the ledger,
and the normalization by \(A_\delta B_I^{\rm pred}\) used above gives
\eqref{eq:maximal-implies-root-gate} for this retained source and all its later
restrictions.  Consequently the Wang--Zahl set theorem closes the scalar
physical gate in \(d=3\), despite not proving the maximal conjecture.  This is
an import of the full three-dimensional Kakeya set theorem, not a derivation
from planar C\'ordoba; it neither closes dimensions \(d\ge4\) nor supplies the
typed Hall routing.

At scalar subpower resolution, the exact content can therefore be compressed
to one refinement statement.  On every regularized weighted tube cell there
must be a single point mask \(H\) such that
\begin{equation}\label{eq:scalar-refinement-kakeya-interface}
 \boxed{
  \int_HM\ge s_\delta\int M,
  \qquad
  \|M{\bf1}_H\|_\infty\le K_\delta,
  \qquad
  s_\delta^{-1},K_\delta=\delta^{-o(1)}.}
\end{equation}
For a dyadic cell this is equivalent, up to subpower factors, to the shaded
Kakeya union-volume statement.  Indeed,
\eqref{eq:scalar-refinement-kakeya-interface} implies
\[
 |U|\ge|\supp(M{\bf1}_H)|
 \ge\frac{\int_HM}{\|M{\bf1}_H\|_\infty}
 \ge\delta^{o(1)}\int M.
\]
Conversely, a union lower bound and the average-multiplicity selection
\eqref{eq:wang-zahl-low-multiplicity-mask}--%
\eqref{eq:wang-zahl-retained-load} give
\eqref{eq:scalar-refinement-kakeya-interface}.  Normalizing \(D_T\) by
\(A_\delta B_I^{\rm pred}\) then gives
\eqref{eq:maximal-implies-root-gate}.  Thus the scalar pruning needed in this
paper has the strength of the Kakeya \emph{set} conjecture, not the stronger
maximal conjecture; the latter improves the retained fraction from subpower to
\(1-o(1)\).

This also prevents an automatic dimension lift.  The proper-face/narrow branch
does descend from \(d\) to \(d-1\), but the broad exact-support source remains
in dimension \(d\).  Applying the three-dimensional theorem at the bottom
therefore leaves, already in \(d=4\), precisely
\eqref{eq:scalar-refinement-kakeya-interface} for the four-dimensional broad
source.  Proving that step is the four-dimensional Kakeya set problem in this
language; it is not supplied by the three-dimensional base or by formal
support descent.  Even the narrow branch requires the stronger hereditary
fibre-aware descendant
\eqref{eq:affine-fibre-three-dimensional-refinement}; an unweighted projected
multi-translation predicate is false by
\eqref{eq:false-unweighted-projected-target}.  Thus any upward recursion must
carry the affine fibre coordinate and its confluence history rather than use a
bare lower-dimensional union estimate.

Nor can the all-support faces be glued by projection alone.  If \(\pi_j\)
deletes the \(j\)-th coordinate, Loomis--Whitney has the direction
\[
 |E|^{d-1}\le\prod_{j=1}^d|\pi_jE|;
\]
lower bounds for all face projections give no lower bound for \(|E|\).  For
example, inside a fixed box the slab
\[
 E_\delta:=\{x:|x_1+\cdots+x_d|\le\delta\}
\]
has \(|E_\delta|\simeq\delta\), while every coordinate hyperplane projection
has volume bounded below by a dimensional constant.  This example only tests
formal face gluing, not a broad tube theorem, but it identifies the missing
datum: a reverse passage needs a uniform bound on the load in the common
projection fibres.  In the present construction that datum is exactly the
physical fibre cutoff/root-cross capacity, not an all-support motivic
identity.  Exploiting broad transversality to obtain such a reverse estimate
would require an additional multilinear Kakeya or Brascamp--Lieb type
inequality; it is not a consequence of the lower-dimensional set theorem.

\subsection{Pointwise broadness and forest congestion}

There is a second obstruction before such a multilinear estimate can even be
applied.  Equation \eqref{eq:unrooted-cube-cauchy-binet} is an average over the
incidence tuple law in a cube \(Q\); the root collision measure is supported on
a common physical point \(x\).  Define the pointwise direction Gram operator,
when the denominator is nonzero, by
\[
 G_{Q,x}:=
 \frac{\sum_T{\bf1}_{Y_Q(T)}(x)e_Te_T^{\mathsf T}}
      {\sum_T{\bf1}_{Y_Q(T)}(x)}.
\]
A broad averaged operator \(G_Q\) need not give positive mass to points where
\(G_{Q,x}\) is broad.  Partition a cube into \(d\) disjoint measurable regions
and shade only one of \(d\) transverse direction classes on each region.  The
integrated direction law can have \(\det G_Q\simeq1\), while
\(\det G_{Q,x}=0\) almost everywhere.  Thus the retained packet event
\eqref{eq:cube-broad-packet-event} does not by itself imply the pointwise
broad inequality required by multilinear Kakeya.

The missing transfer would have to produce one retained point set \(E_Q\) with
\begin{equation}\label{eq:cube-to-point-broad-transfer}
 \int_{E_Q}\sum_T{\bf1}_{Y_Q(T)}
 \ge\delta^{o(1)}\int_Q\sum_T{\bf1}_{Y_Q(T)},
 \qquad
 \det G_{Q,x}\ge\delta^{o(1)}\quad(x\in E_Q),
\end{equation}
or else place the complementary pointwise-narrow mass in one coherent carrier
field.  The displayed disjoint-shading example shows that the first
alternative alone is false.  Choosing a different normal on each narrow point
returns to the forbidden local-to-global assembly; proving a coherent
selection with source-subpower loss is another form of the physical packing
problem.  Hence endpoint multilinear Kakeya can estimate a genuinely
pointwise-broad descendant, but the existing cube-average motivic mark does
not construct that descendant.

A sharper scalar condition measures weighted reuse of a forest rather than
the number of chords.  Let \(F_x^\circ\) be any spanning forest of
\(G_x^{\rm br}\), put
\[
 \mathcal C_x(F^\circ):=E(G_x^{\rm br})\setminus E(F_x^\circ),
\]
write \(M_c:=M_{bb'}\) for \(c=\{b,b'\}\), and for
\(c\in\mathcal C_x(F^\circ)\) let \(P_c\subseteq F_x^\circ\) be its unique
fundamental path.  Split the chord weight among its path edges by numbers
\[
 \alpha_{c,f}(x)\ge0,
 \qquad
 \alpha_{c,f}=0\ \text{if }f\notin P_c,
 \qquad
 \sum_{f\in P_c}\alpha_{c,f}=1.
\]

Define the weighted fundamental-cycle congestion, with value zero when there
are no chords, by
\begin{equation}\label{eq:weighted-fundamental-cycle-congestion}
 \boxed{
 \kappa_{F^\circ}(x):=
 \inf_\alpha\max_{f\in F_x^\circ}
 \frac{\displaystyle
       \sum_{\substack{c\in\mathcal C_x(F^\circ)\\f\in P_c}}
       \alpha_{c,f}(x)M_c(x)}
      {M_f(x)}.}
\end{equation}
Every forest edge has positive weight by the definition of
\(G_x^{\rm br}\).  Finite weighted Hall duality gives the exact cut formula
\begin{equation}\label{eq:fundamental-cycle-congestion-cuts}
 \boxed{
 \kappa_{F^\circ}(x)
 =
 \max_{\varnothing\ne A\subseteq\mathcal C_x(F^\circ)}
 \frac{\displaystyle\sum_{c\in A}M_c(x)}
      {\displaystyle\sum_{f\in\cup_{c\in A}P_c}M_f(x)}.}
\end{equation}

Indeed, the chord edges are demands, the forest edges have capacities
\(K M_f(x)\), and a chord is adjacent precisely to the edges of its
fundamental path.  The cut maximum has a connected-core form.  If
\(\mathscr H(F_x^\circ)\) denotes the nonempty connected subtrees of
\(F_x^\circ\), then
\begin{equation}\label{eq:fundamental-cycle-connected-core}
 \boxed{
 \kappa_{F^\circ}(x)
 =
 \max_{H\in\mathscr H(F_x^\circ)}
 \frac{\displaystyle
       \sum_{\substack{c\in\mathcal C_x(F^\circ)\\P_c\subseteq H}}M_c(x)}
      {\displaystyle\sum_{f\in H}M_f(x)}.}
\end{equation}

To see this, start with a chord cut \(A\), replace it by all chords whose
paths lie in \(\bigcup_{c\in A}P_c\), and split that path union into connected
components.  Its ratio is a weighted average of the component ratios, so one
component is at least as large.  The converse follows by taking all chords
whose paths lie in a fixed \(H\).

Consequently a failure at level \(K\) has a measurable connected witness
\(H_x\) satisfying
\begin{equation}\label{eq:dense-cycle-core-witness}
 \sum_{\substack{c\in\mathcal C_x(F^\circ)\\P_c\subseteq H_x}}M_c(x)
 >
 K\sum_{f\in H_x}M_f(x).
\end{equation}
Choose the lexicographically first witness among those with the fewest forest
edges.  Every edge of this minimal \(H_x\) lies on a retained chord path:
otherwise delete that edge, split into components, and retain a component
whose ratio is still larger than \(K\), contradicting minimality.  Hence every
forest edge belongs to the fundamental cycle of some retained chord, and the
graph formed by \(H_x\) and those chords has no bridge.  We call this the
weighted physical cycle core.  This extraction is finite and measurable; it
creates neither a new descent level nor an estimate.

Finally,
\[
 \sum_{c\in\mathcal C_x(F^\circ)}M_c
 =\sum_{f\in F_x^\circ}
   \sum_c\alpha_{c,f}M_c
 \le\kappa_{F^\circ}(x)\sum_{f\in F_x^\circ}M_f.
\]
Minimizing over the finite set of spanning forests gives a measurable
quantity
\[
 \kappa_{\rm cyc}(x):=\min_{F_x^\circ}\kappa_{F^\circ}(x).
\]
Consequently the hereditary condition
\begin{equation}\label{eq:fundamental-cycle-congestion-bound}
 \boxed{
 \operatorname*{ess\,sup}_x\kappa_{\rm cyc}(x)
 \le\delta^{-o(1)}}
\end{equation}
implies \eqref{eq:broad-cycle-budget}: the minimum in
\eqref{eq:broad-weighted-cycle-excess} is no larger than the chord mass for a
forest minimizing \(\kappa_{F^\circ}\), while that forest's mass is at most
\(B_I^{\rm pred}\sum_bU_b(x)\).  Unlike
\eqref{eq:broad-cycle-rank-condition}, this condition permits many chords
when their demands can be spread across sufficiently large path weights.  It
is a scalar graph criterion only; it does not construct a faithfully labelled
motivic or physical predecessor lift.

For the full-motive route the forest must satisfy an additional compatibility
which is invisible in the scalar graph.  Call a spanning forest \(F_x\)
\emph{motivic-admissible} if, for every physical chord atom \(z\) with child
edge \(c\), the existing operation-history pair complex contains a two-chain
\(\Sigma_{z,F}^{\rm mot}\) whose structural boundary satisfies
\begin{equation}\label{eq:motivic-admissible-forest}
 \boxed{
 \pi_{\rm alg}\partial_{\rm mot}\Sigma_{z,F}^{\rm mot}
 =
 [c]-\sum_{f\in P_F(c)}\epsilon_{c,f}[f],}
\end{equation}
where \(\pi_{\rm alg}\) forgets the operation-history labels down to the
root-child graph, and with every remaining uncompressed boundary face lying
in a strict typed predecessor.  Let \(\mathfrak F_x^{\rm alg}\) be the finite
set of these structural candidates.  Physical common-point realization is
not included in this definition; it is tested later by
\eqref{eq:coherent-physical-capacity-system}.  Thus structural admissibility
does not presuppose the physical conclusion.

The distinction is essential.  The full pointed Hall/root--Bass model in
\cite{UnifiedTower} retains operation histories at the depth required by the
fixed support, while a fixed-depth channel or face-minimal derived
compression need not retain the fundamental-cycle chain
\eqref{eq:motivic-admissible-forest}.  The prime-power all-support comparison
preserves the structural cells once present, but it does not turn an arbitrary
physical graph forest into an admissible one.  Thus define
\[
 \kappa_{\rm mot}(x):=
 \min_{F_x\in\mathfrak F_x^{\rm alg}}\kappa_F(x),
\]
with value \(+\infty\) when \(\mathfrak F_x^{\rm alg}=\varnothing\).
The unrestricted \(\kappa_{\rm cyc}\) remains a valid purely scalar
criterion; \(\kappa_{\rm mot}\) is the corresponding criterion available to
the motivic fundamental-cycle route.  In particular,
\begin{equation}\label{eq:motivic-cycle-congestion-bound}
 \boxed{
 \operatorname*{ess\,sup}_x\kappa_{\rm mot}(x)
 \le\delta^{-o(1)}}
\end{equation}
closes the scalar gate by the same forest-allocation proof as
\eqref{eq:fundamental-cycle-congestion-bound}.  If this condition fails but
\(\mathfrak F_x^{\rm alg}\ne\varnothing\), the dense connected core of an
admissible forest is the source passed to the hybrid comparison below.

\subsection{Typed chord routing}

There is nevertheless a genuine reduction in the typed comparison that must
be proved.  Write \(b(e)\) for the unique root child whose descendant set
contains \(\gamma(e)\), and fix an order on the root children.  Let
\(F_x^{\min}\) be the measurable minimizing forest chosen in
\eqref{eq:broad-weighted-cycle-excess}, and let \(E_x^{\rm ch,min}\) be the
once-oriented physical broad pairs \((e,e')\) whose child edge
\(\{b(e),b(e')\}\) is a chord of \(F_x^{\min}\).  For an admissible strict
type \(\xi\), write \(\mathcal P_{\rm mot}(e,e';\xi)\) for the finite set of
face/residue-support provenances in the existing pair complex which realize
\(\xi\) from \((e,e')\).  Let
\[
 \mathcal P_{\xi,x}^{\rm ch}
 :=\bigcup_{(e,e')\in E_x^{\rm ch,min}}
      \mathcal P_{\rm mot}(e,e';\xi),
 \qquad
 \mathscr Y_x^{\rm ch,lab}
 :=\bigsqcup_\xi\bigsqcup_{\rho\in\mathcal P_{\xi,x}^{\rm ch}}
      \mathscr Y_{\xi,x}^{[2]}.
\]
It is enough to construct a measurable joint lift on this chord source alone,
\[
 \Theta_x^{\rm ch}:E_x^{\rm ch,min}\longrightarrow
 \mathscr Y_x^{\rm ch,lab},
\]
which preserves \(x\), uses an admissible strict type
\(\xi\in\mathcal S_{\rm alg}(e,e')\), retains \(\rho\) until the branch trace,
and satisfies \(\rho\in\mathcal P_{\rm mot}(e,e';\xi)\) on each source pair.
Let
\(\operatorname{pr}_{\rm phys}(\xi,\rho;h,h'):=(\xi,h,h')\).

The required coalesced weighted-fibre estimate is
\begin{equation}\label{eq:chord-joint-typed-fibre-capacity}
 \boxed{
 \sum_{\substack{(e,e')\in E_x^{\rm ch,min}:\\
       \operatorname{pr}_{\rm phys}\Theta_x^{\rm ch}(e,e')
       =(\xi,h,h')}}
 u_e(x)u_{e'}(x)
 \le C_\delta v_{\xi,h}(x)v_{\xi,h'}(x),
 \qquad C_\delta=\delta^{-o(1)}.}
\end{equation}
The projection in the sum forgets \(\rho\) only after all of its labelled
contributions to the same physical destination have been coalesced.
Summing \eqref{eq:chord-joint-typed-fibre-capacity}, then using predecessor R6
and \eqref{eq:coalesced-predecessor-unary-ledger}, gives
\[
 \begin{aligned}
  \mathfrak C_{I,k}^{\rm br}
  &=2\int\sum_{(e,e')\in E_x^{\rm ch,min}}u_eu_{e'}\,dx\\
  &\le 2C_\delta\sum_\xi\int\mathsf N_\xi^2
   \lesssim
   \delta^{-o(1)}B_I^{\rm pred}
   \int\mathsf N_{I,k}.
 \end{aligned}
\]
Hence \eqref{eq:chord-joint-typed-fibre-capacity} proves
\eqref{eq:broad-cycle-budget} and closes the scalar root gate.  No routing is
needed on the forest edges for this implication.  If a deterministic chord
lift is too rigid, the same proof uses the measurable Hall cuts restricted to
the chord-source measure.

The chord-to-predecessor target can be weakened once more by using the whole
fundamental cycle.  For this typed route choose measurably
\(F_x^{\rm mot}\in\mathfrak F_x^{\rm alg}\); if the admissible class is empty,
this route has a structural obstruction before any estimate is attempted.
With the fixed child order, set
\[
 \lambda_{I,k,x}^{\rm br}
 :=
 \sum_{\substack{b(e)<b(e'),\ T(e)\ne T(e')\\
                  e,e'\ {\rm broad\ at}\ x}}
 u_e(x)u_{e'}(x)\,\delta_{(e,e')}.
\]
Let
\[
 E_x^{\rm ch,mot}
 :=\{(e,e')\in E_x^{\rm br}:
       \{b(e),b(e')\}\notin F_x^{\rm mot}\}
\]
be the chord source belonging to this selected admissible forest.  It need not
equal \(E_x^{\rm ch,min}\), which belongs to the unrestricted scalar
minimizer.  Disintegrate this once-oriented root measure into
\[
 \lambda_{I,k,x}^{\rm br}
 =\lambda_{I,k,x}^{F}+\lambda_{I,k,x}^{\rm ch},
\]
according as the child edge lies in \(F_x^{\rm mot}\) or its complement, and
put
\[
 d\Lambda_{I,k}^{F}:=2\,d\lambda_{I,k,x}^{F}\,dx,
 \qquad
 d\Lambda_{I,k}^{\rm ch}:=2\,d\lambda_{I,k,x}^{\rm ch}\,dx.
\]
Then
\begin{equation}\label{eq:forest-reservoir-capacity}
 \boxed{
 \|\Lambda_{I,k}^{F}\|
 \le 2B_I^{\rm pred}\int\mathsf N_{I,k},
 \qquad
 \|\Lambda_{I,k}^{\rm ch}\|
 =:\mathfrak C_{I,k}^{\rm br}(F^{\rm mot})
 \ge\mathfrak C_{I,k}^{\rm br}.}
\end{equation}

For a chord \(\{b,b'\}\), let \(P_{F_x^{\rm mot}}(b,b')\) be its unique
forest path.
For a physical chord atom \(z\), declare \(z\rightsquigarrow_F y\) precisely
when \(y\) occurs as a boundary face of a trace-compatible faithfully labelled
motivic filling of its fundamental cycle, and \(y\) is either a compatible
physical pair over an edge of \(P_{F_x^{\rm mot}}(b,b')\) or a strict typed
predecessor.
Thus
\[
 z\rightsquigarrow_F y,
 \qquad
 y\in
 \mathscr Z_{I,k}^{F}
 \sqcup\bigsqcup_\xi\mathscr Y_\xi^{[2]},
\]
where \(\mathscr Z_{I,k}^{F}:=\supp\Lambda_{I,k}^{F}\), and compatibility
includes the common physical point, source histories, orientation line and
residue-support label \(\rho\).  If no such filling exists, the corresponding
allowed neighbourhood is empty.  For a measurable chord set \(E\), write
\(\mathcal A_F(E)\) and \(\mathcal A_\xi^{\rm ch}(E)\) for its allowed
neighbourhoods in the forest and predecessor targets.  The exact hybrid
capacity target is
\begin{equation}\label{eq:chord-hybrid-hall-capacity}
 \boxed{\begin{aligned}
 \Lambda_{I,k}^{\rm ch}(E)
 &\le C_\delta\left[
       \Lambda_{I,k}^{F}(\mathcal A_F(E))
       +\sum_\xi
        \Lambda_\xi^{\rm pred}(\mathcal A_\xi^{\rm ch}(E))
      \right]\\
 &\hspace{7em}\text{for every measurable }E,
 \qquad C_\delta=\delta^{-o(1)}.
 \end{aligned}}
\end{equation}

As in \eqref{eq:pair-routing-hall-cut}, fibrewise max-flow/min-cut shows that
this is equivalent to one measurable positive splitting kernel into the
combined target, chosen once and then restricted hereditarily.  Taking the
whole chord source in \eqref{eq:chord-hybrid-hall-capacity}, using
\eqref{eq:forest-reservoir-capacity}, predecessor R6 and
\eqref{eq:coalesced-predecessor-unary-ledger}, yields
\[
 \mathfrak C_{I,k}^{\rm br}(F^{\rm mot})
 \lesssim
 \delta^{-o(1)}B_I^{\rm pred}\int\mathsf N_{I,k}.
\]

Thus the hybrid Hall estimate also proves
\eqref{eq:broad-cycle-budget} and closes the root gate for the whole
\(F^{\rm mot}\)-decomposition.  It is weaker than
\eqref{eq:chord-joint-typed-fibre-capacity}: forest boundary faces may absorb
chord mass, but only up to their already bounded global capacity.  A formal
fundamental-cycle identity only names candidate faces.  Constructing a
trace-compatible physical filling is part of the realization interface, and
even such a filling proves neither the Hall cuts nor a bound for their
congestion.

\subsection{Faithful positive cycle realization}

Structural cycle closure and positive physical control must be kept separate.
A coordinate lift may construct cycle classes and their Gersten boundaries,
while triangular branch curvature survives compression; closure does not bound
the faithfully labelled coefficient current.  In the present dictionary,
\(\mathfrak R_{\rm alg}\) is the structural cycle lift, whereas
\(\mathfrak C_{I,k}^{\rm br}\) measures the positive physical mass left on its
independent chords.  The full motivic filtration names these typed cycles, but
algebraic \(K\)-theory proves neither
\eqref{eq:broad-cycle-budget} nor Hall capacity.  Its legitimate role is to
organize the cycle-space obstruction whose positive physical mass must still
be estimated.

\subsection{A labelled chain model}

More precisely, importing that interface here would require no new descent
coordinate.  The once-oriented measure introduced above satisfies
\[
 d\Lambda^{\rm root}_{I,k}
 =2\,d\lambda_{I,k,x}^{\rm br}\,dx,
 \qquad
 J_x^{\rm br}
 :=(b(e),b(e'))_\ast\lambda_{I,k,x}^{\rm br}
 =\sum_{b<b'}M_{bb'}(x)[b,b'].
\]

Let \(\widetilde C_\bullet^{\rm mot}\) be the subcomplex of the existing pair
base change \(\mathcal K^{[2]}\) generated by the typed root-collision cells
and by the fundamental-cycle two-cells belonging to the selected
motivic-admissible forest.
The provenance set \(\mathcal P_{\rm mot}(e,e';\xi)\) just introduced is a
label inside the all-support filtration, not another filtration coordinate.
Its physical root-collision current is the weighted degree-one chain
\[
 \mathbf c_{I,k,x}^{\rm root}
 :=\sum_{\substack{b(e)<b(e'),\ T(e)\ne T(e')\\
                   e,e'\ {\rm broad\ at}\ x}}
 u_e(x)u_{e'}(x)[e,e']_{\rm mot}.
\]

For \(A\) in the finite physical pair fibre, write
\(\mathbf c_{I,k,x}^{\rm root}|_A\) and
\(\lambda_{I,k,x}^{\rm br}|_A\) for the corresponding restrictions.  These
restrictions are taken inside one fixed realization; no new routing is chosen
after \(A\) is specified.

Let \(C_\bullet^{\rm str,lab}(x)\) be the structural labelled chain complex
whose one-generators are admissible quadruples
\[
 (e,e';\xi,\rho),
 \qquad
 \xi\in\mathcal S_{\rm alg}(e,e'),
 \qquad
 \rho\in\mathcal P_{\rm mot}(e,e';\xi).
\]

Let \(C_\bullet^{\rm lab}(x)\) be its physical refinement whose one-generators
are admissible sextuples
\[
 (e,e';\xi,\rho;h,h'),
 \qquad
 \xi\in\mathcal S_{\rm alg}(e,e'),
 \qquad
 \rho\in\mathcal P_{\rm mot}(e,e';\xi),
\]
where \((h,h';x)\in\mathscr Y_{\xi,x}^{[2]}\) is an allowed common-point
physical lift of \((e,e';x)\).  Thus parallel strict-type labels,
residue-support provenances, source histories and destination histories are
not identified.  Write
\[
 \tau_{\rm phys}:C_\bullet^{\rm lab}(x)
 \longrightarrow C_\bullet^{\rm str,lab}(x)
\]
for the map forgetting only \((h,h')\).  Write \(\pi\) for the map forgetting
all labels except the
two root children, and let \(\operatorname{Tr}_{\rm br}\) forget
\((\xi,\rho;h,h')\) but retain the physical source pair; thus
\(\pi=(b,b')_\ast\circ\operatorname{Tr}_{\rm br}\).


There are then exactly three structural requirements.  First, the motivic
pair-orientation line must lift
to the physical branch-orientation line:
\[
 \jmath_{e,e'}^{\rm mot}:
 \mathfrak o_{\rm mot}(e,e')
 \xrightarrow{\ \sim\ }
 \mathfrak o_{\rm br}(e,e'),
 \qquad
 \jmath_{e,e'}^{\rm mot}\ \text{is branch-exchange equivariant}.
\]

The faithfully labelled structural carrier can be constructed explicitly.
Let \(\mathscr L_x^{\rm str}\) be the finite set of structural labelled cells
in the selected Hall/root--Bass subcomplex and put
\[
 F_x^{\rm flag}
 :=\mathbf Q(z_\sigma:\sigma\in\mathscr L_x^{\rm str}).
\]

For every strict labelled flag
\(\boldsymbol\sigma=(\sigma_0<\cdots<\sigma_j)\), set
\[
 \kappa_x^{\rm flag}(\boldsymbol\sigma)
 :=(-1)^j\{z_{\sigma_0},\ldots,z_{\sigma_j}\}
 \in K_{j+1}^M(F_x^{\rm flag}).
\]
Writing \(x_\sigma^0=[z_\sigma=0]\), its zero-face Gersten boundary is
\begin{equation}\label{eq:coordinate-flag-gersten-boundary}
 \boxed{
 \partial_0\kappa_x^{\rm flag}(\boldsymbol\sigma)
 =
 \sum_{i=0}^j(-1)^i x_{\sigma_i}^0
 \{z_{\sigma_0},\ldots,\widehat{z_{\sigma_i}},\ldots,z_{\sigma_j}\}.}
\end{equation}

Send each residue generator on the right to the labelled flag with
\(\sigma_i\) deleted.  Then
\begin{equation}\label{eq:coordinate-flag-chain-realization}
 \boxed{
 \mathcal R_x^{\rm flag}
  \bigl(\partial_0\kappa_x^{\rm flag}(\boldsymbol\sigma)\bigr)
 =
 \partial_{\rm str,lab}[\boldsymbol\sigma].}
\end{equation}

The independent coordinates retain the deleted-cell label, so this
realization is faithful before compression.  In degree two it uses coordinate
\(K_3\) symbols to realize the fundamental-cycle fillings; coordinate \(K_4\)
symbols provide the coherence between alternative fillings.  Consequently
the structural labelled obstruction vanishes for every
\(F\in\mathfrak F_x^{\rm alg}\).  This is the legitimate algebraic
\(K\)-theory contribution.  It does not produce a lift through
\(\tau_{\rm phys}\), positivity, or a bound on physical coefficients.

After this orientation identification, one must construct, measurably in
\(x\), measure-valued linear maps
\[
 \nu_j^{\rm cyc}(x):
 \widetilde C_j^{\rm mot}\longrightarrow
 \mathcal M(C_j^{\rm lab}(x))
\]
where \(\mathcal M\) denotes finite signed measures on the labelled
generators, and satisfying
\begin{equation}\label{eq:motivic-physical-cycle-interface}
 \boxed{\begin{gathered}
 \partial_{\rm lab}\circ\nu_2^{\rm cyc}
   =\nu_1^{\rm cyc}\circ\partial_{\rm mot},\\
 \nu_1^{\rm cyc}
   (\mathbf c_{I,k,x}^{\rm root}|_A)\ge0,\\
 \operatorname{Tr}_{\rm br}
   \nu_1^{\rm cyc}(\mathbf c_{I,k,x}^{\rm root}|_A)
   =\lambda_{I,k,x}^{\rm br}|_A
 \quad\text{for every }A.
 \end{gathered}}
\end{equation}

The compressed identity
\[
 \pi_\ast\nu_1^{\rm cyc}
   (\mathbf c_{I,k,x}^{\rm root}|_A)
 =(b,b')_\ast(\lambda_{I,k,x}^{\rm br}|_A)
\]
then follows rather than being imposed separately.

The orientation lift, the first identity, and the last trace identity are,
respectively, the double-orientation, zero matching-defect, and
branch-mass-compatible trace conditions.  A coordinate flag/\(K\)-theory lift
\eqref{eq:coordinate-flag-gersten-boundary}--%
\eqref{eq:coordinate-flag-chain-realization} constructs their structural
labelled part.  Retaining \(\rho\), positivity of the realized degree-one
current, and the branch-trace identity prevent compression from
discarding a cycle detected by the physical branch trace.  In particular, the
map is not allowed to coalesce two equal \(\xi\)'s before their residue-support
labels have been traced.  These properties are not consequences of a signed
Gersten lift.  Even after
\eqref{eq:motivic-physical-cycle-interface} is proved, the quantitative
estimate \eqref{eq:broad-cycle-budget} remains independent.  Thus the
cycle-space version of the full-motive route has two separate obligations:
a faithful positive cycle realization and a weighted physical cycle budget,
not merely formal vanishing of a motivic boundary.  For the scalar root gate
only the latter is needed; for the hereditary typed statement the exact
capacity requirement remains \eqref{eq:pair-routing-hall-cut}.

For a structural candidate satisfying
\eqref{eq:motivic-admissible-forest}, the full Hall/root--Bass two-chain and the
coordinate flag carrier already supply the faithfully provenance-labelled
structural boundary.  The remaining question is whether that boundary has a
common-point positive physical lift with the prescribed branch trace and
destination capacity.  No separate obstruction layer is needed: all of these
requirements enter the single finite system below.  For a structural one-generator
\(g=(e,e';\xi,\rho)\), define its common-point physical fibre
\[
 \mathscr L_x(g)
 :=
 \{(h,h'):(e,e';\xi,\rho;h,h')
             \in C_1^{\rm lab}(x)\},
\]
and put
\begin{equation}\label{eq:physical-pair-support}
 \boxed{
 \mathcal S_{\rm phys}(e,e';x)
 :=
 \{(\xi,\rho):
   \xi\in\mathcal S_{\rm alg}(e,e'),\
   \rho\in\mathcal P_{\rm mot}(e,e';\xi),\
   \mathscr L_x(e,e';\xi,\rho)\ne\varnothing\}.}
\end{equation}
Because the fibres are finite, nonemptiness of
\(\mathcal S_{\rm phys}(e,e';x)\) is exactly the measurable atomic support
test.  Existing insertion/deletion/confluence spans prove this test for the
labels they realize, but support alone gives no pair RN bound.

\subsection{External-product and correlated capacities}

Let
\(\mathcal S_{\boxtimes}(e,e';x)\subseteq
\mathcal S_{\rm phys}(e,e';x)\) consist of labels realized by an endpointwise
external product of two unary physical kernels, and put
\[
 E_x^{\boxtimes}
 :=\{(e,e')\in E_x^{\rm br}:
       \mathcal S_{\boxtimes}(e,e';x)\ne\varnothing\},
 \qquad
 E_x^{\rm corr}:=E_x^{\rm br}\setminus E_x^{\boxtimes}.
\]

Fix one external-product label
\(\theta=(\xi,\rho,K_\theta^{(1)},K_\theta^{(2)})\), and let
\(E_x(\theta)\subseteq E_x^{\boxtimes}\) be the pairs assigned to it by the
fixed label order.  Unary R5 gives
\[
 \sum_e u_e(x)K_\theta^{(1)}(h\mid e)
 \le A_{\theta,1}v_{\xi,h}(x),
 \qquad
 \sum_{e'}u_{e'}(x)K_\theta^{(2)}(h'\mid e')
 \le A_{\theta,2}v_{\xi,h'}(x).
\]
Nonnegativity and the external-product formula then give the pair estimate
\begin{equation}\label{eq:external-span-pair-capacity}
 \boxed{
 \begin{aligned}
 &\sum_{(e,e')\in E_x(\theta)}
  u_eu_{e'}K_\theta^{(1)}(h\mid e)
             K_\theta^{(2)}(h'\mid e')\\
 &\qquad\le
 \left(\sum_eu_eK_\theta^{(1)}(h\mid e)\right)
 \left(\sum_{e'}u_{e'}K_\theta^{(2)}(h'\mid e')\right)
 \le
 A_{\theta,1}A_{\theta,2}
 v_{\xi,h}v_{\xi,h'} .
 \end{aligned}}
\end{equation}
The tensor product of \eqref{eq:physical-kernel-chain-rule} on the two
endpoints supplies the labelled matching identity, and its coefficients are
positive with exact branch trace.  Coalescing the finitely many
\(\theta\)-labels costs only the fixed operation-history alphabet and the
subpower carrier factor already counted in
\eqref{eq:global-coarse-carrier-count}.  Hence
\eqref{eq:coherent-physical-capacity-system} is proved on
\(E_x^{\boxtimes}\).  The genuinely unresolved typed source is
\(E_x^{\rm corr}\).  The diagonal confluence example above shows why the same
calculation cannot be extended to a correlated joint kernel.

\subsection{Mode decomposition and fractional capacity}

A concrete sufficient route on \(E_x^{\rm corr}\) is to split it among
coherent deterministic modes which retain \((e,e';\rho)\), common-point
support and the same motivic boundary.  This is genuinely additional to the
merged comparison theorem: the Hall/root--Bass realization in
\cite{UnifiedTower} is faithful on structural operation words but does not
assert local fullness or produce positive physical kernels.

Fix a labelled type \(\theta=(\xi,\rho)\) and a coherent deterministic output
\[
 F_\theta(e,e')=(h_\theta(e,e'),h'_\theta(e,e'))
\]
on every positive source pair.  Its mode-conflict graph
\(\mathfrak G_{\theta,x}^{\rm mode}\) has the source pairs as vertices; two
vertices are adjacent when they share the first endpoint but have different
first outputs, or share the second endpoint but have different second outputs.

A color class is exactly a source set on which the output separates as
\((\phi(e),\psi(e'))\).  Indeed, a convex combination equal to a point mass
can use only product factors supported at that same point, so the active
source of each factor is independent; conversely each color class defines
the two endpoint maps.  Consequently
\begin{equation}\label{eq:deterministic-mode-conflict-rank}
 \boxed{
 \mathfrak r_{+}^{\rm mode}(F_\theta;x)
 =\chi(\mathfrak G_{\theta,x}^{\rm mode}).}
\end{equation}

For the diagonal confluence map
\((i,j)\mapsto(i+j,i+j)\) on \((\mathbf Z/R)^2\), every color class is a
matching and the \(R\) level sets \(i+j=s\) are perfect matchings.  Hence
\(\mathfrak r_{+}^{\rm mode}=R\), although there is only one unlabelled
structural operation type.  Structural word count and the \(q^2\) provenance
therefore do not bound the physical mode number.

If several coherent outputs are allowed, minimize the right-hand side over
the source-local choices before coalescing destination capacities.  Removing
\(\beta_1(\mathfrak G_{\theta,x}^{\rm mode})\) edges leaves a forest; choosing
one endpoint of each removed edge gives a feedback vertex set of size at most
\(\beta_1\).  Color the remaining forest with two colors and give those
vertices separate colors.  Hence
\begin{equation}\label{eq:mode-conflict-cycle-rank-bound}
 \boxed{
 \chi(\mathfrak G_{\theta,x}^{\rm mode})
 \le 2+\beta_1(\mathfrak G_{\theta,x}^{\rm mode}).}
\end{equation}
Thus a subpower cycle-rank bound on the faithfully labelled mode-conflict
graph is a valid sufficient condition for the positive mode gate.  It is
not the cycle rank of the compressed root-child interaction graph
\(G_x^{\rm br}\), and the weighted destination capacities must still be
coalesced as in \eqref{eq:coherent-physical-capacity-system}.

Positive splitting permits a sharper weighted fractional form.  Let
\(\mathcal I(\mathfrak G_{\theta,x}^{\rm mode})\) be the independent source
sets and write \(a_C:=A_{C,1}A_{C,2}\).  Put
\begin{equation}\label{eq:fractional-mode-capacity}
 \boxed{
 \begin{aligned}
 \kappa_{+}^{\rm mode}(\theta,x)
 &:=
 \min\left\{
  \sum_{C\in\mathcal I}a_Ct_C:
  t_C\ge0,\ \sum_{C\ni z}t_C\ge1\ \forall z
 \right\},\\
 \kappa_{+}^{\rm mode}(\theta,x)
 &\le\delta^{-o(1)}
 \quad\Longrightarrow\quad
 \text{\eqref{eq:coherent-physical-capacity-system} on this selector}.
 \end{aligned}}
\end{equation}

For a feasible family, set
\[
 \alpha_C(z)
 :=
 \frac{t_C{\bf1}_{C}(z)}
      {\sum_{C'\ni z}t_{C'}}.
\]
Then \(\sum_C\alpha_C(z)=1\), \(\alpha_C(z)\le t_C\), and the output separates
endpointwise on \(C\).  Applying
\eqref{eq:external-span-pair-capacity} on every \(C\) proves the implication
in \eqref{eq:fractional-mode-capacity}.
The exact weighted finite LP dual is
\begin{equation}\label{eq:mode-conflict-fractional-dual}
 \boxed{
\kappa_{+}^{\rm mode}(\theta,x)
 =
 \max\left\{
  \sum_z\omega_z:
  \omega_z\ge0,\qquad
  \sum_{z\in C}\omega_z\le a_C
  \text{ for every independent }C
 \right\}.}
\end{equation}
A large dual value therefore extracts one coalesced weighted physical source
on which no coherent separated mode carries more than its unary RN cost.  This is the
correct positive mode-core obstruction; estimating independently selected
local color classes would again be a local-to-global error.  An integral
coloring gives
\(\kappa_+^{\rm mode}\le
\chi(\mathfrak G_{\theta,x}^{\rm mode})\max_Ca_C\), so
\eqref{eq:mode-conflict-cycle-rank-bound} remains a stronger sufficient
shadow.

Within this deterministic-mode route, the remaining sufficient statement is
\[
 \operatorname*{ess\,sup}_x
 \sum_\theta\kappa_{+}^{\rm mode}(\theta,x)
 \le\delta^{-o(1)}.
\]
This is the physical no-new-communication bound for the joint
first-obstruction selector, with provenance and common-point labels retained.
It cannot be replaced by formal termination, operation-word faithfulness, or
the \(q^2\) line.

\subsection{Finite linear programming formulation}

Now let \(\nu_\bullet^{\rm str}\) be a chosen coordinate-flag structural chain
map.  Keep the source atom as part of every index.  Order all available
degree-one physical labels over the positive source atoms, put their
coefficients in \(a_x\), and put the available source-labelled signed
higher-cell coefficients in \(b_x\).  Missing physical support simply removes
the corresponding column and makes the system below infeasible if that
structural coefficient is required.  The
boundary, fixed structural pushforward and source-by-source branch-trace
equations have the form \(M_xa_x+N_xb_x=c_x\).  No source labels are
coalesced, so deleting source blocks implements every later restriction
without a new choice.
Let \(D_xa_x\) be the vector obtained by coalescing all degree-one labelled
coefficients at each physical predecessor destination, and let \(q_x\) be the
corresponding vector of capacities
\(v_{\xi,h}(x)v_{\xi,h'}(x)\); for the hybrid route append the forest
reservoir coordinates.  The single typed realization target is
\begin{equation}\label{eq:coherent-physical-capacity-system}
 \boxed{
 \begin{gathered}
  \text{find }a_x\ge0,\quad b_x\text{ signed, such that}\\
  M_xa_x+N_xb_x=c_x,\qquad
  D_xa_x\le C_\delta q_x,\qquad C_\delta=\delta^{-o(1)} .
 \end{gathered}}
\end{equation}

The first equation simultaneously imposes the fixed structural pushforward,
the labelled chain identity and the source-by-source branch trace.  Hence a
solution of \eqref{eq:coherent-physical-capacity-system} is already a
source-local coherent positive physical lift with the required RN bound.
Conversely every such lift gives these coefficients.

There is also one exact dual test.  Introduce a nonnegative slack vector in the
capacity inequality.  Finite-dimensional Farkas duality says that
\eqref{eq:coherent-physical-capacity-system} is feasible exactly when
\begin{equation}\label{eq:coherent-capacity-dual-cut}
 \boxed{
  \langle y,c_x\rangle
   +C_\delta\langle\lambda,q_x\rangle\ge0
  \quad\text{whenever}\quad
  N_x^{\mathsf T}y=0,\quad
  M_x^{\mathsf T}y+D_x^{\mathsf T}\lambda\ge0,\quad
  \lambda\ge0.}
\end{equation}
Because all source labels remain uncoalesced, the lexicographically first
feasible basic solution is measurable in \(x\), and later restriction deletes
source blocks from that same solution.  If a coherent filling has already
been fixed on every allowed source--destination edge, the variables \(b_x\)
can be eliminated and \eqref{eq:coherent-capacity-dual-cut} reduces to the
ordinary Hall cuts \eqref{eq:pair-routing-hall-cut}.  In general it is the
correct replacement: support \eqref{eq:physical-pair-support} and Hall
capacity are merely projections; graded matching is encoded by the linear rows
and positivity by \(a_x\ge0\).

Taking \(A=E_x^{\rm ch,mot}\) in
\eqref{eq:motivic-physical-cycle-interface} supplies the structural labelled
chord lift required above, possibly as a positive splitting kernel rather than
a deterministic \(\Theta_x^{\rm ch}\); the matching identity for the
two-cells supplies its trace-compatible candidate boundary faces.  What it
does not supply is the hybrid
destination-capacity estimate
\eqref{eq:chord-hybrid-hall-capacity}; the direct chord estimate
\eqref{eq:chord-joint-typed-fibre-capacity} is a stronger sufficient
alternative.  Thus restriction compatibility connects the full motivic
interface to the scalar cycle reduction without turning a collection of
locally reselected maps into a global one.

For the selected \(F_x^{\rm mot}\), if its congestion exceeds a proposed
level \(K\), let \(H_x\subseteq F_x^{\rm mot}\) be the minimal connected
witness obtained from \eqref{eq:dense-cycle-core-witness} with
\(F^\circ=F_x^{\rm mot}\), and restrict the physical source to
\[
 E_x^{\rm core}(H_x)
 :=
 \{(e,e')\in E_x^{\rm ch,mot}:
   P_{F_x^{\rm mot}}(b(e),b(e'))\subseteq H_x\}.
\]
Applying the same source-local realization to this single restriction retains
the whole bridgeless cycle core and disintegrates it by the existing
\((\xi,\rho;h,h')\) labels.  This is the correct input to the full
all-support descent: the first failed typed obstruction is taken on the core
current, not independently on its chord atoms.  No conclusion follows merely
from the existence of those labels.  Different cores can have overlapping
predecessor neighbourhoods, so their capacities must remain coalesced in the
single Hall ledger; summing separately proved core estimates would reintroduce
the forbidden local-to-global step.

In particular, \(\beta_1(G_x^{\rm br})\) is not a Hall theorem.  It counts
excess scalar interactions in the compressed root-child graph.  Cycles in the
faithfully labelled routing multigraph instead describe possible
redistributions among admissible destinations; without target capacities they
give neither a positive routing nor its Radon--Nikodym bound.

Thus
\eqref{eq:root-collision-realization-target}, rather than another fixed-root
C\'ordoba estimate, is the natural full-motive target of the
Selmer--Cartan route.  The scalar gate could of course be proved by a different
physical argument; the pair realization is the form compatible with the
descent filtration.  The broad packet selection has already been normalized in
the source-mass ledger; it does not by itself prove the pair-measure domination
in \eqref{eq:root-collision-realization-target}.

\subsection{The algebraic candidate and the remaining physical gate}

The root-cross quantity does have a canonical algebraic controller.  The
divided-power identity
\[
 \gamma_2(y+z)-\gamma_2(y)-\gamma_2(z)=yz
\]
gives, after the real realization \(\gamma_2(t)=t^2/2\),
\begin{equation}\label{eq:root-pd-polarization}
 \mathsf P^{\rm root}_{I,k}(x)
 :=
 \gamma_2(W_{r_\ast}(x))
 -
 \sum_{b\in\operatorname{ch}(r_\ast)}\gamma_2(W_b(x)),
\end{equation}
and the history partition yields
\begin{equation}\label{eq:root-cross-is-pd}
 \boxed{
 \mathsf P^{\rm root}_{I,k}
 =\sum_{b<b'}W_bW_{b'}\ge0,
 \qquad
 \mathsf X^{\rm root}_{I,k}
 =2\int\mathsf P^{\rm root}_{I,k}.}
\end{equation}
The factor \(1/2\) is used only over the real physical leg; the integral
controller requires no division by \(2\).  The positivity in
\eqref{eq:root-cross-is-pd} appears only after the three PD terms have been
combined.  The term \(\gamma_2(y+z)\) depends jointly on both inputs, while the
two unary terms are subtracted.  Hence the polarization identity is not a
termwise positive decomposition and cannot itself define the kernels
\(\mathcal R_\xi\), prove the joint weighted-fibre estimate
\eqref{eq:joint-typed-fibre-capacity}, or verify the Hall cuts.

This is intrinsically a second-moment object.  Its physical source is the pair
diagonal
\begin{equation}\label{eq:physical-pair-diagonal}
 \Omega^{[2]}_{I,k}
 :=
 \bigsqcup_{\gamma,T,\gamma',T'}
 \{(x;\gamma,T;\gamma',T'):
       u_{\gamma,T}(x)u_{\gamma',T'}(x)>0\},
 \qquad
 d\mu^{[2]}_{I,k}
 :=
 u_{\gamma,T}(x)u_{\gamma',T'}(x)\,dx,
\end{equation}
the fibre product of two marked history sources over the same physical point.
Accordingly the relevant realization is the external pair base change
\begin{equation}\label{eq:pair-filtered-realization}
 \mathcal K^{[2]}:=\mathcal K\boxtimes_{\R^d}\mathcal K
\end{equation}
over the fixed \((\eta,\nu)\)-fibre, with the existing support, confluence and
Cartan filtrations on both factors.  The kernels in
\eqref{eq:root-collision-realization-target} must positively route its broad
root-collision part by the first failed typed obstruction.  This is not another
motivic level: it is the physical realization forced by the quadratic
estimate R6.

Two firewalls are important.  First, both factors have the same operation
support, so support-squarefree local faithfulness does not apply; the
unrestricted PD controller records the repeated input but supplies no positive
measure partition.  Second, the physical product
\(u_\gamma u_{\gamma'}\) is polarization in \(\Gamma^2(D_I)\), whereas
\(L^{\rm conf}_{q^2}=W_q/q^2W_q\) records the supplied primitive filtered
confluence provenance.  The product neither creates that lift nor changes
coefficient depth.  The merged prime-power all-support comparison already
handles the latter, and algebraic \(K\)-theory is not an analytic input here.

For a root-collision pair the PD hull nevertheless gives a finite algebraic
candidate.  Put
\[
 \operatorname{Pol}(\gamma,\gamma')
 :=
 \gamma_2(h_\gamma+h_{\gamma'})
 -\gamma_2(h_\gamma)-\gamma_2(h_{\gamma'}),
\]
apply the support-functorial filtration, and take its first nonzero typed
graded component; if it vanishes, use the finite Reedy reconstruction.  The
result is
\begin{equation}\label{eq:algebraic-root-collision-routing}
 \mathfrak R_{\rm alg}(\gamma,\gamma')
 \in
 \bigoplus_{\xi\prec(\eta;I,\chi,\nu,n)}\mathbf Z[\xi].
\end{equation}
This chain lowers \eqref{eq:descent-complexity}, but it may have signs.
Therefore the routing kernels must be an additional positive measurable
refinement of \(\mathfrak R_{\rm alg}\), with the RN ledger
\eqref{eq:root-collision-realization-target}.  Formal Reedy termination does
not imply that inequality.  Concretely, the missing refinement is a positive
joint physical lift whose weighted fibres satisfy
\eqref{eq:joint-typed-fibre-capacity}, or a splittable lift satisfying
\eqref{eq:pair-routing-hall-cut}.  Signed PD cancellation proves neither.
\section{Where the two-dimensional theorem enters}

For a rooted \(d\)-packet, the full direction support has cardinality \(d\):
one root slot and \(d-1\) side slots.  The unified tower's exact-support comparison is
used for \(|I|\ge3\).  Its last root projection reaches \(|I|=2\), the ordinary
Kummer/Heisenberg low sector, containing one root and one side direction.  The
physical realization of this terminal two-dimensional piece is the hereditary
planar C\'ordoba estimate \cite{Cordoba77}
\[
 \int_{\R^2}m_Z^2\lesssim N\delta\bigl(1+\log(1/\delta)\bigr),
 \qquad
 m_Z\lesssim_\zeta
 \delta^{-\zeta}\frac{1+\log(1/\delta)}{\lambda}
\]
after deletion of an \(O(\delta^\zeta)\) fraction of shaded incidence mass.
With the slowly varying choice of \(\zeta\) made after
\eqref{eq:root-fibre-hard-cutoff}, this is the base case for
\eqref{eq:closed-support-recurrence}.  In particular, when
\(\lambda_2\ge\delta^{o(1)}\), it gives
\(B_{K_2}\lesssim\delta^{-o(1)}\).
The proper-predecessor recombination above reduces R6 to
\eqref{eq:root-cross-energy-gate}.  Conditional on that estimate, ordinary
finite induction on \(|I|\), using the latching, confluence/load and fresh
exact-support terms at each stage, propagates the estimate through the support
filtration.  The forgotten root coordinate is restored by
\eqref{eq:root-fibre-hard-cutoff} before the next extension step.  This is
induction in the physical realization of the already constructed typed
coefficient filtration, not an intrinsic construction of motives.

\section{Finite descent}

Fix the finite support, decoration, history, and Cartan ceilings from the
standing scope.  If
\(N_s=\prod_{p\in\Pi}p^{a_p}\), set
\begin{equation}\label{eq:descent-complexity}
 \mathfrak c(\eta;I,\chi,\nu,n;N_s)
 :=
 |\eta|_{\rm val}+|I|+|\nu|_1+\ell(\chi)+n+\sum_{p\in\Pi}a_p.
\end{equation}
A proper support face lowers \(|I|\), a strict predecessor lowers
\(\ell(\chi)\), a decoration reduction lowers its own depth, and a Cartan
predecessor lowers \(n\).  Exact-order loss replaces \(N_s\) by a proper
divisor.  The structural interchange identities ensure that lowering one
coordinate does not increase another, so every branch is finite.

At each state the audit \eqref{eq:complete-obstruction-list} determines the
next proper piece.  On the physical leg, a zero determinant is reconstructed
from proper faces; a nonzero determinant passes to the root antecedent and
then to one lower support.  The identity
\eqref{eq:latching-root-comparison} records the parallel motivic
normalization but is not used to infer physical nonvanishing.  Passive
decoration fibres disappear under \eqref{eq:physical-vertical-augmentation};
history/load fibres remain in \eqref{eq:rn-corrected-evaluation}.

Consequently formal termination carries no analytic gain.  Conditional on
\eqref{eq:root-cross-energy-gate}, the closed support recurrence and the
two-dimensional base estimate control every support stage.  The root-cross
gate is indispensable: neither the complexity decrease nor motivic
reconstruction proves a global incidence estimate.

\section{Closure of the filtered route}

\subsection{Structural bookkeeping}

The fixed-root determinant comparison
\eqref{eq:normalized-controller-comparison}, literal side faces, root
projection, insertion kernels, and the symmetric packet law
\eqref{eq:symmetric-root-tuple-law} verify the finite physical filtration.
Global gluing and every hereditary restriction are controlled by
\eqref{eq:global-local-fsc-sum} and
\eqref{eq:global-restriction-rn}.  The operators
\eqref{eq:physical-koszul-operator} give the structural representation
\eqref{eq:global-fsc-physical-koszul-comparison}; the passive decoration base
change \eqref{eq:physical-fsc-full-base-change} does not alter physical
multiplicity.

After the augmentation \eqref{eq:physical-vertical-augmentation}, the only
analytic data are the exact-support determinant, the confluence/load fibre,
and the root antecedent.  The reduced-history tree and
\eqref{eq:lca-positive-charge} control every strict predecessor.  Thus the
global confluence estimate R6 is equivalent to the remaining physical
root-cross gate \eqref{eq:root-cross-energy-gate}; no formal motivic
reconstruction is used to prove it.

\subsection{Scalar analytic closure}

In dimension two the hereditary planar C\'ordoba estimate \cite{Cordoba77} supplies the base
case.  In dimension three,
\eqref{eq:weighted-planar-cordoba}--%
\eqref{eq:affine-fibre-three-dimensional-refinement}, together with the sticky
inputs \eqref{eq:terminal-sticky-carrier-gate}, close the root gate.

At the three-to-four-dimensional handoff, the factorized and fibre-disjoint
blocks close by \eqref{eq:factorized-fibre-lift} and
\eqref{eq:fibre-disjoint-lift}.  Directionwise sampling and the unbiased lift
close every shell covered by
\eqref{eq:convexified-high-occupancy-energy}.  Proper-history fibre collisions
are absorbed by \eqref{eq:marked-fibre-proper-lca-energy}; the residual
low-occupancy source is exactly
\eqref{eq:low-occupancy-marked-source}.  Its broad part is multilinear, its
narrow part retains the affine fibre mark, and the uniform marked theorem
\cite{StickyContact4D} supplies
\eqref{eq:terminal-marked-sticky-gate}.  Hence the scalar route closes through
dimension four.

For \(d\ge5\), the analogous source-hereditary marked root gate remains an
input.  Assuming \eqref{eq:root-cross-energy-gate} at every retained support
and load stage, the support induction
\eqref{eq:support-induction-closure} gives
\[
 B_{K_I}\lesssim_d\delta^{-o(1)}.
\]
The final weighted-to-unweighted cutoff
\eqref{eq:terminal-unweighting} then retains subpower incidence mass.  For
sparse initial shading the same argument keeps the explicit density factors in
\eqref{eq:closed-support-recurrence}; the filtration creates no density.

\subsection{The separate typed requirement}

The hereditary typed statement requires a positive common-point routing on the
correlated pair source.  Its exact forms are the Hall/Radon--Nikodym condition
\eqref{eq:pair-routing-hall-cut} and the equivalent capacity system
\eqref{eq:coherent-capacity-dual-cut}; the fractional-mode bound
\eqref{eq:fractional-mode-capacity} is sufficient.  The structural tower
supplies labelled destinations and boundaries, but neither this positivity nor
the weighted capacity.  Consequently the scalar theorem above does not claim
the typed comparison.

\subsection{Terminal incidence count}

Finally convert the incidence set \(E\subset\Omega_{\rm inc}\) into the
terminal ordinary descendant
\[
 Y'(T)
 :=\{x\in Y(T):(Q(x),T,x)\in E\},
 \qquad
 m'(x):=\sum_T{\bf1}_{Y'(T)}(x)=m_E(x).
\]
The incidence identity gives
\begin{equation}\label{eq:terminal-incidence-volume}
 \sum_T|Y'(T)|
 =\int m'(x)\,dx
 \le
 \|m'\|_\infty
 \left|\bigcup_TY'(T)\right|.
\end{equation}
If \(\mathcal I_{\rm in}:=\sum_T|Y(T)|\), then
\(\mathcal I_{\rm in}=W\) and the cell partition gives
\[
 \sum_T|Y'(T)|
 =W\mu_{\rm inc}(E).
\]
Thus \eqref{eq:terminal-unweighting} gives
\[
 \sum_T|Y'(T)|
 \ge\frac\alpha2\,\mathcal I_{\rm in}.
\]
Under the same root-cross hypothesis, the mass ledger and terminal cutoff
therefore imply
\[
 \left|\bigcup_TY'(T)\right|
 \gtrsim_d
 \delta^{o(1)}\mathcal I_{\rm in}.
\]
In particular, for \(N_T\) direction-separated tubes with
\(|Y(T)|\ge\lambda_{\rm in}\delta^{d-1}\),
\[
 \left|\bigcup_TY'(T)\right|
 \gtrsim_d
 \delta^{o(1)}N_T\lambda_{\rm in}\delta^{d-1}.
\]
This last passage uses only incidence counting.  After taking the scalar root
gate as a physical hypothesis, the unproved analytic content is confined to
\eqref{eq:root-cross-energy-gate}, rather than hidden in a local-to-global root
assembly.  The already realized \(q^2\) confluence provenance is structural
data and is not used as an analytic estimate in this implication.

\begin{thebibliography}{99}
\bibitem{UnifiedTower}
\emph{Selmer--Cartan and Motivic Moore--Reedy Towers: Source Geometry,
Full-Motive Realization, and a (31)-adic Witness}, unified core draft v0.1,
2026.
\bibitem{Cordoba77}
A. C\'ordoba,
\emph{The Kakeya maximal function and the spherical summation multipliers},
Amer. J. Math. 99 (1977), no. 1, 1--22.
\bibitem{BCT06}
J. Bennett, A. Carbery and T. Tao,
\emph{On the multilinear restriction and Kakeya conjectures},
Acta Math. 196 (2006), no. 2, 261--302; arXiv:math/0509262.
\bibitem{Wolff95}
T. Wolff,
\emph{An improved bound for Kakeya type maximal functions},
Rev. Mat. Iberoam. 11 (1995), no. 3, 651--674.
\bibitem{KatzTao02}
N. H. Katz and T. Tao,
\emph{New bounds for Kakeya problems},
J. Anal. Math. 87 (2002), 231--263; arXiv:math/0102135.
\bibitem{WangZahl3D}
H. Wang and J. Zahl,
\emph{Volume estimates for unions of convex sets, and the Kakeya set
conjecture in three dimensions}, arXiv:2502.17655, 2025.
\bibitem{StickyKakeya}
H. Wang and J. Zahl,
\emph{Sticky Kakeya sets and the sticky Kakeya conjecture},
arXiv:2210.09581, 2022.
\bibitem{ChoudhuriSticky4D}
M. Rai Choudhuri,
\emph{An improved bound on the Hausdorff dimension of sticky Kakeya sets
in \(\mathbb R^4\)}, arXiv:2410.23579, 2024.
\bibitem{StickyContact4D}
\emph{Sticky Kakeya in \(\mathbb R^4\) as a Contact--Symplectic Incidence
Problem: Lossless Edge-Flow Closure and a Marked Finite-Scale Estimate},
companion manuscript v0.82, 2026.
\bibitem{GuthWangZahl3D}
L. Guth, H. Wang and J. Zahl,
\emph{A streamlined proof of the Kakeya set conjecture in
\(\mathbb R^3\)}, arXiv:2601.14411, 2026.
\bibitem{ZahlKakeyaSurvey}
J. Zahl,
\emph{A survey of the Kakeya conjecture, 2000--2025}, arXiv:2512.09397,
2025.
\bibitem{Guth10}
L. Guth,
\emph{The endpoint case of the Bennett--Carbery--Tao multilinear Kakeya
conjecture}, Acta Math. 205 (2010), 263--286; arXiv:0811.2251.
\end{thebibliography}

\end{document}
